\chapter{The Six Operations in Brave New Motivic Homotopy Theory}
\label{chap:six-operations}

\section{Purpose, scope, and standing conventions}
\label{sec:six-purpose}

For every quasi-compact quasi-separated spectral scheme \(S\),
Chapter~\ref{chap:stable-brave-new-motivic} constructed a compactly generated
stable presentably symmetric monoidal category
\[
    \SH(S),
\]
together with an adjunction
\[
    f^*:\SH(Y)
    \rightleftarrows
    \SH(X):f_*
\]
for every morphism
\[
    f:X\longrightarrow Y
\]
between quasi-compact quasi-separated spectral schemes.  For every smooth
morphism \(p:X\to Y\) of finite presentation, it also constructed an
adjunction
\[
    p_\sharp:\SH(X)
    \rightleftarrows
    \SH(Y):p^*,
\]
stable smooth base change, the smooth projection formula, stable
localization and recollement, closed exchange, stable homotopy purity, and
virtual Thom twists
\[
    \Sigma^\xi,
    \qquad
    \xi\in K(X).
\]

Chapter~\ref{chap:proper-ambidexterity-compactification} supplied the proper
geometric input required for exceptional functoriality.  In particular, it
proved that arbitrary direct image preserves small colimits, that proper
direct image satisfies arbitrary base change and the projection formula,
that open--proper exchange and proper support hold, and that every separated
morphism of finite presentation admits a compactification.  It further
constructed nonempty cofiltered \(\infty\)-categories of compactifications
and of finite compactification grids, together with all refinement and
composition coherences.

The purpose of the present chapter is to assemble those results into the
exceptional adjunction
\[
    f_!:\SH(X)
    \rightleftarrows
    \SH(Y):f^!
\]
for every separated morphism of finite presentation.  We prove coherent
composition, exceptional base change, the exceptional projection formula,
internal-Hom identities, smooth purity, smooth base change for arbitrary
direct image, localization in six-functor
notation, and the associated Borel--Moore and duality formalisms.  We then
package the construction as a lax symmetric monoidal functor on a category of
correspondences.

\begin{convention}[Exceptional morphisms]
\label{conv:six-exceptional-morphisms}
All spectral schemes occurring in a motivic coefficient category are
quasi-compact and quasi-separated.  We retain the class
\[
    \mathcal E
\]
of separated morphisms of finite presentation from
Convention~\ref{conv:proper-geometric-classes}.

The symbol \(f_!\) is used only when
\[
    f\in\mathcal E.
\]
For such an \(f\), the notation \(f^!\) denotes the right adjoint of the
exceptional direct image \(f_!\) constructed in this chapter.

Independently, if
\[
    i:Z\hookrightarrow X
\]
is a closed immersion with quasi-compact open complement, we retain the
notation
\[
    i^!:\SH(X)\longrightarrow\SH(Z)
\]
for the closed exceptional inverse image constructed in
Chapter~\ref{chap:stable-brave-new-motivic} as the right adjoint of
\(i_*\).  When \(i\) is of finite presentation and hence belongs to
\(\mathcal E\), Proposition~\ref{prop:exceptional-inverse-normalization}
canonically identifies this pre-existing closed exceptional inverse image
with the right adjoint of the exceptional direct image
\[
    i_!\simeq i_*.
\]
Thus the two uses of the notation \(i^!\) agree whenever both are defined.

For an arbitrary morphism \(f\), the adjunction
\[
    f^*:\SH(Y)
    \rightleftarrows
    \SH(X):f_*
\]
remains available.  Since
Chapter~\ref{chap:proper-ambidexterity-compactification}
proved that \(f_*\) preserves all small colimits, it has a further right
adjoint by the adjoint functor theorem.  We do not introduce notation for
that auxiliary right adjoint.
\end{convention}

\begin{convention}[Compactifications]
\label{conv:six-compactifications}
For \(f:X\to Y\) in \(\mathcal E\), a compactification means a factorization
\[
    X\xrightarrow{j}\overline X\xrightarrow{p}Y
\]
with \(j\) a quasi-compact open immersion and \(p\) a morphism whose
classical truncation is proper.  We write
\[
    \operatorname{Comp}(f)
\]
for the \(\infty\)-category of compactifications and proper strict
refinements constructed in
Definition~\ref{def:category-compactifications}.
\end{convention}

\begin{convention}[Universes and compactification colimits]
\label{conv:six-compactification-universes}
Fix Grothendieck universes
\[
    \mathcal U\in\mathcal V.
\]
All spectral schemes, morphisms, and geometric diagrams are
\(\mathcal U\)-small.  The coefficient categories
\(\SH(S)\) and their functor categories are regarded as
\(\mathcal V\)-categories.  The categories
\[
    \operatorname{Comp}(f),
    \qquad
    \operatorname{Comp}(f_1,\ldots,f_n),
\]
and the finite compactification-grid categories are consequently regarded
as \(\mathcal V\)-small indexing categories.

All compactification colimits are formed in the corresponding
\(\mathcal V\)-presentable functor category.  We suppress the associated
change-of-universe functors from the notation.
\end{convention}

\begin{remark}[Absolute purity]
\label{rem:six-absolute-purity-warning}
The six operations and smooth purity constructed in this chapter are
unconditional under the standing geometric scope.

For every quasi-smooth closed immersion with finite locally free conormal
module in the scope treated below---including regular closed immersions in
the discrete case when they satisfy the standing finiteness
hypotheses---there is a canonical coefficientwise purity transformation
\[
    \Sigma^{-\nu_i}i^*
    \longrightarrow
    i^!.
\]
Its invertibility on every object is the additional property usually called
\emph{absolute purity}.

Integral motivic stable homotopy theory does not supply such an equivalence
for every regular closed immersion over arbitrary bases merely as a formal
consequence of the six operations.  Absolute purity is a property of the
chosen coefficient object or coefficient theory.  Accordingly, the chapter
constructs the transformation and proves the formal properties stated
below, but asserts invertibility only for objects that are pure for the
immersion.

This distinction is already present in ordinary motivic homotopy theory;
see, for example, \cite{HoyoisSixOperations} and
\cite{DegliseJinKhanFundamentalClasses}.
\end{remark}

\subsection{The chapter-level output}

The constructions below culminate in the following data.
\begin{enumerate}
    \item For every \(f\in\mathcal E\), a colimit-preserving exceptional
    direct image
    \[
        f_!:\SH(X)\longrightarrow\SH(Y)
    \]
    independent of compactification through a contractible space of
    coherent equivalences.

    \item Coherent identity and composition equivalences
    \[
        (\id_X)_!\simeq\id_{\SH(X)},
        \qquad
        g_!f_!\simeq(gf)_!.
    \]

    \item Exceptional base change and the exceptional projection formula.

    \item A right adjoint
    \[
        f_!\dashv f^!
    \]
    with coherent contravariant composition and the corresponding
    right-adjoint exchange identities.

    \item Normalizations
    \[
        j_!\simeq j_\sharp
        \quad\text{for quasi-compact open }j,
        \qquad
        p_!\simeq p_*
        \quad\text{for proper }p.
    \]

    \item Smooth purity
    \[
        p^!\simeq\Sigma^{T_p}p^*,
        \qquad
        p_!\simeq p_\sharp\Sigma^{-T_p},
    \]
    and smooth base change for arbitrary direct image along every smooth,
    separated morphism \(g\) of finite presentation,
    \[
        g^*f_*
        \simeq
        f'_*(g')^*.
    \]

    \item A lax symmetric monoidal correspondence action
    \[
        \bigl(X\xleftarrow{u}K\xrightarrow{v}Y\bigr)
        \longmapsto
        v_!u^*.
    \]

    \item Six-functor localization, Borel--Moore objects, relative
    dualizing objects, and Grothendieck--Verdier exchange.
\end{enumerate}

\section{Exceptional direct image by compactification}
\label{sec:exceptional-direct-image}

Let
\[
    f:X\longrightarrow Y
\]
belong to \(\mathcal E\).  Chapter~\ref{chap:proper-ambidexterity-compactification}
constructed a functor
\[
    \Phi_f:
    \operatorname{Comp}(f)
    \longrightarrow
    \operatorname{Fun}^L
    \bigl(
      \SH(X),
      \SH(Y)
    \bigr)
\]
by the formula
\[
    \Phi_f
    \bigl(
      X\xrightarrow{j}\overline X\xrightarrow{p}Y
    \bigr)
    :=
    p_*j_\sharp.
\]
Every refinement is carried to an equivalence by proper support, and all
comparison maps satisfy coherent composition.

\begin{definition}[Exceptional direct image]
\label{def:exceptional-direct-image}
Define
\[
    f_!
    :=
    \operatorname*{colim}_{c\in\operatorname{Comp}(f)}
    \Phi_f(c)
\]
in
\[
    \operatorname{Fun}^L
    \bigl(
      \SH(X),
      \SH(Y)
    \bigr).
\]
\end{definition}

\begin{theorem}[Evaluation on a compactification]
\label{thm:exceptional-evaluation-compactification}
Let
\[
    c=
    \bigl(
      X\xrightarrow{j}\overline X\xrightarrow{p}Y
    \bigr)
\]
be any compactification of \(f\).  The canonical cocone morphism is an
equivalence
\[
    \epsilon_c:
    p_*j_\sharp
    \xrightarrow{\simeq}
    f_!.
\]
The space of compatible inverses to the family of equivalences
\(\epsilon_c\) is contractible.  In particular, every compactification
computes \(f_!\) canonically up to a contractible space of coherent choices.
\end{theorem}

\begin{proof}
By Proposition~\ref{prop:compactification-direct-image-diagram}, the functor
\(\Phi_f\) carries every morphism of \(\operatorname{Comp}(f)\) to an
equivalence and therefore factors through the groupoid completion
\[
    \bigl|\operatorname{Comp}(f)\bigr|.
\]
Theorem~\ref{thm:compactification-category-contractible} gives
\[
    \bigl|\operatorname{Comp}(f)\bigr|
    \simeq
    *.
\]
A diagram of equivalences indexed by a contractible groupoid has colimit
canonically equivalent to each of its values.  Hence the cocone morphism
from \(\Phi_f(c)=p_*j_\sharp\) to the colimit is an equivalence.  The space
of inverses to an equivalence is contractible, and the same remains true
after imposing compatibility with the contractible indexing groupoid.
\end{proof}

\begin{corollary}[Cocontinuity]
\label{cor:exceptional-direct-image-cocontinuous}
For every \(f\in\mathcal E\), the functor
\[
    f_!:\SH(X)\longrightarrow\SH(Y)
\]
preserves all small colimits.
\end{corollary}

\begin{proof}
The colimit in Definition~\ref{def:exceptional-direct-image} is formed in
\(\operatorname{Fun}^L(\SH(X),\SH(Y))\).  Equivalently, choose a
compactification \(f=pj\).  The functor \(j_\sharp\) preserves colimits as a
left adjoint, and \(p_*\) preserves colimits by
Theorem~\ref{thm:arbitrary-direct-image-cocontinuous}.  Thus
\(p_*j_\sharp\), and hence \(f_!\), is colimit-preserving.
\end{proof}

\begin{proposition}[Functoriality under refinement]
\label{prop:exceptional-refinement-functoriality}
Let
\[
    r:
    \bigl(
      X\xrightarrow{j_2}\overline X_2\xrightarrow{p_2}Y
    \bigr)
    \longrightarrow
    \bigl(
      X\xrightarrow{j_1}\overline X_1\xrightarrow{p_1}Y
    \bigr)
\]
be a proper refinement.  Under the evaluation equivalences of
Theorem~\ref{thm:exceptional-evaluation-compactification}, the comparison
\[
    p_{2*}j_{2\sharp}
    \simeq
    p_{1*}r_*j_{2\sharp}
    \xrightarrow{\simeq}
    p_{1*}j_{1\sharp}
\]
is identified with the identity transformation of \(f_!\).
\end{proposition}

\begin{proof}
The displayed map is precisely \(\Phi_f(r)\).  The cocone defining the
colimit is natural in the indexing category, so
\[
    \epsilon_{c_1}\Phi_f(r)
    =
    \epsilon_{c_2}.
\]
Conjugating by the evaluation equivalences identifies \(\Phi_f(r)\) with
\(\id_{f_!}\).
\end{proof}

\section{Identity and normalization}
\label{sec:exceptional-normalization}

\begin{proposition}[Identity normalization]
\label{prop:exceptional-identity}
For every quasi-compact quasi-separated spectral scheme \(X\), there is a
canonical equivalence
\[
    (\id_X)_!
    \xrightarrow{\simeq}
    \id_{\SH(X)}.
\]
These equivalences are compatible with the composition equivalences
constructed below.
\end{proposition}

\begin{proof}
The factorization
\[
    X\xrightarrow{\id_X}X\xrightarrow{\id_X}X
\]
is a compactification of \(\id_X\).  Theorem~\ref{thm:exceptional-evaluation-compactification}
therefore gives
\[
    (\id_X)_!
    \simeq
    (\id_X)_*(\id_X)_\sharp
    \simeq
    \id_{\SH(X)}.
\]
Compatibility with composition follows from the identity restrictions of
finite compactification grids.
\end{proof}

\begin{theorem}[Open normalization]
\label{thm:exceptional-open-normalization}
Let
\[
    j:U\hookrightarrow X
\]
be a quasi-compact open immersion.  There is a canonical equivalence
\[
    j_!
    \xrightarrow{\simeq}
    j_\sharp.
\]
It is compatible with composition and arbitrary base change of
quasi-compact open immersions.
\end{theorem}

\begin{proof}
The factorization
\[
    U\xrightarrow{j}X\xrightarrow{\id_X}X
\]
is a compactification of \(j\).  Evaluation gives
\[
    j_!
    \simeq
    (\id_X)_*j_\sharp
    \simeq
    j_\sharp.
\]
Composition and base-change compatibility follow from the corresponding
properties of the compactification evaluation maps and stable smooth base
change, since open immersions are smooth.
\end{proof}

\begin{theorem}[Proper normalization]
\label{thm:exceptional-proper-normalization}
Let
\[
    p:X\longrightarrow Y
\]
be proper and belong to \(\mathcal E\).  There is a canonical equivalence
\[
    p_!
    \xrightarrow{\simeq}
    p_*.
\]
It is compatible with composition of proper morphisms and with arbitrary
base change.
\end{theorem}

\begin{proof}
The factorization
\[
    X\xrightarrow{\id_X}X\xrightarrow{p}Y
\]
is a compactification.  Hence
\[
    p_!
    \simeq
    p_*(\id_X)_\sharp
    \simeq
    p_*.
\]
For a composite of proper morphisms, the associated compactification grid
has only identity open immersions, so the composition equivalence below
reduces to the canonical direct-image composition equivalence.  Base-change
compatibility reduces to proper base change.
\end{proof}

\begin{corollary}[Closed normalization]
\label{cor:exceptional-closed-normalization}
For every closed immersion \(i:Z\hookrightarrow X\) of finite presentation,
there is a canonical equivalence
\[
    i_!\simeq i_*.
\]
\end{corollary}

\begin{proof}
A closed immersion belongs to \(\mathcal P\), so this is
Theorem~\ref{thm:exceptional-proper-normalization}.
\end{proof}

\section{Composition of exceptional direct image}
\label{sec:exceptional-composition}

Let
\[
    X\xrightarrow{f}Y\xrightarrow{g}Z
\]
belong to \(\mathcal E\).  Chapter~\ref{chap:proper-ambidexterity-compactification}
constructed the nonempty cofiltered \(\infty\)-category
\[
    \operatorname{Comp}(f,g)
\]
of compatible composition data.  Every object determines
compactifications of \(f\), \(g\), and \(gf\), together with an
open--proper Cartesian square.

\begin{construction}[Binary composition]
\label{constr:exceptional-binary-composition}
Let \(d\in\operatorname{Comp}(f,g)\).  In the notation of
Definition~\ref{def:compactification-composition-data}, open--proper
exchange gives
\[
    r_*k_\sharp
    \xrightarrow{\simeq}
    j_{g\sharp}p_{f*}.
\]
Taking its inverse and composing with the proper direct-image composition
constraint gives an equivalence
\[
\begin{aligned}
    p_{g*}j_{g\sharp}p_{f*}j_{f\sharp}
    &\xrightarrow{\simeq}
    p_{g*}r_*k_\sharp j_{f\sharp} \\
    &\xrightarrow{\simeq}
    (p_gr)_*(kj_f)_\sharp.
\end{aligned}
\]
Using Theorem~\ref{thm:exceptional-evaluation-compactification} on the three
compactifications defines
\[
    \mu_{g,f}(d):
    g_!f_!
    \xrightarrow{\simeq}
    (gf)_!.
\]
\end{construction}

\begin{theorem}[Coherent composition]
\label{thm:exceptional-composition}
For every composable pair
\[
    X\xrightarrow{f}Y\xrightarrow{g}Z
\]
in \(\mathcal E\), Construction~\ref{constr:exceptional-binary-composition}
induces a canonical equivalence
\[
    \mu_{g,f}:
    g_!f_!
    \xrightarrow{\simeq}
    (gf)_!.
\]
It is independent of the chosen composition datum and is compatible with
proper refinements.

For every finite composable string
\[
    X_n\xrightarrow{f_n}X_{n-1}
    \longrightarrow\cdots\longrightarrow
    X_1\xrightarrow{f_1}X_0,
\]
all parenthesized composites
\[
    f_{1!}f_{2!}\cdots f_{n!}
\]
are canonically and coherently equivalent to
\[
    (f_1f_2\cdots f_n)_!.
\]
The equivalences satisfy unitality, associativity, and all higher
composition coherences.
\end{theorem}

\begin{proof}
Naturality of open--proper exchange and its compatibility with proper
support show that a refinement of composition data carries the equivalence
\(\mu_{g,f}(d)\) to the corresponding equivalence for the refined datum.
Proposition~\ref{prop:compactification-composition-data} gives
\[
    \bigl|\operatorname{Comp}(f,g)\bigr|
    \simeq
    *.
\]
Hence the compatible family \(\mu_{g,f}(d)\) descends to a canonical
comparison independent of \(d\).

For a finite string, choose a compactification grid.  Proposition~\ref{prop:finite-grid-composition-formula}
identifies every parenthesized composite of the associated compactification
functors with the functor associated to the full interval.  Theorem~\ref{thm:finite-compactification-grids}
shows that the groupoid completion of the grid category is contractible.
Restriction of grids along monotone maps of finite ordinals identifies
faces, degeneracies, and pasted subdivisions.  The compatibility in
Proposition~\ref{prop:finite-grid-composition-formula} therefore gives all
identity, associativity, and higher simplicial coherences after passage to
the canonical exceptional direct images.
\end{proof}

\begin{corollary}[Unital exceptional functoriality]
\label{cor:exceptional-unital-functoriality}
The assignments
\[
    X\longmapsto\SH(X),
    \qquad
    f\longmapsto f_!
\]
form a homotopy-coherent covariant functor on the subcategory of
quasi-compact quasi-separated spectral schemes and morphisms in
\(\mathcal E\).
\end{corollary}

\begin{proof}
Combine Proposition~\ref{prop:exceptional-identity} with
Theorem~\ref{thm:exceptional-composition}.
\end{proof}

\section{Exceptional base change}
\label{sec:exceptional-base-change}

Consider a Cartesian square
\[
\begin{tikzcd}[column sep=large,row sep=large]
    X' \arrow[r,"g'"] \arrow[d,"f'"']
        & X \arrow[d,"f"] \\
    Y' \arrow[r,"g"']
        & Y,
\end{tikzcd}
\]
where \(f\in\mathcal E\).  Then \(f'\in\mathcal E\).

\begin{construction}[Exceptional Beck--Chevalley transformation]
\label{constr:exceptional-base-change}
Choose a compactification
\[
    X\xrightarrow{j}\overline X\xrightarrow{p}Y
\]
of \(f\).  Its pullback is a compactification
\[
    X'\xrightarrow{j'}\overline X'\xrightarrow{p'}Y'
\]
of \(f'\).  Write
\[
    \overline g:\overline X'\longrightarrow\overline X
\]
for the base-change morphism.  Proper base change and stable smooth base
change for the open immersion \(j\) give
\[
\begin{aligned}
    g^*p_*j_\sharp
    &\xrightarrow{\simeq}
    p'_*\overline g^*j_\sharp \\
    &\xrightarrow{\simeq}
    p'_*j'_\sharp(g')^*.
\end{aligned}
\]
Conjugating by the evaluation equivalences defines
\[
    \operatorname{Ex}^{*!}_{f,g}:
    g^*f_!
    \longrightarrow
    f'_!(g')^*.
\]
\end{construction}

\begin{theorem}[Exceptional base change]
\label{thm:exceptional-base-change}
In every Cartesian square as above, the canonical transformation
\[
    \operatorname{Ex}^{*!}_{f,g}:
    g^*f_!
    \xrightarrow{\simeq}
    f'_!(g')^*
\]
is an equivalence.  It is independent of the compactification used in its
construction and satisfies the horizontal and vertical Beck--Chevalley
pasting laws.  It is compatible with composition of exceptional morphisms
and with arbitrary iterated base change.
\end{theorem}

\begin{proof}
The transformation of Construction~\ref{constr:exceptional-base-change} is
a composite of equivalences, so it is an equivalence for every chosen
compactification.  If
\[
    r:\overline X_2\longrightarrow\overline X_1
\]
is a proper refinement, its pullback is a proper refinement of the
base-changed compactifications.  Naturality of proper base change, smooth
base change, and proper support identifies the two resulting composites.
The contractibility of \(\operatorname{Comp}(f)\) therefore makes the
transformation independent of the compactification.

The Beck--Chevalley pasting laws for the two elementary ingredients imply
the horizontal and vertical pasting laws on every compactification.  Since
the evaluation equivalences and compactification base-change functors are
coherent, these pasting identities descend to \(f_!\).  Compatibility with
composition of exceptional morphisms follows by applying the same argument
to finite compactification grids and using
Proposition~\ref{prop:finite-grid-base-change}.
\end{proof}

\begin{corollary}[Normalization of exceptional base change]
\label{cor:exceptional-base-change-normalization}
Under the normalization equivalences of
Section~\ref{sec:exceptional-normalization}, exceptional base change agrees
with:
\begin{enumerate}
    \item proper base change when \(f\) is proper;
    \item stable closed base change when \(f\) is a closed immersion;
    \item stable smooth base change when \(f\) is a quasi-compact open
    immersion.
\end{enumerate}
The identifications are equalities of canonical adjoint mates.
\end{corollary}

\begin{proof}
Use respectively the compactifications
\[
    X\xrightarrow{\id}X\xrightarrow{f}Y,
\]
for proper and closed \(f\), and
\[
    X\xrightarrow{f}Y\xrightarrow{\id}Y
\]
for open \(f\).  Construction~\ref{constr:exceptional-base-change} then
reduces to the stated canonical Beck--Chevalley transformation.
\end{proof}

\section{Exceptional projection and external products}
\label{sec:exceptional-projection}

\begin{lemma}[Module-linearity of open--proper exchange]
\label{lem:open-proper-exchange-module-linearity}
Consider a Cartesian square
\[
\begin{tikzcd}[column sep=large,row sep=large]
    U' \arrow[r,"j'"] \arrow[d,"q"']
        & X' \arrow[d,"p"] \\
    U \arrow[r,"j"']
        & X,
\end{tikzcd}
\]
where \(j\) is a quasi-compact open immersion and \(p\) has proper
classical truncation.  The open--proper exchange equivalence
\[
    p_*j'_\sharp
    \xrightarrow{\simeq}
    j_\sharp q_*
\]
is linear over \(\SH(X)\).

More explicitly, for
\[
    A\in\SH(U'),
    \qquad
    B\in\SH(X),
\]
the diagram
\[
\begin{tikzcd}[column sep=huge,row sep=large]
    p_*j'_\sharp
    \bigl(
      A\otimes_{U'}(pj')^*B
    \bigr)
        \arrow[r,"\simeq"]
        \arrow[d,"\simeq"']
        &
    j_\sharp q_*
    \bigl(
      A\otimes_{U'}q^*j^*B
    \bigr)
        \arrow[d,"\simeq"]
    \\
    p_*j'_\sharp A\otimes_XB
        \arrow[r,"\simeq"']
        &
    j_\sharp q_*A\otimes_XB
\end{tikzcd}
\]
commutes.  The vertical maps are the composites of the open and proper
projection-formula equivalences, and the horizontal maps are open--proper
exchange.

These module-linearity identifications are compatible with proper
refinement, Cartesian base change, and horizontal and vertical pasting of
open--proper squares.
\end{lemma}

\begin{proof}
Since
\[
    pj'=jq,
\]
there is a canonical identification
\[
    (pj')^*B
    \simeq
    q^*j^*B.
\]

Proper support identifies both
\[
    p_*j'_\sharp
    \bigl(
      A\otimes(pj')^*B
    \bigr)
\]
and
\[
    p_*j'_\sharp A\otimes B
\]
with objects in the essential image of the fully faithful functor
\[
    j_\sharp:\SH(U)\longrightarrow\SH(X).
\]
It is therefore enough to verify commutativity after applying \(j^*\).

Proper base change and full faithfulness of \(j'_\sharp\) give
\[
\begin{aligned}
    j^*p_*j'_\sharp
    \bigl(
      A\otimes(pj')^*B
    \bigr)
    &\simeq
    q_*j'^*j'_\sharp
    \bigl(
      A\otimes q^*j^*B
    \bigr) \\
    &\simeq
    q_*
    \bigl(
      A\otimes q^*j^*B
    \bigr).
\end{aligned}
\]
Likewise,
\[
\begin{aligned}
    j^*
    \bigl(
      p_*j'_\sharp A\otimes_XB
    \bigr)
    &\simeq
    j^*p_*j'_\sharp A\otimes_Uj^*B \\
    &\simeq
    q_*A\otimes_Uj^*B \\
    &\simeq
    q_*
    \bigl(
      A\otimes_{U'}q^*j^*B
    \bigr),
\end{aligned}
\]
where the final equivalence is the proper projection formula for \(q\).

Under these identifications, both paths around the displayed square become
the same proper projection-formula transformation.  Hence the square
commutes after applying \(j^*\), and therefore commutes in \(\SH(X)\).

Compatibility with refinement, base change, and pasting follows from the
corresponding compatibilities of proper base change, proper support, the
open and proper projection formulas, and full faithfulness of open direct
image.
\end{proof}

\begin{construction}[Exceptional projection transformation]
\label{constr:exceptional-projection}
Let
\[
    f:X\longrightarrow Y
\]
belong to \(\mathcal E\), and let
\[
    A\in\SH(X),
    \qquad
    B\in\SH(Y).
\]
Choose a compactification
\[
    X\xrightarrow{j}\overline X\xrightarrow{p}Y.
\]
Since
\[
    f^*B\simeq j^*p^*B,
\]
the open projection formula and the proper projection formula give
\[
\begin{aligned}
    p_*j_\sharp
    \bigl(
      A\otimes_Xf^*B
    \bigr)
    &\simeq
    p_*j_\sharp
    \bigl(
      A\otimes_Xj^*p^*B
    \bigr) \\
    &\xrightarrow{\simeq}
    p_*
    \bigl(
      j_\sharp A\otimes_{\overline X}p^*B
    \bigr) \\
    &\xrightarrow{\simeq}
    p_*j_\sharp A\otimes_YB.
\end{aligned}
\]
Conjugating by compactification evaluation defines the exceptional
projection transformation
\[
    \operatorname{PF}^!_f(A,B):
    f_!
    \bigl(
      A\otimes_Xf^*B
    \bigr)
    \longrightarrow
    f_!A\otimes_YB.
\]
\end{construction}

\begin{theorem}[Exceptional projection formula]
\label{thm:exceptional-projection-formula}
For every \(f\in\mathcal E\), the canonical exceptional projection
transformation is an equivalence:
\[
    f_!
    \bigl(
      A\otimes_Xf^*B
    \bigr)
    \xrightarrow{\simeq}
    f_!A\otimes_YB.
\]
It is natural in \(A\) and \(B\), compatible with composition of
exceptional morphisms, and compatible with exceptional base change.
Thus \(f_!\) is a strong \(\SH(Y)\)-module functor.
\end{theorem}

\begin{proof}
For every compactification
\[
    X\xrightarrow{j}\overline X\xrightarrow{p}Y
\]
of \(f\), the transformation of
Construction~\ref{constr:exceptional-projection}
is the composite
\[
\begin{aligned}
    p_*j_\sharp
    \bigl(
      A\otimes_Xj^*p^*B
    \bigr)
    &\xrightarrow{\simeq}
    p_*
    \bigl(
      j_\sharp A\otimes_{\overline X}p^*B
    \bigr) \\
    &\xrightarrow{\simeq}
    p_*j_\sharp A\otimes_YB.
\end{aligned}
\]
The first arrow is the open projection formula and the second is the proper
projection formula.  It is therefore an equivalence.

Let
\[
    r:
    \bigl(
      X\xrightarrow{j_2}\overline X_2\xrightarrow{p_2}Y
    \bigr)
    \longrightarrow
    \bigl(
      X\xrightarrow{j_1}\overline X_1\xrightarrow{p_1}Y
    \bigr)
\]
be a proper refinement.  Naturality of the open and proper projection
formulas, compatibility of proper support with tensor products, and
Lemma~\ref{lem:open-proper-exchange-module-linearity}
identify the projection transformation associated with the first
compactification with the transformation associated with the second.
Thus the compactificationwise transformations form a compatible family over
\(\operatorname{Comp}(f)\).

Since the groupoid completion of \(\operatorname{Comp}(f)\) is
contractible, this compatible family descends to a canonical equivalence
\[
    f_!
    \bigl(
      A\otimes_Xf^*B
    \bigr)
    \xrightarrow{\simeq}
    f_!A\otimes_YB.
\]
Naturality in \(A\) and \(B\) is inherited from the open and proper
projection formulas.

Now let
\[
    X\xrightarrow{f}Y\xrightarrow{g}Z
\]
be composable exceptional morphisms.  Choose a finite compactification grid
for the pair \((f,g)\).  On every cell of the grid, the exceptional
composition comparison is assembled from proper direct-image composition
and open--proper exchange.

The projection formula for the full interval is the pasting of the
projection formulas for the two subintervals.  This follows from:

\begin{enumerate}
    \item composition compatibility of the open projection formula;

    \item composition compatibility of the proper projection formula;

    \item module-linearity of open--proper exchange from
    Lemma~\ref{lem:open-proper-exchange-module-linearity}.
\end{enumerate}

Consequently, on every compactification grid, the diagram comparing
\[
    g_!f_!
    \bigl(
      A\otimes_X(gf)^*B
    \bigr)
\]
with
\[
    (gf)_!A\otimes_ZB
\]
commutes.  The finite-grid coherence theorem and contractibility of the
grid category descend this equality to the canonical exceptional direct
images.  This proves compatibility with composition of exceptional
morphisms.

For a Cartesian base-change square, choose a compactification of \(f\) and
its pulled-back compactification.  Compatibility of the open and proper
projection formulas with their respective Beck--Chevalley transformations,
together with the base-change compatibility of
Lemma~\ref{lem:open-proper-exchange-module-linearity},
shows that the exceptional projection formula commutes with exceptional
base change on each compactification.  Contractibility again descends the
comparison to \(f_!\).

Thus \(f_!\) is a strong \(\SH(Y)\)-module functor.
\end{proof}

\begin{corollary}[Projection with target Thom twists]
\label{cor:exceptional-target-thom-projection}
Let \(f:X\to Y\) belong to \(\mathcal E\), let \(\xi\in K(Y)\), and let
\(A\in\SH(X)\).  There is a canonical equivalence
\[
    f_!\Sigma^{f^*\xi}A
    \xrightarrow{\simeq}
    \Sigma^\xi f_!A.
\]
It is compatible with addition in \(K(Y)\), composition in \(f\), and
base change.
\end{corollary}

\begin{proof}
Apply Theorem~\ref{thm:exceptional-projection-formula} to the invertible
Thom object \(\operatorname{th}_Y(\xi)\) and use pullback naturality of the
motivic \(J\)-homomorphism.
\end{proof}

\subsection{External products}

Let
\[
    f:X\longrightarrow Y,
    \qquad
    g:X'\longrightarrow Y'
\]
belong to \(\mathcal E\).  For
\[
    A\in\SH(X),
    \qquad
    B\in\SH(X'),
\]
write
\[
    A\boxtimes B
    :=
    \operatorname{pr}_X^*A
    \otimes
    \operatorname{pr}_{X'}^*B
    \in
    \SH(X\times X').
\]

\begin{proposition}[Exceptional external products]
\label{prop:exceptional-external-products}
There is a canonical equivalence
\[
    f_!A\boxtimes g_!B
    \xrightarrow{\simeq}
    (f\times g)_!
    \bigl(
      A\boxtimes B
    \bigr).
\]
It is symmetric, unital, associative, compatible with base change, and
compatible with composition of exceptional morphisms.
\end{proposition}

\begin{proof}
Factor
\[
    f\times g
    =
    (f\times\id_{Y'})
    (\id_X\times g).
\]
Exceptional base change for the Cartesian square obtained from \(f\) by
pullback along \(Y\times Y'\to Y\) gives
\[
    \operatorname{pr}_Y^*f_!A
    \simeq
    (f\times\id_{Y'})_!
    \operatorname{pr}_X^*A
\]
on \(Y\times Y'\).  Hence the exceptional projection formula gives
\[
\begin{aligned}
    f_!A\boxtimes g_!B
    &\simeq
    (f\times\id_{Y'})_!
    \Bigl(
      \operatorname{pr}_X^*A
      \otimes
      \operatorname{pr}_{Y'}^*g_!B
    \Bigr).
\end{aligned}
\]
On \(X\times Y'\), exceptional base change for \(g\) gives
\[
    \operatorname{pr}_{Y'}^*g_!B
    \simeq
    (\id_X\times g)_!
    \operatorname{pr}_{X'}^*B.
\]
A second application of the exceptional projection formula therefore
identifies the preceding object with
\[
\begin{aligned}
    (f\times\id_{Y'})_!
    (\id_X\times g)_!
    \Bigl(
      \operatorname{pr}_X^*A
      \otimes
      \operatorname{pr}_{X'}^*B
    \Bigr).
\end{aligned}
\]
Exceptional composition gives the asserted equivalence with
\[
    (f\times g)_!(A\boxtimes B).
\]
All coherences follow from the Beck--Chevalley and projection-formula
pasting laws.
\end{proof}

\section{Exceptional inverse image}
\label{sec:exceptional-inverse-image}

\begin{definition}[Exceptional inverse image]
\label{def:exceptional-inverse-image}
Let \(f:X\to Y\) belong to \(\mathcal E\).  Since \(f_!\) is
colimit-preserving between presentable \(\infty\)-categories, it admits a
right adjoint.  Define
\[
    f^!:\SH(Y)\longrightarrow\SH(X)
\]
by the adjunction
\[
    f_!:\SH(X)
    \rightleftarrows
    \SH(Y):f^!.
\]
\end{definition}

\begin{proposition}[Composition of exceptional inverse image]
\label{prop:exceptional-inverse-composition}
For composable morphisms
\[
    X\xrightarrow{f}Y\xrightarrow{g}Z
\]
in \(\mathcal E\), taking right adjoints of
Theorem~\ref{thm:exceptional-composition} gives a canonical equivalence
\[
    (gf)^!
    \xrightarrow{\simeq}
    f^!g^!.
\]
These equivalences are coherently unital and associative.
\end{proposition}

\begin{proof}
The right adjoint of \(g_!f_!\) is \(f^!g^!\), while the right adjoint of
\((gf)_!\) is \((gf)^!\).  Take the right adjoint mate of
\(\mu_{g,f}\).  Coherent unitality and associativity follow from the
corresponding coherences for exceptional direct image and the higher
functoriality of adjoint mates.
\end{proof}

\begin{proposition}[Normalization of exceptional inverse image]
\label{prop:exceptional-inverse-normalization}
The following canonical equivalences hold.

\begin{enumerate}
    \item If \(j:U\hookrightarrow X\) is a quasi-compact open immersion,
    then
    \[
        j^!\simeq j^*.
    \]

    \item If \(p:X\to Y\) is proper and belongs to \(\mathcal E\), then
    \(p^!\) is canonically the right adjoint of \(p_*\).

    \item If \(i:Z\hookrightarrow X\) is a closed immersion of finite
    presentation, then the exceptional inverse image of
    Definition~\ref{def:exceptional-inverse-image} is canonically
    equivalent to the closed exceptional inverse image constructed in
    Chapter~\ref{chap:stable-brave-new-motivic}.
\end{enumerate}

The three identifications are compatible with composition.  They identify
the canonical right-adjoint base-change transformations for the normalized
functors.  They also identify the pullback--exceptional exchange
transformations whenever those transformations are invertible; in
particular, this holds for smooth base change.
\end{proposition}

\begin{proof}
For a quasi-compact open immersion, open normalization gives
\[
    j_!\simeq j_\sharp.
\]
The right adjoint of \(j_\sharp\) is \(j^*\), so uniqueness of right
adjoints gives
\[
    j^!\simeq j^*.
\]

For a proper morphism,
\[
    p_!\simeq p_*
\]
by proper normalization.  Their right adjoints are therefore canonically
equivalent.

A closed immersion of finite presentation is proper and belongs to
\(\mathcal E\).  Its exceptional direct image is consequently
\[
    i_!\simeq i_*.
\]
The earlier closed exceptional inverse image was defined as the right
adjoint of \(i_*\).  Uniqueness of right adjoints gives the canonical
identification with the present \(i^!\).

Composition compatibility follows by taking right adjoints of the
composition-compatible normalization equivalences for \(j_!\), \(p_!\),
and \(i_!\).

Right-adjoint base-change compatibility follows by taking mates of the
corresponding normalized base-change equivalences for the left adjoints.
For pullback--exceptional exchange, normalization identifies the canonical
transformation but does not make it invertible for an arbitrary base
change.  Its invertibility under smooth base change follows from
Proposition~\ref{prop:smooth-pullback-exceptional-exchange}.
\end{proof}

\begin{theorem}[Right-adjoint exceptional base change]
\label{thm:right-adjoint-exceptional-base-change}
For every Cartesian square
\[
\begin{tikzcd}[column sep=large,row sep=large]
    X' \arrow[r,"g'"] \arrow[d,"f'"']
        & X \arrow[d,"f"] \\
    Y' \arrow[r,"g"']
        & Y
\end{tikzcd}
\]
with \(f\in\mathcal E\), there is a canonical equivalence
\[
    (g')_*f'^!
    \xrightarrow{\simeq}
    f^!g_*.
\]
It is the right adjoint mate of exceptional base change and satisfies the
horizontal and vertical pasting laws.
\end{theorem}

\begin{proof}
The left adjoint of \((g')_*f'^!\) is \(f'_!(g')^*\), while the left
adjoint of \(f^!g_*\) is \(g^*f_!\).  Take the right adjoint mate of the
equivalence
\[
    g^*f_!
    \xrightarrow{\simeq}
    f'_!(g')^*.
\]
The pasting laws are preserved by the calculus of mates.
\end{proof}


\section{Internal-Hom identities}
\label{sec:exceptional-internal-hom}

\begin{theorem}[Exceptional Hom--projection identity]
\label{thm:exceptional-hom-projection}
Let \(f:X\to Y\) belong to \(\mathcal E\).  For
\[
    A\in\SH(X),
    \qquad
    B\in\SH(Y),
\]
there is a canonical equivalence
\[
    \underline{\operatorname{Hom}}_Y
    \bigl(
      f_!A,
      B
    \bigr)
    \xrightarrow{\simeq}
    f_*
    \underline{\operatorname{Hom}}_X
    \bigl(
      A,
      f^!B
    \bigr).
\]
It is natural in \(A\), \(B\), and \(f\), and compatible with composition.
For every base-change square, it commutes with the canonical comparison
transformations for direct image, exceptional inverse image, and internal
Hom.  Whenever those comparison transformations are equivalences, the
resulting base-change comparison is an equivalence.
\end{theorem}

\begin{proof}
Let \(C\in\SH(Y)\).  Using the adjunctions, closedness of the tensor
products, and the exceptional projection formula, there are natural
equivalences
\[
\begin{aligned}
    \Map_Y
    \Bigl(
      C,
      f_*\underline{\operatorname{Hom}}_X(A,f^!B)
    \Bigr)
    &\simeq
    \Map_X
    \Bigl(
      f^*C,
      \underline{\operatorname{Hom}}_X(A,f^!B)
    \Bigr) \\
    &\simeq
    \Map_X
    \bigl(
      f^*C\otimes_XA,
      f^!B
    \bigr) \\
    &\simeq
    \Map_Y
    \bigl(
      f_!(f^*C\otimes_XA),
      B
    \bigr) \\
    &\simeq
    \Map_Y
    \bigl(
      C\otimes_Yf_!A,
      B
    \bigr) \\
    &\simeq
    \Map_Y
    \Bigl(
      C,
      \underline{\operatorname{Hom}}_Y(f_!A,B)
    \Bigr).
\end{aligned}
\]
Yoneda gives the equivalence.  Naturality and composition compatibility
follow from the coherent adjunction and projection-formula data used in the
calculation.  Applying the canonical comparison transformations to the same
mapping-space calculation gives the asserted base-change compatibility; no
invertibility is asserted unless the individual comparison transformations
are equivalences.
\end{proof}

\begin{theorem}[Exceptional pullback--Hom identity]
\label{thm:exceptional-pullback-hom}
Let \(f:X\to Y\) belong to \(\mathcal E\).  For
\[
    B,C\in\SH(Y),
\]
there is a canonical equivalence
\[
    f^!
    \underline{\operatorname{Hom}}_Y(B,C)
    \xrightarrow{\simeq}
    \underline{\operatorname{Hom}}_X
    \bigl(
      f^*B,
      f^!C
    \bigr).
\]
\end{theorem}

\begin{proof}
Let \(A\in\SH(X)\).  Then
\[
\begin{aligned}
    \Map_X
    \Bigl(
      A,
      f^!\underline{\operatorname{Hom}}_Y(B,C)
    \Bigr)
    &\simeq
    \Map_Y
    \Bigl(
      f_!A,
      \underline{\operatorname{Hom}}_Y(B,C)
    \Bigr) \\
    &\simeq
    \Map_Y
    \bigl(
      f_!A\otimes_YB,
      C
    \bigr) \\
    &\simeq
    \Map_Y
    \bigl(
      f_!(A\otimes_Xf^*B),
      C
    \bigr) \\
    &\simeq
    \Map_X
    \bigl(
      A\otimes_Xf^*B,
      f^!C
    \bigr) \\
    &\simeq
    \Map_X
    \Bigl(
      A,
      \underline{\operatorname{Hom}}_X(f^*B,f^!C)
    \Bigr).
\end{aligned}
\]
Yoneda gives the assertion.
\end{proof}

\begin{corollary}[Dualizable coefficient formula]
\label{cor:exceptional-dualizable-coefficient}
Let \(B\in\SH(Y)\) be strongly dualizable.  For every \(C\in\SH(Y)\),
there is a canonical equivalence
\[
    f^!C\otimes_Xf^*B
    \xrightarrow{\simeq}
    f^!(C\otimes_YB).
\]
\end{corollary}

\begin{proof}
For dualizable \(B\), tensoring by \(B\) is canonically identified with
internal Hom from \(B^\vee\).  Apply
Theorem~\ref{thm:exceptional-pullback-hom} to \(B^\vee\).
\end{proof}

\section{Classical comparison of exceptional operations}
\label{sec:exceptional-classical-comparison}

\begin{convention}[Ordinary exceptional operations]
\label{conv:ordinary-exceptional-operations}
For every ordinary quasi-compact quasi-separated scheme \(S\), we fix the
ordinary stable motivic category
\[
    \SH_{\mathrm{cl}}(S)
\]
with its standard exceptional formalism for separated morphisms of finite
presentation.  Thus, for every such morphism
\[
    h:T\longrightarrow S,
\]
the symbols
\[
    h_!:
    \SH_{\mathrm{cl}}(T)
    \rightleftarrows
    \SH_{\mathrm{cl}}(S):h^!
\]
denote an adjunction fixed independently of the spectral construction in
this chapter.

The ordinary exceptional direct image is normalized by compactification:
for every factorization
\[
    T\xrightarrow{j}\overline T\xrightarrow{p}S
\]
with \(j\) a quasi-compact open immersion and \(p\) proper, there is a
canonical equivalence
\[
    h_!
    \simeq
    p_*j_\sharp.
\]
These equivalences are compatible with proper refinements, composition,
base change, projection formulas, and open and proper normalization.  The
ordinary smooth-purity normalization is
\[
    h^!
    \simeq
    \Sigma^{T_h}h^*,
    \qquad
    h_!
    \simeq
    h_\sharp\Sigma^{-T_h}
\]
for smooth separated \(h\) of finite presentation.  We use this fixed
ordinary formalism, including its coherent adjunction and exchange data,
throughout the comparison and specialization arguments below; see
\cite{HoyoisSixOperations} and
\cite{DegliseJinKhanFundamentalClasses}.
\end{convention}

Let
\[
    \Theta_S^{\mathrm{st}}:
    \SH(S)
    \xrightarrow{\simeq}
    \SH_{\mathrm{cl}}(S_{\mathrm{cl}})
\]
be the stable comparison equivalence of
Theorem~\ref{thm:stable-classical-comparison}.

\begin{theorem}[Comparison for exceptional direct image]
\label{thm:comparison-exceptional-direct-image}
Let \(f:X\to Y\) belong to \(\mathcal E\).  There is a canonical
equivalence
\[
    \Theta_Y^{\mathrm{st}}f_!
    \xrightarrow{\simeq}
    (f_{\mathrm{cl}})_!
    \Theta_X^{\mathrm{st}}.
\]
It is compatible with identities, composition, exceptional base change,
the exceptional projection formula, open and proper normalization, and
external products.
\end{theorem}

\begin{proof}
Choose a compactification
\[
    X\xrightarrow{j}\overline X\xrightarrow{p}Y.
\]
Compatibility of stable comparison with open direct image and arbitrary
direct image gives
\[
\begin{aligned}
    \Theta_Y^{\mathrm{st}}p_*j_\sharp
    &\simeq
    (p_{\mathrm{cl}})_*
    \Theta_{\overline X}^{\mathrm{st}}j_\sharp \\
    &\simeq
    (p_{\mathrm{cl}})_*
    (j_{\mathrm{cl}})_\sharp
    \Theta_X^{\mathrm{st}}.
\end{aligned}
\]
Since \(p_{\mathrm{cl}}\) is proper and \(j_{\mathrm{cl}}\) is open, the
right-hand side is the ordinary compactification formula for
\((f_{\mathrm{cl}})_!\).  The comparison is natural under proper
refinements because stable comparison preserves proper support and
open--proper exchange.  It therefore descends through the contractible
compactification category to the asserted equivalence.

Each compatibility follows by applying stable comparison to the
compactification or grid construction of the corresponding spectral
transformation.  On every compactification the result is the canonical
ordinary transformation, and contractibility descends the identification.
\end{proof}

\begin{corollary}[Comparison for exceptional inverse image]
\label{cor:comparison-exceptional-inverse-image}
For every \(f\in\mathcal E\), there is a canonical equivalence
\[
    \Theta_X^{\mathrm{st}}f^!
    \xrightarrow{\simeq}
    (f_{\mathrm{cl}})^!
    \Theta_Y^{\mathrm{st}}.
\]
It is compatible with composition, right-adjoint base change, and the
internal-Hom identities.
\end{corollary}

\begin{proof}
Take right adjoints of the equivalence in
Theorem~\ref{thm:comparison-exceptional-direct-image}.  The compatibility
statements follow from the calculus of adjoint mates.
\end{proof}

\begin{corollary}[Derived nilpotent invariance of the exceptional package]
\label{cor:exceptional-derived-invariance}
The exceptional operations and every transformation constructed in
Sections~\ref{sec:exceptional-direct-image}--\ref{sec:exceptional-internal-hom}
are invariant under passage to classical truncation.  In particular, they
are determined by the corresponding ordinary six-operation data under the
stable comparison equivalences.
\end{corollary}

\begin{proof}
Combine Theorem~\ref{thm:comparison-exceptional-direct-image} and
Corollary~\ref{cor:comparison-exceptional-inverse-image} with stable derived
nilpotent invariance.
\end{proof}


\section{Smooth purity}
\label{sec:six-smooth-purity}

Let
\[
    p:X\longrightarrow Y
\]
be smooth, separated, and of finite presentation.  Recall the tangent Thom
class
\[
    T_p=[\Omega_{X/Y}]\in K(X)
\]
from Convention~\ref{conv:tangent-normal-classes}.

\begin{construction}[Smooth-purity transformation]
\label{constr:smooth-purity-transformation}
Ordinary stable motivic homotopy theory has a canonical smooth-purity
equivalence
\[
    \mathfrak p^{\mathrm{cl}}_p:
    \Sigma^{T_{p_{\mathrm{cl}}}}
    p_{\mathrm{cl}}^*
    \xrightarrow{\simeq}
    p_{\mathrm{cl}}^!.
\]

There is a canonical cotangent-comparison morphism
\[
    \iota_X^*\Omega_{X/Y}
    \longrightarrow
    \Omega_{X_{\mathrm{cl}}/Y_{\mathrm{cl}}}.
\]
Since \(p\) is smooth, the relative cotangent module
\[
    \Omega_{X/Y}
\]
is finite locally free.  Its pullback
\[
    \iota_X^*\Omega_{X/Y}
\]
is therefore a discrete finite locally free
\(\mathcal O_{X_{\mathrm{cl}}}\)-module.  The classical truncation
\(p_{\mathrm{cl}}\) is smooth, so
\[
    \Omega_{X_{\mathrm{cl}}/Y_{\mathrm{cl}}}
\]
is also a discrete finite locally free
\(\mathcal O_{X_{\mathrm{cl}}}\)-module.

The degree-zero cotangent comparison identifies their underlying classical
modules.  Hence the canonical comparison morphism is an equivalence:
\[
    \iota_X^*\Omega_{X/Y}
    \xrightarrow{\simeq}
    \Omega_{X_{\mathrm{cl}}/Y_{\mathrm{cl}}}.
\]
It determines a canonical path in \(K(X_{\mathrm{cl}})\)
\[
    \iota_X^*T_p
    \simeq
    T_{p_{\mathrm{cl}}}.
\]

Compatibility of stable comparison with the motivic
\(J\)-homomorphism therefore gives a canonical equivalence
\[
    \Theta_X^{\mathrm{st}}
    \Sigma^{T_p}p^*
    \simeq
    \Sigma^{T_{p_{\mathrm{cl}}}}
    p_{\mathrm{cl}}^*
    \Theta_Y^{\mathrm{st}}.
\]

Theorem~\ref{thm:comparison-exceptional-direct-image} and
Corollary~\ref{cor:comparison-exceptional-inverse-image}
give a canonical equivalence
\[
    \Theta_X^{\mathrm{st}}p^!
    \simeq
    p_{\mathrm{cl}}^!\Theta_Y^{\mathrm{st}}.
\]
Consequently, conjugation by
\[
    \Theta_X^{\mathrm{st}}
    \quad\text{and}\quad
    \Theta_Y^{\mathrm{st}}
\]
induces an equivalence between the space of natural transformations
\[
    \operatorname{Nat}
    \bigl(
      \Sigma^{T_p}p^*,
      p^!
    \bigr)
\]
and the space of natural transformations
\[
    \operatorname{Nat}
    \bigl(
      \Sigma^{T_{p_{\mathrm{cl}}}}p_{\mathrm{cl}}^*,
      p_{\mathrm{cl}}^!
    \bigr).
\]

Define
\[
    \mathfrak p_p:
    \Sigma^{T_p}p^*
    \longrightarrow
    p^!
\]
to be the transformation corresponding to
\[
    \mathfrak p^{\mathrm{cl}}_p
\]
under this equivalence of transformation spaces.
\end{construction}

\begin{theorem}[Smooth purity]
\label{thm:six-smooth-purity}
Let \(p:X\to Y\) be smooth, separated, and of finite presentation.  The
smooth-purity transformation is an equivalence:
\[
    \mathfrak p_p:
    \Sigma^{T_p}p^*
    \xrightarrow{\simeq}
    p^!.
\]
It is compatible with arbitrary base change, composition of smooth
separated morphisms of finite presentation, and virtual Thom twists pulled
back from the target.
\end{theorem}

\begin{proof}
Stable comparison carries \(\mathfrak p_p\) to the ordinary smooth-purity
equivalence by construction.  Since \(\Theta_X^{\mathrm{st}}\) is an
equivalence, it is conservative, and therefore \(\mathfrak p_p\) is an
equivalence.

The ordinary smooth-purity equivalences satisfy base-change and composition
compatibility.  Under the path
\[
    T_{qp}
    \simeq
    T_p+p^*T_q
\]
induced by cotangent transitivity and the motivic \(J\)-homomorphism, the
ordinary purity map for \(qp\) is the pasted composite of the purity maps
for \(p\) and \(q\).  Stable comparison preserves each operation and all
adjunction data in these diagrams.  Transport therefore gives the spectral
compatibilities, including the higher associativity coherences.  Target
Thom-twist compatibility follows from the exceptional projection formula
and pullback naturality of Thom twists.
\end{proof}

\begin{corollary}[Smooth exceptional direct image]
\label{cor:smooth-exceptional-direct-image}
Let \(p:X\to Y\) be smooth, separated, and of finite presentation.  There
is a canonical equivalence
\[
    p_!
    \xrightarrow{\simeq}
    p_\sharp\Sigma^{-T_p}.
\]
It is compatible with base change and composition.
\end{corollary}

\begin{proof}
The functor \(p_\sharp\Sigma^{-T_p}\) is left adjoint to
\[
    \Sigma^{T_p}p^*.
\]
By Theorem~\ref{thm:six-smooth-purity}, the latter is equivalent to \(p^!\),
whose left adjoint is \(p_!\).  Uniqueness of left adjoints gives the
canonical equivalence.  Compatibility is inherited from smooth purity.
\end{proof}

\begin{corollary}[\'{E}tale purity]
\label{cor:six-etale-purity}
Let \(p:X\to Y\) be \'{e}tale, separated, and of finite presentation.  Then
\[
    p^!\simeq p^*,
    \qquad
    p_!\simeq p_\sharp.
\]
\end{corollary}

\begin{proof}
For an \'{e}tale morphism,
\[
    T_p=[\Omega_{X/Y}]=0.
\]
Apply Theorem~\ref{thm:six-smooth-purity} and
Corollary~\ref{cor:smooth-exceptional-direct-image}.
\end{proof}

\begin{theorem}[Compatibility with smooth proper ambidexterity]
\label{thm:smooth-purity-ambidexterity-compatibility}
Let \(p:X\to Y\) be smooth, proper, and of finite presentation.  Under the
proper normalization
\[
    p_!\simeq p_*,
\]
the smooth exceptional equivalence
\[
    p_\sharp\Sigma^{-T_p}
    \xrightarrow{\simeq}
    p_!
\]
agrees with the ambidexterity equivalence
\[
    p_\sharp\Sigma^{-T_p}
    \xrightarrow{\simeq}
    p_*
\]
constructed in
Theorem~\ref{thm:smooth-proper-ambidexterity}.
\end{theorem}

\begin{proof}
Stable comparison carries both transformations to the canonical ordinary
smooth-proper ambidexterity transformation.  The comparison functor is an
equivalence and therefore induces an equivalence on transformation spaces.
The two spectral transformations consequently agree.
\end{proof}

\begin{lemma}[Purity compatibility with Beck--Chevalley mates]
\label{lem:purity-beck-chevalley-mates}
Consider a Cartesian square
\[
\begin{tikzcd}[column sep=large,row sep=large]
    X' \arrow[r,"g'"] \arrow[d,"f'"']
        & X \arrow[d,"f"] \\
    Y' \arrow[r,"g"']
        & Y
\end{tikzcd}
\]
with \(f\in\mathcal E\), and suppose that \(g\) is smooth, separated, and
of finite presentation.

\begin{enumerate}
    \item Take the mate of exceptional Beck--Chevalley
    \[
        g^*f_!
        \xrightarrow{\simeq}
        f'_!(g')^*
    \]
    with respect to the adjunctions
    \[
        f_!\dashv f^!,
        \qquad
        f'_!\dashv f'^!,
    \]
    leaving \(g^*\) and \((g')^*\) unchanged.  The resulting transformation
    is
    \[
        (g')^*f^!
        \longrightarrow
        f'^!g^*.
    \]
    Under the smooth-purity equivalences for \(g\) and \(g'\), it is the
    Thom-cancellation equivalence
    \[
        (g')^*f^!
        \xrightarrow{\simeq}
        f'^!g^*.
    \]
    It is therefore equal to
    \(\operatorname{Ex}^{!*}_{f,g}\).

    \item Regard the square in transposed form:
    \[
    \begin{tikzcd}[column sep=large,row sep=large]
        X' \arrow[r,"f'"] \arrow[d,"g'"']
            & Y' \arrow[d,"g"] \\
        X \arrow[r,"f"']
            & Y.
    \end{tikzcd}
    \]
    Take the right-adjoint mate of
    \[
        f^*g_!
        \xrightarrow{\simeq}
        g'_!f'^*
    \]
    with respect to
    \[
        f^*\dashv f_*,
        \qquad
        f'^*\dashv f'_*,
        \qquad
        g_!\dashv g^!,
        \qquad
        g'_!\dashv g'^!.
    \]
    The resulting equivalence is
    \[
        f'_*g'^!
        \xrightarrow{\simeq}
        g^!f_*.
    \]
    After conjugation by smooth purity, cotangent base change, the
    dualizable projection formula, and cancellation of
    \(\Sigma^{T_g}\), its inverse is the canonical ordinary
    Beck--Chevalley transformation
    \[
        g^*f_*
        \longrightarrow
        f'_*(g')^*.
    \]
\end{enumerate}

Both identifications are compatible with horizontal and vertical pasting.
\end{lemma}

\begin{proof}
For part~(1), let
\[
    \beta_{f,g}:
    (g')^*f^!
    \longrightarrow
    f'^!g^*
\]
be the mate of exceptional Beck--Chevalley with respect to
\[
    f_!\dashv f^!,
    \qquad
    f'_!\dashv f'^!.
\]
Explicitly, it is the composite
\[
\begin{aligned}
    (g')^*f^!
    &\longrightarrow
    f'^!f'_!(g')^*f^! \\
    &\xrightarrow{\simeq}
    f'^!g^*f_!f^! \\
    &\longrightarrow
    f'^!g^*.
\end{aligned}
\]
The first arrow is the unit of \(f'_!\dashv f'^!\), the middle arrow is
induced by exceptional Beck--Chevalley, and the final arrow is induced by
the counit of \(f_!\dashv f^!\).  Thus
\[
    \beta_{f,g}
    =
    \operatorname{Ex}^{!*}_{f,g}.
\]

Smooth purity, exceptional inverse-image composition, and cotangent base
change also define a Thom-cancellation equivalence
\[
    \tau_{f,g}:
    (g')^*f^!
    \xrightarrow{\simeq}
    f'^!g^*
\]
by the composite
\[
\begin{aligned}
    (g')^*f^!
    &\xrightarrow{\simeq}
    \Sigma^{-T_{g'}}g'^!f^! \\
    &\xrightarrow{\simeq}
    \Sigma^{-T_{g'}}(fg')^! \\
    &\xrightarrow{\simeq}
    \Sigma^{-T_{g'}}(gf')^! \\
    &\xrightarrow{\simeq}
    \Sigma^{-T_{g'}}f'^!g^! \\
    &\xrightarrow{\simeq}
    \Sigma^{-T_{g'}}f'^!\Sigma^{T_g}g^* \\
    &\xrightarrow{\simeq}
    \Sigma^{-T_{g'}}\Sigma^{f'^*T_g}f'^!g^* \\
    &\xrightarrow{\simeq}
    f'^!g^*.
\end{aligned}
\]
The penultimate equivalence is the dualizable coefficient formula, and the
last equivalence uses
\[
    T_{g'}
    \simeq
    f'^*T_g.
\]

Apply stable classical comparison to both transformations.  Compatibility
of comparison with exceptional direct and inverse image, the units and
counits of their adjunction, exceptional Beck--Chevalley, exceptional
inverse-image composition, smooth purity, and Thom twists carries
\(\beta_{f,g}\) to the ordinary mate of exceptional Beck--Chevalley and
carries \(\tau_{f,g}\) to the ordinary Thom-cancellation transformation.
In the ordinary six-operation formalism, smooth purity is compatible with
exceptional Beck--Chevalley under the calculus of mates; hence those two
ordinary transformations agree.  See
\cite{HoyoisSixOperations}.

Stable comparison is an equivalence and therefore induces an equivalence on
spaces of natural transformations.  It follows that
\[
    \beta_{f,g}
    =
    \tau_{f,g}.
\]
This proves part~(1).

For part~(2), let
\[
    \delta_{f,g}:
    f'_*g'^!
    \xrightarrow{\simeq}
    g^!f_*
\]
be the right-adjoint mate of
\[
    f^*g_!
    \xrightarrow{\simeq}
    g'_!f'^*.
\]
Conjugating \(\delta_{f,g}\) by smooth purity gives
\[
    f'_*\Sigma^{T_{g'}}(g')^*
    \xrightarrow{\simeq}
    \Sigma^{T_g}g^*f_*.
\]
Cotangent base change and the dualizable projection formula identify its
source by
\[
\begin{aligned}
    f'_*\Sigma^{T_{g'}}(g')^*
    &\simeq
    f'_*\Sigma^{f'^*T_g}(g')^* \\
    &\simeq
    \Sigma^{T_g}f'_*(g')^*.
\end{aligned}
\]
Cancelling \(\Sigma^{T_g}\) therefore gives an equivalence
\[
    f'_*(g')^*
    \xrightarrow{\simeq}
    g^*f_*.
\]
Let
\[
    \kappa_{f,g}:
    g^*f_*
    \xrightarrow{\simeq}
    f'_*(g')^*
\]
be its inverse.

Stable comparison carries \(\delta_{f,g}\), all of the displayed
conjugating equivalences, and \(\kappa_{f,g}\) to their ordinary
counterparts.  In the ordinary six-operation formalism, the same
purity--mate compatibility identifies the resulting inverse with the
canonical ordinary Beck--Chevalley transformation.  Since stable comparison
is fully faithful on transformation spaces,
\(\kappa_{f,g}\) is the canonical Beck--Chevalley transformation in the
spectral theory.

The ordinary identifications are compatible with horizontal and vertical
pasting.  Stable comparison preserves the pasting diagrams, cotangent
base-change and transitivity paths, the dualizable projection formula, and
Thom cancellation.  Since it reflects equality of transformations, both
identifications satisfy the asserted horizontal and vertical pasting laws.
\end{proof}


\section{Smooth exchange consequences}
\label{sec:smooth-exchange-consequences}

\begin{construction}[Pullback--exceptional exchange]
\label{constr:pullback-exceptional-exchange}
For a Cartesian square
\[
\begin{tikzcd}[column sep=large,row sep=large]
    X' \arrow[r,"g'"] \arrow[d,"f'"']
        & X \arrow[d,"f"] \\
    Y' \arrow[r,"g"']
        & Y
\end{tikzcd}
\]
with \(f\in\mathcal E\), define a canonical transformation
\[
    \operatorname{Ex}^{!*}_{f,g}:
    (g')^*f^!
    \longrightarrow
    f'^!g^*
\]
by the composite
\[
\begin{aligned}
    (g')^*f^!
    &\longrightarrow
    f'^!f'_!(g')^*f^! \\
    &\xrightarrow{\simeq}
    f'^!g^*f_!f^! \\
    &\longrightarrow
    f'^!g^*.
\end{aligned}
\]
The first arrow is the unit of \(f'_!\dashv f'^!\), the middle equivalence
is exceptional base change, and the last arrow is induced by the counit of
\(f_!\dashv f^!\).
\end{construction}

\begin{proposition}[Smooth pullback--exceptional exchange]
\label{prop:smooth-pullback-exceptional-exchange}
In Construction~\ref{constr:pullback-exceptional-exchange}, suppose that
\(g\) is smooth, separated, and of finite presentation.  Then
\[
    \operatorname{Ex}^{!*}_{f,g}:
    (g')^*f^!
    \xrightarrow{\simeq}
    f'^!g^*
\]
is an equivalence.

These equivalences satisfy horizontal and vertical pasting and are
compatible with composition in \(f\).
\end{proposition}

\begin{proof}
The base change \(g'\) is smooth, separated, and of finite presentation.
Smooth purity and exceptional inverse-image composition give
\[
\begin{aligned}
    (g')^*f^!
    &\simeq
    \Sigma^{-T_{g'}}(g')^!f^! \\
    &\simeq
    \Sigma^{-T_{g'}}(fg')^! \\
    &\simeq
    \Sigma^{-T_{g'}}(gf')^! \\
    &\simeq
    \Sigma^{-T_{g'}}f'^!g^! \\
    &\simeq
    \Sigma^{-T_{g'}}f'^!\Sigma^{T_g}g^*.
\end{aligned}
\]
Cotangent base change gives
\[
    T_{g'}
    \simeq
    f'^*T_g.
\]
Since the Thom object \(\operatorname{th}_{Y'}(T_g)\) is invertible, the
dualizable coefficient formula gives
\[
    f'^!\Sigma^{T_g}g^*
    \simeq
    \Sigma^{f'^*T_g}f'^!g^*.
\]
Cancelling the two inverse Thom twists produces an equivalence
\[
    (g')^*f^!
    \xrightarrow{\simeq}
    f'^!g^*.
\]

By
Lemma~\ref{lem:purity-beck-chevalley-mates}(1),
this equivalence is precisely the canonical transformation
\(\operatorname{Ex}^{!*}_{f,g}\) of
Construction~\ref{constr:pullback-exceptional-exchange}.

Horizontal and vertical pasting, as well as compatibility with composition
in \(f\), now follow from the corresponding pasting laws in
Lemma~\ref{lem:purity-beck-chevalley-mates}.
\end{proof}

\begin{lemma}[Dualizable projection formula for arbitrary direct image]
\label{lem:arbitrary-direct-image-dualizable-projection}
Let
\[
    f:X\longrightarrow Y
\]
be any morphism of quasi-compact quasi-separated spectral schemes.  For
\(A\in\SH(X)\) and every strongly dualizable object \(L\in\SH(Y)\), the
canonical projection transformation is an equivalence
\[
    f_*A\otimes_YL
    \xrightarrow{\simeq}
    f_*
    \bigl(
      A\otimes_Xf^*L
    \bigr).
\]
It is coherent under composition in \(f\) and under pullback of the
dualizable coefficient.
\end{lemma}

\begin{proof}
Let \(C\in\SH(Y)\).  Since \(L\) is strongly dualizable and \(f^*\) is
strong symmetric monoidal, \(f^*L\) is strongly dualizable with dual
\(f^*L^\vee\).  There are natural equivalences
\[
\begin{aligned}
    \Map_Y
    \bigl(
      C,
      f_*A\otimes_YL
    \bigr)
    &\simeq
    \Map_Y
    \bigl(
      C\otimes_YL^\vee,
      f_*A
    \bigr) \\
    &\simeq
    \Map_X
    \bigl(
      f^*(C\otimes_YL^\vee),
      A
    \bigr) \\
    &\simeq
    \Map_X
    \bigl(
      f^*C\otimes_Xf^*L^\vee,
      A
    \bigr) \\
    &\simeq
    \Map_X
    \bigl(
      f^*C,
      A\otimes_Xf^*L
    \bigr) \\
    &\simeq
    \Map_Y
    \bigl(
      C,
      f_*(A\otimes_Xf^*L)
    \bigr).
\end{aligned}
\]
Yoneda gives the asserted equivalence.  The composite is built only from
the adjunction \(f^*\dashv f_*\), the strong monoidal constraint of
\(f^*\), and the evaluation and coevaluation maps of \(L\).  It is therefore
the canonical projection transformation and inherits the stated
coherences.
\end{proof}

\begin{theorem}[Smooth base change for arbitrary direct image]
\label{thm:smooth-base-change-arbitrary-direct-image}
Consider a Cartesian square
\[
\begin{tikzcd}[column sep=large,row sep=large]
    X' \arrow[r,"g'"] \arrow[d,"f'"']
        & X \arrow[d,"f"] \\
    Y' \arrow[r,"g"']
        & Y.
\end{tikzcd}
\]
If \(g\) is smooth, separated, and of finite presentation, then the
canonical ordinary Beck--Chevalley transformation is an equivalence
\[
    g^*f_*
    \xrightarrow{\simeq}
    f'_*(g')^*.
\]

These equivalences satisfy horizontal and vertical pasting, are compatible
with composition in \(f\), and are compatible with composition of smooth
separated morphisms of finite presentation in \(g\).
\end{theorem}

\begin{proof}
Regard the square in transposed form:
\[
\begin{tikzcd}[column sep=large,row sep=large]
    X' \arrow[r,"f'"] \arrow[d,"g'"']
        & Y' \arrow[d,"g"] \\
    X \arrow[r,"f"']
        & Y.
\end{tikzcd}
\]
Since \(g\) is smooth, separated, and of finite presentation, it belongs
to \(\mathcal E\).  Exceptional base change for \(g\) gives
\[
    f^*g_!
    \xrightarrow{\simeq}
    g'_!f'^*.
\]
Taking right adjoints gives
\[
    f'_*g'^!
    \xrightarrow{\simeq}
    g^!f_*.
\]

Smooth purity identifies this equivalence with
\[
    f'_*
    \Sigma^{T_{g'}}(g')^*
    \xrightarrow{\simeq}
    \Sigma^{T_g}g^*f_*.
\]
Cotangent base change gives
\[
    T_{g'}
    \simeq
    f'^*T_g.
\]
Applying
Lemma~\ref{lem:arbitrary-direct-image-dualizable-projection}
to \(f'\) and to the invertible object
\(\operatorname{th}_{Y'}(T_g)\) gives
\[
    f'_*
    \Sigma^{T_{g'}}(g')^*
    \simeq
    \Sigma^{T_g}f'_*(g')^*.
\]
Thus there is a canonical equivalence
\[
    \Sigma^{T_g}f'_*(g')^*
    \xrightarrow{\simeq}
    \Sigma^{T_g}g^*f_*.
\]
Cancelling \(\Sigma^{T_g}\) and inverting gives
\[
    g^*f_*
    \xrightarrow{\simeq}
    f'_*(g')^*.
\]

By
Lemma~\ref{lem:purity-beck-chevalley-mates}(2),
the resulting map is exactly the canonical ordinary Beck--Chevalley
transformation, rather than merely an abstract equivalence between its
source and target.

The horizontal and vertical pasting laws and the composition
compatibilities follow from
Lemma~\ref{lem:purity-beck-chevalley-mates},
exceptional Beck--Chevalley pasting, composition compatibility of smooth
purity, cotangent base change and transitivity, and coherence of the
dualizable projection formula.
\end{proof}


\section{Correspondences and coherent exceptional functoriality}
\label{sec:exceptional-correspondences}

\subsection{The correspondence category}

\begin{definition}[Exceptional correspondences]
\label{def:exceptional-correspondences}
Let
\[
    \operatorname{Corr}_{\mathcal E}
    \bigl(
      \operatorname{Sch}^{\mathrm{sp}}_{\mathrm{qcqs}}
    \bigr)
\]
be the category enriched in \(\operatorname{Cat}_\infty\) defined as
follows.

Its objects are quasi-compact quasi-separated spectral schemes.

For two objects \(X\) and \(Y\), let
\[
    \mathcal A_{\mathcal E,\mathcal P}(X,Y)
\]
be the \(\infty\)-category whose objects are morphisms
\[
    (u,v):
    K\longrightarrow X\times Y
\]
with \(K\) quasi-compact and quasi-separated and with
\[
    v:K\longrightarrow Y
\]
belonging to \(\mathcal E\).  A morphism
\[
    (K,u,v)
    \longrightarrow
    (K',u',v')
\]
in
\(\mathcal A_{\mathcal E,\mathcal P}(X,Y)\)
is a morphism
\[
    a:K\longrightarrow K'
\]
over \(X\times Y\) whose classical truncation is proper.

Define the hom \(\infty\)-category by
\[
    \operatorname{Corr}_{\mathcal E}(X,Y)
    :=
    \mathcal A_{\mathcal E,\mathcal P}(X,Y)^{\mathrm{op}}.
\]
Thus an object of
\(\operatorname{Corr}_{\mathcal E}(X,Y)\)
is a span
\[
    X\xleftarrow{u}K\xrightarrow{v}Y
\]
with \(v\in\mathcal E\), and a morphism
\[
    \bigl(
      X\xleftarrow{u'}K'\xrightarrow{v'}Y
    \bigr)
    \longrightarrow
    \bigl(
      X\xleftarrow{u}K\xrightarrow{v}Y
    \bigr)
\]
is represented by a proper morphism
\[
    a:K\longrightarrow K'
\]
over \(X\times Y\).

Composition
\[
    \operatorname{Corr}_{\mathcal E}(X,Y)
    \times
    \operatorname{Corr}_{\mathcal E}(Y,Z)
    \longrightarrow
    \operatorname{Corr}_{\mathcal E}(X,Z)
\]
is induced by derived fibre product of apices:
\[
    \bigl(
      X\xleftarrow{u}K\xrightarrow{v}Y
    \bigr),
    \quad
    \bigl(
      Y\xleftarrow{u'}L\xrightarrow{v'}Z
    \bigr)
\]
compose to
\[
    X
    \xleftarrow{u\operatorname{pr}_K}
    K\times_YL
    \xrightarrow{v'\operatorname{pr}_L}
    Z.
\]
The right leg belongs to \(\mathcal E\), since
\(\operatorname{pr}_L\) is a base change of \(v\) and
\(\mathcal E\) is stable under base change and composition.

Proper maps of apices are preserved by derived fibre product, so
composition is functorial on the hom \(\infty\)-categories.  The coherent
associativity and unit equivalences are those of derived fibre products.

The complete Segal nerve of this
\(\operatorname{Cat}_\infty\)-enriched category is denoted by the same
symbol and is regarded as the exceptional correspondence
\((\infty,2)\)-category.

Cartesian products of spectral schemes and spans give it a symmetric
monoidal structure.
\end{definition}

\begin{construction}[Action of a correspondence]
\label{constr:exceptional-correspondence-action}
To a correspondence
\[
    X\xleftarrow{u}K\xrightarrow{v}Y
\]
associate the colimit-preserving functor
\[
    v_!u^*:
    \SH(X)
    \longrightarrow
    \SH(Y).
\]
For composable correspondences
\[
    X\xleftarrow{u}K\xrightarrow{v}Y,
    \qquad
    Y\xleftarrow{u'}L\xrightarrow{v'}Z,
\]
put
\[
    M:=K\times_YL.
\]
Exceptional base change in the Cartesian square
\[
\begin{tikzcd}[column sep=large,row sep=large]
    M \arrow[r,"\widetilde u'"] \arrow[d,"\widetilde v"']
        & K \arrow[d,"v"] \\
    L \arrow[r,"u'"']
        & Y
\end{tikzcd}
\]
and exceptional composition give
\[
\begin{aligned}
    v'_!(u')^*v_!u^*
    &\xrightarrow{\simeq}
    v'_!\widetilde v_!(\widetilde u')^*u^* \\
    &\xrightarrow{\simeq}
    (v'\widetilde v)_!
    (u\widetilde u')^*.
\end{aligned}
\]
This is the comparison with the composite correspondence.

Now let
\[
    a:
    \bigl(
      X\xleftarrow{u'}K'\xrightarrow{v'}Y
    \bigr)
    \Longrightarrow
    \bigl(
      X\xleftarrow{u}K\xrightarrow{v}Y
    \bigr)
\]
be a \(2\)-morphism.  The map \(a\) belongs to \(\mathcal E\).  Indeed, the
graph
\[
    \Gamma_a:
    K\longrightarrow K\times_YK'
\]
is a section of the projection
\[
    K\times_YK'\longrightarrow K,
\]
which is the base change of \(v'\).  Hence \(\Gamma_a\) is a closed
immersion of finite presentation.  The projection
\[
    K\times_YK'\longrightarrow K'
\]
is a base change of \(v\), and their composite is \(a\).

Since \(a\) is proper, proper normalization identifies
\[
    a_!\simeq a_*.
\]
The unit of \(a^*\dashv a_*\) therefore gives a natural transformation
\[
\begin{aligned}
    v'_!u'^*
    &\longrightarrow
    v'_!a_*a^*u'^* \\
    &\simeq
    v'_!a_!u^* \\
    &\simeq
    v_!u^*.
\end{aligned}
\]
This is the action on \(2\)-morphisms.
\end{construction}

\begin{lemma}[Enriched pseudofunctor-extension criterion]
\label{lem:enriched-correspondence-extension}
Let \(\mathcal C\) and \(\mathcal D\) be
\(\operatorname{Cat}_\infty\)-enriched categories.  To define a
pseudofunctorial \((\infty,2)\)-functor
\[
    F:\mathcal C\longrightarrow\mathcal D
\]
it is enough to give the following data.

\begin{enumerate}
    \item An object \(F(X)\) of \(\mathcal D\) for every object \(X\) of
    \(\mathcal C\).

    \item For every pair \(X,Y\), a functor of hom \(\infty\)-categories
    \[
        F_{X,Y}:
        \mathcal C(X,Y)
        \longrightarrow
        \mathcal D(F(X),F(Y)).
    \]

    \item For every finite composable string
    \[
        X_n\xrightarrow{f_n}X_{n-1}
        \longrightarrow\cdots\longrightarrow
        X_1\xrightarrow{f_1}X_0
    \]
    in \(\mathcal C\), a natural equivalence
    \[
        F(f_1)\cdots F(f_n)
        \xrightarrow{\simeq}
        F(f_1\cdots f_n).
    \]

    \item Compatibility of these equivalences with every monotone map of
    finite ordinals, including faces, degeneracies, and subdivisions.
\end{enumerate}

The resulting \((\infty,2)\)-functor is normal when the comparison attached
to an identity string is the identity equivalence.
\end{lemma}

\begin{proof}
Apply the complete Segal nerve to the two
\(\operatorname{Cat}_\infty\)-enriched categories.  The \(n\)-simplices of
the complete Segal nerve encode strings of \(n\) composable
\(1\)-morphisms, morphisms between such strings, and the coherences among
their composites.

The object assignment and the functors on hom \(\infty\)-categories define
the degree-zero and degree-one parts of the desired simplicial map.  The
finite-string equivalences define its values on Segal composites in every
higher degree.  Compatibility with every monotone map of finite ordinals is
exactly the condition that these degreewise assignments commute with all
face and degeneracy operators.

Consequently, the data assemble to a map of complete Segal objects and
therefore to a pseudofunctorial \((\infty,2)\)-functor
\[
    F:\mathcal C\longrightarrow\mathcal D.
\]
The identity condition gives normality.
\end{proof}

\begin{proposition}[Correspondence extension]
\label{prop:exceptional-correspondence-extension}
For every pair \(X,Y\), Construction~\ref{constr:exceptional-correspondence-action}
defines an \(\infty\)-functor
\[
    F_{X,Y}:
    \operatorname{Corr}_{\mathcal E}(X,Y)
    \longrightarrow
    \operatorname{Fun}^L
    \bigl(
      \SH(X),
      \SH(Y)
    \bigr).
\]

For every finite composable string of correspondences, the iterated
exceptional Beck--Chevalley and exceptional-composition equivalences give
a natural equivalence from the composite of the assigned functors to the
functor assigned to the iterated derived fibre product.  These equivalences
are natural in proper maps of all apices and are compatible with every face,
degeneracy, and subdivision of the string.
\end{proposition}

\begin{proof}
Let
\[
    a:
    \bigl(
      X\xleftarrow{u'}K'\xrightarrow{v'}Y
    \bigr)
    \longrightarrow
    \bigl(
      X\xleftarrow{u}K\xrightarrow{v}Y
    \bigr)
\]
be a morphism in
\(\operatorname{Corr}_{\mathcal E}(X,Y)\).
Thus it is represented by a proper morphism
\[
    a:K\longrightarrow K'
\]
over \(X\times Y\).

We first verify that \(a\in\mathcal E\).  The graph
\[
    \Gamma_a:
    K\longrightarrow K\times_YK'
\]
is a section of the projection
\[
    K\times_YK'\longrightarrow K,
\]
which is the base change of the separated finitely presented morphism
\(v'\).  Hence \(\Gamma_a\) is a closed immersion of finite presentation.
The second projection
\[
    K\times_YK'\longrightarrow K'
\]
is a base change of \(v\), hence belongs to \(\mathcal E\).  Their composite
is \(a\), so \(a\in\mathcal E\).

Since \(a\) is proper, proper normalization gives
\[
    a_!\simeq a_*.
\]
The adjunction unit
\[
    \id\longrightarrow a_*a^*
\]
therefore gives
\[
\begin{aligned}
    v'_!u'^*
    &\longrightarrow
    v'_!a_*a^*u'^* \\
    &\simeq
    v'_!a_!u^* \\
    &\simeq
    v_!u^*.
\end{aligned}
\]
This assignment is natural in the mapping space of proper maps over
\(X\times Y\).

If
\[
    K\xrightarrow{a}K'\xrightarrow{b}K''
\]
are composable proper maps of apices, the unit for \(ba\) is the canonical
pasting of the units for \(a\) and \(b\):
\[
    \id
    \longrightarrow
    b_*b^*
    \longrightarrow
    b_*a_*a^*b^*.
\]
Compatibility of proper normalization and exceptional composition
identifies the resulting transformation with the composite of the
transformations assigned to \(a\) and \(b\).  Identity maps give identity
transformations.  Thus the construction defines the required functor on
each hom \(\infty\)-category.

Now consider a finite composable string
\[
    X_0
    \xleftarrow{u_1}
    K_1
    \xrightarrow{v_1}
    X_1
    \xleftarrow{u_2}
    K_2
    \xrightarrow{v_2}
    \cdots
    \xleftarrow{u_n}
    K_n
    \xrightarrow{v_n}
    X_n.
\]
Let
\[
    K_{[n]}
    :=
    K_1
    \times_{X_1}
    K_2
    \times_{X_2}\cdots
    \times_{X_{n-1}}
    K_n
\]
be its iterated derived fibre product.  Repeated exceptional base change
and exceptional composition give a canonical equivalence
\[
    v_{n!}u_n^*\cdots v_{1!}u_1^*
    \xrightarrow{\simeq}
    v_{[n]!}u_{[n]}^*,
\]
where
\[
    X_0
    \xleftarrow{u_{[n]}}
    K_{[n]}
    \xrightarrow{v_{[n]}}
    X_n
\]
is the composite correspondence.

A collection of proper maps of the apices induces a proper map of the
iterated fibre products.  Naturality of exceptional Beck--Chevalley in the
resulting Cartesian cubes and multiplicativity of the adjunction units
identify the comparison for the induced map with the corresponding
horizontal and vertical composites.

Restriction of the string along a monotone map of finite ordinals either
forgets a correspondence, inserts an identity, or combines adjacent
substrings.  The horizontal and vertical Beck--Chevalley pasting laws,
coherent exceptional composition, and identity normalization identify the
associated comparisons.  Hence the finite-string comparisons are
compatible with all faces, degeneracies, and subdivisions.
\end{proof}

\begin{theorem}[The exceptional correspondence functor]
\label{thm:exceptional-correspondence-functor}
Construction~\ref{constr:exceptional-correspondence-action} extends to a
normal pseudofunctorial \((\infty,2)\)-functor
\[
    \SH_{\mathrm{corr}}:
    \operatorname{Corr}_{\mathcal E}
    \bigl(
      \operatorname{Sch}^{\mathrm{sp}}_{\mathrm{qcqs}}
    \bigr)
    \longrightarrow
    \operatorname{Pr}^L_{\mathrm{st}},
\]
characterized on \(1\)-morphisms by
\[
    \SH_{\mathrm{corr}}
    \bigl(
      X\xleftarrow{u}K\xrightarrow{v}Y
    \bigr)
    =
    v_!u^*.
\]

It has a canonical lax symmetric monoidal refinement.  For every finite
family
\[
    X_1,\ldots,X_n
\]
its \(n\)-ary structural morphism is the colimit-preserving functor
\[
    \SH(X_1)
    \otimes_{\operatorname{Pr}^L}
    \cdots
    \otimes_{\operatorname{Pr}^L}
    \SH(X_n)
    \longrightarrow
    \SH(X_1\times\cdots\times X_n)
\]
induced by
\[
    (A_1,\ldots,A_n)
    \longmapsto
    \operatorname{pr}_1^*A_1
    \otimes\cdots\otimes
    \operatorname{pr}_n^*A_n.
\]
For the empty family, the structural unit is the colimit-preserving functor
\[
    \Sp
    \longrightarrow
    \SH(\operatorname{Spec}\mathbb S)
\]
carrying the sphere spectrum to the motivic unit.
\end{theorem}

\begin{proof}
The object assignment is
\[
    X\longmapsto\SH(X).
\]
Proposition~\ref{prop:exceptional-correspondence-extension} gives functors
on all hom \(\infty\)-categories.  It also gives, for every finite
composable string of correspondences, a natural equivalence from the
composite of the assigned functors to the functor assigned to the iterated
derived fibre product.  Those equivalences are compatible with all faces,
degeneracies, and subdivisions.

Lemma~\ref{lem:enriched-correspondence-extension} therefore first
produces a coherently unital pseudofunctorial \((\infty,2)\)-functor.  On
the identity correspondence of \(X\), its value is \((\id_X)_!\), and
Proposition~\ref{prop:exceptional-identity} gives a canonical equivalence
\[
    u_X:
    (\id_X)_!
    \xrightarrow{\simeq}
    \id_{\SH(X)}.
\]
The compatibility of \(u_X\) with left and right composition follows from
the degeneracy restrictions of the finite compactification-grid
coherences.  Normalize the pseudofunctor by replacing its value on every
identity correspondence by the literal identity functor and conjugating
the corresponding left- and right-unit comparisons by \(u_X\).  The
unitality and associativity coherences already proved for the unnormalized
pseudofunctor imply the normalized coherences.  The space of such
normalizations is contractible, because the space of compatible inverses to
the identity-evaluation equivalences is contractible.  We thereby obtain
the asserted normal pseudofunctorial \((\infty,2)\)-functor
\[
    \SH_{\mathrm{corr}}:
    \operatorname{Corr}_{\mathcal E}
    \bigl(
      \operatorname{Sch}^{\mathrm{sp}}_{\mathrm{qcqs}}
    \bigr)
    \longrightarrow
    \operatorname{Pr}^L_{\mathrm{st}}.
\]

We construct the lax symmetric monoidal refinement.  For every finite
nonempty set \(I\) and every family \((X_i)_{i\in I}\), the functor
\[
    \prod_{i\in I}\SH(X_i)
    \longrightarrow
    \SH\Bigl(\prod_{i\in I}X_i\Bigr),
\]
defined by
\[
    (A_i)_{i\in I}
    \longmapsto
    \bigotimes_{i\in I}\operatorname{pr}_i^*A_i,
\]
preserves colimits separately in every variable.  By the universal property
of the tensor product in \(\operatorname{Pr}^L\), it induces
\[
    \bigotimes_{i\in I}^{\operatorname{Pr}^L}\SH(X_i)
    \longrightarrow
    \SH\Bigl(\prod_{i\in I}X_i\Bigr).
\]

For the empty family, the monoidal unit of
\(\operatorname{Pr}^L_{\mathrm{st}}\) is \(\Sp\).  The motivic unit
\[
    \mathbf1_{\operatorname{Spec}\mathbb S}
    \in
    \SH(\operatorname{Spec}\mathbb S)
\]
therefore determines the colimit-preserving unit morphism
\[
    \Sp
    \longrightarrow
    \SH(\operatorname{Spec}\mathbb S)
\]
which sends the sphere spectrum to
\(\mathbf1_{\operatorname{Spec}\mathbb S}\).

Let
\[
    X_i\xleftarrow{u_i}K_i\xrightarrow{v_i}Y_i,
    \qquad
    i\in I,
\]
be a finite family of exceptional correspondences.  Repeated application
of Proposition~\ref{prop:exceptional-external-products} gives a canonical
equivalence
\[
    \boxtimes_{i\in I}
    \bigl(
      v_{i!}u_i^*
    \bigr)
    \xrightarrow{\simeq}
    \Bigl(
      \prod_{i\in I}v_i
    \Bigr)_!
    \Bigl(
      \prod_{i\in I}u_i
    \Bigr)^*.
\]
The associativity, symmetry, and unit coherences of this equivalence are
obtained from the corresponding coherences of exceptional external
products, exceptional composition, Cartesian products, and the symmetric
monoidal structures on pullback.

For proper maps of apices, naturality follows from proper base change,
proper normalization, and multiplicativity of the adjunction units
\[
    \id\longrightarrow a_*a^*.
\]
The finite-arity exterior-product equivalences are compatible with
bijections, insertions of the monoidal unit, partitions of finite sets, and
iterated Cartesian products.  These are precisely the coherence conditions
for a lax symmetric monoidal refinement.

No assertion that
\[
    \SH(X)\otimes_{\operatorname{Pr}^L}\SH(Y)
    \longrightarrow
    \SH(X\times Y)
\]
is an equivalence is used.
\end{proof}

\begin{remark}[Role of correspondences]
\label{rem:correspondences-role}
The correspondence category packages coherence already proved by
compactification grids and Beck--Chevalley pasting.  It is not used to
define \(f_!\).  This order avoids invoking the expression \(v_!u^*\) before
the exceptional direct image exists.
\end{remark}

\section{Closed purity transformations and pure coefficients}
\label{sec:six-closed-purity}

Let
\[
    i:Z\hookrightarrow X
\]
be a quasi-smooth closed immersion whose conormal module is finite locally
free.  Put
\[
    N_i^*:=N^*_{Z/X}=L_{Z/X}[-1],
    \qquad
    \nu_i=[N_i^*]\in K(Z).
\]
Under the coordinate-module convention, the inverse normal Thom twist is
\[
    \Sigma^{-\nu_i}.
\]

\subsection{The flat spectral affine line}

\begin{construction}[The flat spectral affine line]
\label{constr:flat-spectral-affine-line}
Let
\[
    \mathbb S[\mathbb N]
    :=
    \Sigma^\infty_+\mathbb N
\]
be the spherical monoid algebra of the additive monoid \(\mathbb N\).  For a
spectral scheme \(X\), define
\[
    \mathbb A^{1,\mathrm{fl}}_X
    :=
    \operatorname{Spec}_X
    \bigl(
      \mathcal O_X\otimes\mathbb S[\mathbb N]
    \bigr).
\]
The generator \(1\in\mathbb N\) determines a coordinate
\[
    t\in
    \pi_0
    \Gamma
    \bigl(
      \mathbb A^{1,\mathrm{fl}}_X,
      \mathcal O
    \bigr),
\]
a zero section
\[
    0_X:X\longrightarrow\mathbb A^{1,\mathrm{fl}}_X,
\]
and an open complement
\[
    \mathbb G^{\mathrm{fl}}_{m,X}
    :=
    \mathbb A^{1,\mathrm{fl}}_X\setminus0_X(X).
\]
This flat affine line is used only as the deformation parameter below.  It
is not the brave new differential affine line used as the motivic interval.
\end{construction}

\begin{proposition}[Basic properties of the flat affine line]
\label{prop:flat-affine-line-properties}
The construction of
\(\mathbb A^{1,\mathrm{fl}}_X\) commutes with arbitrary base change.  The
structural morphism
\[
    \mathbb A^{1,\mathrm{fl}}_X
    \longrightarrow
    X
\]
is affine, flat, separated, and quasi-compact.  If \(X\) is quasi-compact
and quasi-separated, then so is \(\mathbb A^{1,\mathrm{fl}}_X\).

Its classical truncation is the ordinary affine line:
\[
    \bigl(
      \mathbb A^{1,\mathrm{fl}}_X
    \bigr)_{\mathrm{cl}}
    \simeq
    \mathbb A^1_{X_{\mathrm{cl}}}.
\]
Likewise,
\[
    \bigl(
      \mathbb G^{\mathrm{fl}}_{m,X}
    \bigr)_{\mathrm{cl}}
    \simeq
    \mathbb G_{m,X_{\mathrm{cl}}}.
\]

If
\[
    T\longrightarrow\mathbb A^{1,\mathrm{fl}}_X
\]
is a spectral scheme and
\[
    T_0:=T\times_{\mathbb A^{1,\mathrm{fl}}_X}X,
\]
where \(X\to\mathbb A^{1,\mathrm{fl}}_X\) is the zero section, then there is
a canonical cofiber sequence in \(\operatorname{QCoh}(T)\)
\[
    \mathcal O_T
    \xrightarrow{t}
    \mathcal O_T
    \longrightarrow
    \mathcal O_{T_0}.
\]
\end{proposition}

\begin{proof}
All assertions are local on \(X\), so let
\[
    X=\operatorname{Spec}(A)
\]
with \(A\) connective, and put
\[
    B:=A\otimes\mathbb S[\mathbb N].
\]
For every morphism \(A\to A'\), there is a canonical equivalence
\[
\begin{aligned}
    B\otimes_AA'
    &\simeq
    \bigl(
      A\otimes\mathbb S[\mathbb N]
    \bigr)
    \otimes_AA' \\
    &\simeq
    A'\otimes\mathbb S[\mathbb N].
\end{aligned}
\]
This proves arbitrary base-change compatibility.

As an \(A\)-module, \(B\) has a canonical decomposition
\[
    B
    \simeq
    \bigoplus_{n\geq0}A\cdot t^n.
\]
It is therefore a coproduct of copies of \(A\), and hence is flat as an
\(A\)-module.  Consequently,
\[
    \operatorname{Spec}(B)
    \longrightarrow
    \operatorname{Spec}(A)
\]
is flat.  Since the morphism is affine, it is separated and quasi-compact.
If \(X\) is quasi-compact and quasi-separated, the same is therefore true
of \(\mathbb A^{1,\mathrm{fl}}_X\).

The decomposition of the underlying \(A\)-module gives
\[
    \pi_0(B)
    \simeq
    \pi_0(A)[t].
\]
Therefore
\[
    \bigl(
      \mathbb A^{1,\mathrm{fl}}_X
    \bigr)_{\mathrm{cl}}
    \simeq
    \mathbb A^1_{X_{\mathrm{cl}}}.
\]
Localizing at the coordinate \(t\) gives
\[
    \pi_0(B[t^{-1}])
    \simeq
    \pi_0(A)[t,t^{-1}],
\]
and hence
\[
    \bigl(
      \mathbb G^{\mathrm{fl}}_{m,X}
    \bigr)_{\mathrm{cl}}
    \simeq
    \mathbb G_{m,X_{\mathrm{cl}}}.
\]

Let
\[
    \epsilon_0:
    \mathbb S[\mathbb N]
    \longrightarrow
    \mathbb S
\]
be the zero-evaluation morphism determined by sending the generator
\(1\in\mathbb N\) to the zero element of \(\pi_0(\mathbb S)\).  Its base
change along \(\mathbb S\to A\) is the algebra morphism
\[
    B
    \longrightarrow
    A
\]
classifying the zero section.

Multiplication by \(t\) shifts the decomposition
\[
    \bigoplus_{n\geq0}A\cdot t^n
\]
one place to the right.  Its cofiber is therefore canonically the degree-zero
summand \(A\).  Thus there is a canonical cofiber sequence of \(B\)-modules
\[
    B
    \xrightarrow{t}
    B
    \longrightarrow
    A,
\]
where \(A\) is regarded as a \(B\)-module through zero evaluation.

Pulling this cofiber sequence back along
\[
    T\longrightarrow\operatorname{Spec}(B)
\]
gives
\[
    \mathcal O_T
    \xrightarrow{t}
    \mathcal O_T
    \longrightarrow
    \mathcal O_{T_0},
\]
as asserted.
\end{proof}

\subsection{Spectral deformation to the normal bundle}

\begin{construction}[The canonical deformation prestack]
\label{constr:spectral-deformation-prestack}
For a spectral scheme
\[
    T\longrightarrow X\times\mathbb A^{1,\mathrm{fl}},
\]
write
\[
    T_0
    :=
    T\times_{\mathbb A^{1,\mathrm{fl}}}0.
\]
Define a prestack
\[
    D_{Z/X}
    \longrightarrow
    X\times\mathbb A^{1,\mathrm{fl}}
\]
by the functor of points
\[
    \operatorname{Map}_{X\times\mathbb A^{1,\mathrm{fl}}}
    \bigl(
      T,
      D_{Z/X}
    \bigr)
    :=
    \operatorname{Map}_X(T_0,Z).
\]
\end{construction}

\begin{theorem}[Representability, finiteness, and fibres of the deformation]
\label{thm:spectral-deformation-normal-bundle}
The prestack \(D_{Z/X}\) is represented by a spectral scheme, and its
structural morphism
\[
    D_{Z/X}
    \longrightarrow
    X\times\mathbb A^{1,\mathrm{fl}}
\]
is separated and of finite presentation.  In particular,
\(D_{Z/X}\) is quasi-compact and quasi-separated.

It has a canonical quasi-smooth closed immersion
\[
    c_i:
    Z\times\mathbb A^{1,\mathrm{fl}}
    \longrightarrow
    D_{Z/X}
\]
whose conormal module is canonically equivalent to the pullback of
\(N_i^*\):
\[
    N^*_{c_i}
    \simeq
    \operatorname{pr}_Z^*N_i^*.
\]

There are canonical fibre identifications
\[
    D_{Z/X}
    \times_{\mathbb A^{1,\mathrm{fl}}}
    \mathbb G_m^{\mathrm{fl}}
    \simeq
    X\times\mathbb G_m^{\mathrm{fl}}
\]
and
\[
    D_{Z/X}
    \times_{\mathbb A^{1,\mathrm{fl}}}0
    \simeq
    \mathbb V_Z(N_i^*).
\]
On the generic fibre, \(c_i\) is identified with
\[
    i\times\id_{\mathbb G_m^{\mathrm{fl}}},
\]
and on the special fibre it is identified with the zero section of
\(\mathbb V_Z(N_i^*)\).

The construction, its structural morphism, the closed immersion \(c_i\),
its conormal identification, and both fibre identifications commute
coherently with arbitrary base change of \(i\).
\end{theorem}

\begin{proof}
Representability is local on \(X\).  By the local zero-locus theorem for
quasi-smooth closed immersions, there is a Zariski cover by opens
\(U\subseteq X\) on which
\[
    Z_U:=Z\times_XU
\]
is the derived zero locus of a morphism
\[
    s:E\longrightarrow\mathcal O_U
\]
with \(E\) finite locally free.

Put
\[
    P_U
    :=
    \mathbb V_U(E)\times\mathbb A^{1,\mathrm{fl}}.
\]
Let
\[
    \pi:P_U\longrightarrow U\times\mathbb A^{1,\mathrm{fl}}
\]
be the structural morphism, and let
\[
    v_{\mathrm{univ}}:
    \pi^*E\longrightarrow\mathcal O_{P_U}
\]
denote the universal linear coordinate.  Define
\[
    D_s
    :=
    Z
    \bigl(
      t v_{\mathrm{univ}}-\pi^*s
    \bigr)
    \subseteq
    P_U.
\]

Since \(\pi^*E\) is finite locally free, the morphism
\[
    D_s\longrightarrow P_U
\]
is a quasi-smooth closed immersion of finite presentation.  The morphism
\[
    P_U\longrightarrow U\times\mathbb A^{1,\mathrm{fl}}
\]
is affine, separated, smooth, and of finite presentation.  Consequently,
the composite
\[
    D_s\longrightarrow U\times\mathbb A^{1,\mathrm{fl}}
\]
is separated and of finite presentation.

Let
\[
    T\longrightarrow U\times\mathbb A^{1,\mathrm{fl}}
\]
be a test spectral scheme, and put
\[
    T_0
    :=
    T\times_{\mathbb A^{1,\mathrm{fl}}}0.
\]
A map
\[
    T\longrightarrow D_s
\]
is a pair consisting of a morphism
\[
    v:E_T\longrightarrow\mathcal O_T
\]
and a nullhomotopy
\[
    tv\simeq s_T.
\]
Since \(E_T\) is dualizable, applying
\[
    \operatorname{Map}_{\operatorname{QCoh}(T)}(E_T,-)
\]
to the cofiber sequence
\[
    \mathcal O_T
    \xrightarrow{t}
    \mathcal O_T
    \longrightarrow
    \mathcal O_{T_0}
\]
identifies the space of such pairs with the space of nullhomotopies of
\[
    s_T|_{T_0}:
    E_{T_0}\longrightarrow\mathcal O_{T_0}.
\]
The latter is precisely
\[
    \operatorname{Map}_U(T_0,Z_U).
\]
Thus \(D_s\) represents the restriction of the prestack
\(D_{Z/X}\) to \(U\).

Because the represented functor is intrinsic, the local representing
schemes agree on overlaps through a contractible space of equivalences.
They therefore glue by spectral Zariski descent to a spectral scheme
\[
    D_{Z/X}
    \longrightarrow
    X\times\mathbb A^{1,\mathrm{fl}}.
\]
Separatedness and finite presentation are local on the target and hence
descend from the local models.  Since
\(X\times\mathbb A^{1,\mathrm{fl}}\) is quasi-compact and
quasi-separated, a separated morphism of finite presentation over it has
quasi-compact and quasi-separated source.  Thus \(D_{Z/X}\) is qcqs.

We next construct the canonical centre.  Over \(U\), consider the square
\[
\begin{tikzcd}[column sep=large,row sep=large]
    Z_U\times\mathbb A^{1,\mathrm{fl}}
        \arrow[r]
        \arrow[d]
        &
    D_s
        \arrow[d]
    \\
    U\times\mathbb A^{1,\mathrm{fl}}
        \arrow[r,"0_E\times\id"']
        &
    \mathbb V_U(E)\times\mathbb A^{1,\mathrm{fl}},
\end{tikzcd}
\]
where \(0_E\) is the zero section.  This square is Cartesian: after
restricting the equation
\[
    t v_{\mathrm{univ}}-\pi^*s=0
\]
to the zero section \(v_{\mathrm{univ}}=0\), one obtains precisely the
derived zero-locus equation \(s=0\).

The lower horizontal morphism is a quasi-smooth closed immersion whose
conormal module is the pullback of \(E\).  Hence its base change
\[
    Z_U\times\mathbb A^{1,\mathrm{fl}}
    \longrightarrow
    D_s
\]
is a quasi-smooth closed immersion with conormal module
\[
    \operatorname{pr}_{Z_U}^*(E|_{Z_U}).
\]
Under the local zero-locus identification
\[
    E|_{Z_U}
    \simeq
    N_i^*|_{Z_U},
\]
these local conormal identifications agree on overlaps.  The local centres
therefore glue to a canonical quasi-smooth closed immersion
\[
    c_i:
    Z\times\mathbb A^{1,\mathrm{fl}}
    \longrightarrow
    D_{Z/X}
\]
with
\[
    N^*_{c_i}
    \simeq
    \operatorname{pr}_Z^*N_i^*.
\]

Over the open locus \(t\neq0\), multiplication by \(t\) is invertible and
the equation
\[
    tv=s
\]
has the unique solution
\[
    v=t^{-1}s.
\]
This gives the canonical equivalence
\[
    D_{Z/X}
    \times_{\mathbb A^{1,\mathrm{fl}}}
    \mathbb G_m^{\mathrm{fl}}
    \simeq
    X\times\mathbb G_m^{\mathrm{fl}}.
\]

Over \(t=0\), the equation consists of a nullhomotopy of \(s\), which
determines a point of \(Z_U\), together with an arbitrary linear coordinate
in \(E|_{Z_U}\).  Hence
\[
    D_s|_0
    \simeq
    \mathbb V_{Z_U}(E|_{Z_U}).
\]
These identifications glue through
\[
    E|_{Z_U}
    \simeq
    N_i^*|_{Z_U}
\]
to give
\[
    D_{Z/X}|_0
    \simeq
    \mathbb V_Z(N_i^*).
\]
The descriptions of \(c_i\) on the two fibres follow directly from its
local construction.

Finally, the functor-of-points definition, the local equation
\[
    tv=s,
\]
the Cartesian square defining \(c_i\), and the generic- and special-fibre
calculations all commute with extension of scalars.  They therefore give
the asserted coherent compatibility with arbitrary base change.
\end{proof}


\subsection{Bivariant classes and specialization}

\begin{definition}[Twisted bivariant spectra and classes]
\label{def:twisted-bivariant-classes}
Let
\[
    h:W\longrightarrow X
\]
belong to \(\mathcal E\), and let \(\xi\in K(W)\).  Define the twisted
bivariant mapping spectrum
\[
    \mathbf B(W/X,\xi)
    :=
    \operatorname{map}_{\SH(W)}
    \bigl(
      \operatorname{th}_W(\xi),
      h^!\mathbf1_X
    \bigr).
\]
Its underlying infinite-loop space is denoted
\[
    \mathbb B(W/X,\xi)
    :=
    \Omega^\infty\mathbf B(W/X,\xi).
\]
A twisted bivariant class is a point of
\(\mathbb B(W/X,\xi)\), equivalently a morphism
\[
    \operatorname{th}_W(\xi)
    \longrightarrow
    h^!\mathbf1_X
\]
in \(\SH(W)\).

The same notation with a subscript \({\mathrm{cl}}\) denotes the analogous
mapping spectrum and mapping space in ordinary stable motivic homotopy
theory.
\end{definition}

\begin{proposition}[Lax module structure on exceptional inverse image]
\label{prop:exceptional-inverse-module-structure}
Let
\[
    f:X\longrightarrow Y
\]
belong to \(\mathcal E\).  For
\[
    B,C\in\SH(Y),
\]
there is a canonical natural transformation
\[
    \lambda_f(B,C):
    f^!B\otimes_Xf^*C
    \longrightarrow
    f^!(B\otimes_YC).
\]
It is adjoint to the composite
\[
\begin{aligned}
    f_!
    \bigl(
      f^!B\otimes_Xf^*C
    \bigr)
    &\xrightarrow{\simeq}
    f_!f^!B\otimes_YC \\
    &\xrightarrow{\epsilon_f\otimes\id}
    B\otimes_YC,
\end{aligned}
\]
where the first arrow is the exceptional projection formula and
\(\epsilon_f\) is the counit of \(f_!\dashv f^!\).

These transformations are coherently unital and associative in the
coefficient variable.  For composable exceptional morphisms
\[
    X\xrightarrow{f}Y\xrightarrow{g}Z,
\]
the transformation for \(gf\) agrees, under
\[
    (gf)^!\simeq f^!g^!,
\]
with the pasted composite of the transformations for \(g\) and \(f\).
They are compatible with the canonical base-change comparison
transformations.

If \(C\) is strongly dualizable, then
\[
    \lambda_f(B,C):
    f^!B\otimes_Xf^*C
    \xrightarrow{\simeq}
    f^!(B\otimes_YC)
\]
is an equivalence.
\end{proposition}

\begin{proof}
Define \(\lambda_f(B,C)\) by adjunction from
\[
\begin{aligned}
    f_!
    \bigl(
      f^!B\otimes_Xf^*C
    \bigr)
    &\xrightarrow{\simeq}
    f_!f^!B\otimes_YC \\
    &\xrightarrow{\epsilon_f\otimes\id}
    B\otimes_YC.
\end{aligned}
\]
Naturality in \(B\) and \(C\) follows from naturality of the exceptional
projection formula and the counit.

The unit and associativity diagrams for \(\lambda_f\) are the adjoint mates
of the corresponding unit and associativity diagrams for the strong
\(\SH(Y)\)-module structure on \(f_!\).  They therefore commute.

For composable
\[
    X\xrightarrow{f}Y\xrightarrow{g}Z,
\]
the composition compatibility is the mate of the compatibility of the
exceptional projection formula with the composition equivalence
\[
    g_!f_!
    \xrightarrow{\simeq}
    (gf)_!.
\]
The corresponding base-change compatibility is the mate of the
compatibility between exceptional base change and the exceptional
projection formula.

If \(C\) is strongly dualizable, the asserted invertibility is
Corollary~\ref{cor:exceptional-dualizable-coefficient}.
\end{proof}

\begin{proposition}[Composition product]
\label{prop:bivariant-composition-product}
Let
\[
    V\xrightarrow{a}W\xrightarrow{b}X
\]
belong to \(\mathcal E\), and let
\[
    \eta\in K(V),
    \qquad
    \xi\in K(W).
\]
There is a canonical pairing of spectra
\[
    \mathbf B(V/W,\eta)
    \otimes
    \mathbf B(W/X,\xi)
    \longrightarrow
    \mathbf B
    \bigl(
      V/X,
      \eta+a^*\xi
    \bigr).
\]
It is coherently associative and unital and is compatible with pullback,
Thom additivity, and stable classical comparison.

On underlying infinite-loop spaces, this pairing sends
\[
    \alpha\in\mathbb B(V/W,\eta),
    \qquad
    \beta\in\mathbb B(W/X,\xi)
\]
to a class
\[
    \alpha\cdot\beta
    \in
    \mathbb B
    \bigl(
      V/X,
      \eta+a^*\xi
    \bigr).
\]
\end{proposition}

\begin{proof}
The spectral enrichment of the stable coefficient categories gives a
canonical pairing
\[
\begin{aligned}
    &
    \operatorname{map}_V
    \bigl(
      \operatorname{th}_V(\eta),
      a^!\mathbf1_W
    \bigr)
    \otimes
    \operatorname{map}_W
    \bigl(
      \operatorname{th}_W(\xi),
      b^!\mathbf1_X
    \bigr)
    \\
    &\qquad\longrightarrow
    \operatorname{map}_V
    \Bigl(
      \operatorname{th}_V(\eta)
      \otimes_V
      a^*\operatorname{th}_W(\xi),
      a^!\mathbf1_W
      \otimes_V
      a^*b^!\mathbf1_X
    \Bigr).
\end{aligned}
\]
Here the second mapping spectrum is first pulled back along \(a^*\), and
the two resulting maps in \(\SH(V)\) are then tensored.

Apply the lax module transformation of
Proposition~\ref{prop:exceptional-inverse-module-structure}
with
\[
    B=\mathbf1_W,
    \qquad
    C=b^!\mathbf1_X.
\]
It gives
\[
\begin{aligned}
    a^!\mathbf1_W
    \otimes_V
    a^*b^!\mathbf1_X
    &\longrightarrow
    a^!
    \bigl(
      \mathbf1_W\otimes_Wb^!\mathbf1_X
    \bigr) \\
    &\simeq
    a^!b^!\mathbf1_X.
\end{aligned}
\]
Exceptional inverse-image composition then gives
\[
    a^!b^!\mathbf1_X
    \xrightarrow{\simeq}
    (ba)^!\mathbf1_X.
\]

Pullback naturality and coherent Thom additivity give
\[
\begin{aligned}
    \operatorname{th}_V(\eta)
    \otimes_V
    a^*\operatorname{th}_W(\xi)
    &\xrightarrow{\simeq}
    \operatorname{th}_V(\eta)
    \otimes_V
    \operatorname{th}_V(a^*\xi) \\
    &\xrightarrow{\simeq}
    \operatorname{th}_V
    \bigl(
      \eta+a^*\xi
    \bigr).
\end{aligned}
\]
Combining these maps defines the pairing
\[
    \mathbf B(V/W,\eta)
    \otimes
    \mathbf B(W/X,\xi)
    \longrightarrow
    \mathbf B
    \bigl(
      V/X,
      \eta+a^*\xi
    \bigr).
\]

Associativity follows from coherent associativity of:

\begin{enumerate}
    \item the tensor product and spectral enrichment;

    \item the lax module transformations
    \(\lambda_f\) of
    Proposition~\ref{prop:exceptional-inverse-module-structure};

    \item exceptional inverse-image composition;

    \item Thom additivity.
\end{enumerate}

For the identity morphism, the unit bivariant class is represented by
\[
    \mathbf1_W
    \xrightarrow{\id}
    (\id_W)^!\mathbf1_W
    \simeq
    \mathbf1_W.
\]
The unit constraints for the tensor product, the lax module structure, and
exceptional inverse-image composition identify multiplication by this
class with the identity.  Hence the pairing is coherently unital.

Compatibility with pullback follows from compatibility of the lax module
transformations with the pullback--exceptional exchange transformations,
together with pullback compatibility of exceptional inverse-image
composition and Thom twists.  Compatibility with stable classical
comparison follows because comparison preserves the tensor product,
exceptional inverse images, their adjunction data, the projection formula,
and Thom additivity.

Applying \(\Omega^\infty\) and evaluating the pairing on points sends
\[
    \alpha\in\mathbb B(V/W,\eta),
    \qquad
    \beta\in\mathbb B(W/X,\xi)
\]
to the asserted product
\[
    \alpha\cdot\beta
    \in
    \mathbb B
    \bigl(
      V/X,
      \eta+a^*\xi
    \bigr).
\]
\end{proof}

\begin{theorem}[Specialization attached to a deformation datum]
\label{thm:specialization-deformation-datum}
Work in ordinary stable motivic homotopy theory.  Let
\[
    q:
    D\longrightarrow
    X\times\mathbb A^1
\]
be a separated morphism of finite presentation with quasi-compact
quasi-separated source.  Let
\[
    N\xrightarrow{k}D\xleftarrow{j}D_\eta
\]
be the complementary closed and open fibres over \(0\) and
\(\mathbb G_m\), respectively.

Suppose given a specified equivalence
\[
    \alpha_\eta:
    D_\eta
    \xrightarrow{\simeq}
    X\times\mathbb G_m
\]
over \(X\times\mathbb G_m\), where the map
\[
    D_\eta\longrightarrow X\times\mathbb G_m
\]
is the restriction of \(q\).

Suppose also given a finite locally free module \(E\) on a
quasi-compact quasi-separated scheme \(Z\), a separated morphism of finite
presentation
\[
    i:Z\longrightarrow X,
\]
and a specified equivalence
\[
    \alpha_0:
    N
    \xrightarrow{\simeq}
    \mathbb V_Z(E)
\]
over \(X\), where the structural morphism on the right-hand side is
\[
    \mathbb V_Z(E)
    \xrightarrow{\pi}
    Z
    \xrightarrow{i}
    X.
\]

Then the structural morphisms
\[
    D\longrightarrow X,
    \qquad
    D_\eta\longrightarrow X,
    \qquad
    N\longrightarrow X,
    \qquad
    Z\longrightarrow X
\]
belong to \(\mathcal E\), and there is a canonical specialization class
\[
    \eta_D:
    \operatorname{th}_Z(-E)
    \longrightarrow
    i^!\mathbf1_X.
\]

For every Cartesian pullback of the full deformation datum along a
morphism
\[
    g:X'\longrightarrow X,
\]
write
\[
    g_Z:Z'\longrightarrow Z,
    \qquad
    i':Z'\longrightarrow X'
\]
for the induced morphisms and write \(D'\) for the pulled-back deformation
datum.  Under the canonical Thom base-change equivalence
\[
    g_Z^*\operatorname{th}_Z(-E)
    \simeq
    \operatorname{th}_{Z'}(-g_Z^*E),
\]
the specialization classes satisfy
\[
    \operatorname{Ex}^{!*}_{i,g}(\mathbf1_X)
    \circ
    g_Z^*(\eta_D)
    =
    \eta_{D'}.
\]
The class is normalized so that the standard deformation of a zero section
gives the inverse Thom class.
\end{theorem}

\begin{proof}
The projection
\[
    X\times\mathbb A^1
    \longrightarrow
    X
\]
is separated and of finite presentation.  Hence the composite
\[
    d:
    D
    \xrightarrow{q}
    X\times\mathbb A^1
    \longrightarrow
    X
\]
belongs to \(\mathcal E\).  Its restriction to the quasi-compact open
\(D_\eta\) and to the closed fibre \(N\) also belongs to
\(\mathcal E\).  The morphism \(i\) belongs to \(\mathcal E\) by
assumption.

Apply \(d_!\) to the localization triangle
\[
    j_!\mathbf1_{D_\eta}
    \longrightarrow
    \mathbf1_D
    \longrightarrow
    k_*\mathbf1_N
\]
and then apply the mapping-spectrum functor
\[
    \operatorname{map}_{\SH_{\mathrm{cl}}(X)}
    \bigl(
      -,
      \mathbf1_X
    \bigr).
\]
Using exceptional composition, open normalization, and closed
normalization, one obtains a fibre sequence of bivariant mapping spectra
\[
    \mathbf B_{\mathrm{cl}}(N/X,0)
    \longrightarrow
    \mathbf B_{\mathrm{cl}}(D/X,0)
    \longrightarrow
    \mathbf B_{\mathrm{cl}}(D_\eta/X,0).
\]
Its connecting morphism is a canonical map of spectra
\[
    \partial_D:
    \Omega
    \mathbf B_{\mathrm{cl}}(D_\eta/X,0)
    \longrightarrow
    \mathbf B_{\mathrm{cl}}(N/X,0).
\]

For the universal closed--open pair
\[
    X
    \longrightarrow
    X\times\mathbb A^1
    \longleftarrow
    X\times\mathbb G_m,
\]
the coordinate
\[
    t\in\Gamma(\mathbb G_m,\mathcal O^\times)
\]
determines a canonical morphism of spectra
\[
    \gamma_t:
    \mathbf B_{\mathrm{cl}}(X/X,0)
    \longrightarrow
    \Omega
    \mathbf B_{\mathrm{cl}}
    (X\times\mathbb G_m/X,0)
\]
which is a section of the localization boundary.  This is the coordinate
section of
\cite[Paragraph~3.2.2]{DegliseJinKhanFundamentalClasses}.

Because
\[
    \alpha_\eta:
    D_\eta
    \xrightarrow{\simeq}
    X\times\mathbb G_m
\]
is an equivalence over \(X\times\mathbb G_m\), it preserves the structural
map to \(X\) and the distinguished invertible coordinate \(t\).  It
therefore identifies the target of \(\gamma_t\) canonically with
\[
    \Omega\mathbf B_{\mathrm{cl}}(D_\eta/X,0).
\]
Using this identification, define
\[
    \sigma_D
    \in
    \mathbb B_{\mathrm{cl}}(N/X,0)
\]
by
\[
    \sigma_D
    :=
    \Omega^\infty
    \bigl(
      \partial_D\circ\gamma_t
    \bigr)
    \bigl(
      \id_{\mathbf1_X}
    \bigr).
\]

The specified equivalence
\[
    \alpha_0:
    N
    \xrightarrow{\simeq}
    \mathbb V_Z(E)
\]
over \(X\) identifies the structural morphism \(N\to X\) with
\[
    \mathbb V_Z(E)
    \xrightarrow{\pi}
    Z
    \xrightarrow{i}
    X.
\]
There is consequently a canonical equivalence of bivariant mapping spectra
\[
    \boldsymbol{\vartheta}_{\pi,E}:
    \mathbf B_{\mathrm{cl}}(N/X,0)
    \xrightarrow{\simeq}
    \mathbf B_{\mathrm{cl}}(Z/X,-E).
\]

Indeed, smooth purity for the vector-bundle projection gives
\[
    \pi^!
    \simeq
    \Sigma^{\pi^*E}\pi^*,
\]
and vector-bundle homotopy invariance makes
\[
    \pi^*:
    \SH_{\mathrm{cl}}(Z)
    \longrightarrow
    \SH_{\mathrm{cl}}(N)
\]
an equivalence.  Therefore
\[
\begin{aligned}
    \mathbf B_{\mathrm{cl}}(N/X,0)
    &=
    \operatorname{map}_N
    \bigl(
      \mathbf1_N,
      \pi^!i^!\mathbf1_X
    \bigr) \\
    &\simeq
    \operatorname{map}_N
    \bigl(
      \operatorname{th}_N(-\pi^*E),
      \pi^*i^!\mathbf1_X
    \bigr) \\
    &\simeq
    \operatorname{map}_Z
    \bigl(
      \operatorname{th}_Z(-E),
      i^!\mathbf1_X
    \bigr) \\
    &=
    \mathbf B_{\mathrm{cl}}(Z/X,-E).
\end{aligned}
\]

Define
\[
    \eta_D
    \in
    \mathbb B_{\mathrm{cl}}(Z/X,-E)
\]
by
\[
    \eta_D
    :=
    \Omega^\infty
    \boldsymbol{\vartheta}_{\pi,E}
    (\sigma_D).
\]
Equivalently, \(\eta_D\) is the corresponding morphism
\[
    \operatorname{th}_Z(-E)
    \longrightarrow
    i^!\mathbf1_X.
\]

Let
\[
    g:X'\longrightarrow X
\]
be a morphism, and pull back the full deformation datum, including the
specified generic- and special-fibre identifications.  Write
\[
    g_Z:Z'\longrightarrow Z,
    \qquad
    i':Z'\longrightarrow X'
\]
for the induced morphisms and write \(D'\) for the pulled-back deformation
datum.  Exceptional base change identifies the pullback of the localization
fibre sequence for \(D\) with the localization fibre sequence for \(D'\),
with the target comparison given by
\[
    \operatorname{Ex}^{!*}_{i,g}(\mathbf1_X):
    g_Z^*i^!\mathbf1_X
    \longrightarrow
    i'^!\mathbf1_{X'}.
\]
The pullback of \(\alpha_\eta\) remains an equivalence over
\(X'\times\mathbb G_m\) and preserves the coordinate \(t\), so it preserves
the coordinate section \(\gamma_t\).  The pullback of \(\alpha_0\)
identifies the pulled-back special fibre with
\(\mathbb V_{Z'}(g_Z^*E)\).  Smooth purity for the vector-bundle projection,
vector-bundle homotopy invariance, and the inverse Thom equivalence commute
with this pullback.  Therefore, under
\[
    g_Z^*\operatorname{th}_Z(-E)
    \simeq
    \operatorname{th}_{Z'}(-g_Z^*E),
\]
one has
\[
    \operatorname{Ex}^{!*}_{i,g}(\mathbf1_X)
    \circ
    g_Z^*(\eta_D)
    =
    \eta_{D'}.
\]

The zero-section normalization is proved in
Lemma~\ref{lem:specialization-zero-section-normalization}.
\end{proof}


\begin{lemma}[Zero-section normalization of specialization]
\label{lem:specialization-zero-section-normalization}
Let \(E\) be finite locally free on \(Z\), let
\[
    \pi:\mathbb V_Z(E)\longrightarrow Z
\]
be the vector-bundle projection, and let
\[
    s_E:Z\hookrightarrow\mathbb V_Z(E)
\]
be the zero section.  For the standard deformation datum of \(s_E\), the
specialization class of
Theorem~\ref{thm:specialization-deformation-datum} is the inverse Thom
class
\[
    \operatorname{th}_Z(-E)
    \xrightarrow{\simeq}
    s_E^!\mathbf1_{\mathbb V_Z(E)}.
\]
\end{lemma}

\begin{proof}
Write \(X:=\mathbb V_Z(E)\).  The standard deformation of the zero section
is locally described inside
\[
    X_x\times_ZX_v\times\mathbb A^1
\]
by the equation
\[
    x=tv.
\]
The morphism
\[
    X_v\times\mathbb A^1
    \longrightarrow
    D,
    \qquad
    (v,t)\longmapsto(tv,v,t),
\]
is an isomorphism.  Under this identification, the generic-fibre
identification
\[
    D_\eta\simeq X\times\mathbb G_m
\]
is scalar multiplication by \(t\), while the special fibre is
\[
    N\simeq X_v=\mathbb V_Z(E),
\]
and its structural morphism to \(X\) is
\[
    X_v\xrightarrow{\pi}Z\xrightarrow{s_E}X.
\]
Thus the deformation datum, its coordinate section \(\gamma_t\), and its
localization fibre sequence are the standard universal zero-section datum
used to normalize specialization.

For that datum, the coordinate section is normalized so that the
localization boundary is the fundamental class of the zero fibre.  Under
exceptional composition for \(s_E\pi\), smooth purity for the vector-bundle
projection \(\pi\), and vector-bundle homotopy invariance, the equivalence
\[
    \boldsymbol{\vartheta}_{\pi,E}:
    \mathbf B_{\mathrm{cl}}(N/X,0)
    \xrightarrow{\simeq}
    \mathbf B_{\mathrm{cl}}(Z/X,-E)
\]
identifies this zero-fibre class with the inverse Thom class of \(E\).
Equivalently, the image of
\[
    \partial_D\gamma_t(\id_{\mathbf1_X})
\]
under \(\boldsymbol{\vartheta}_{\pi,E}\) is the morphism
\[
    \operatorname{th}_Z(-E)
    \xrightarrow{\simeq}
    s_E^!\mathbf1_X.
\]
This is the zero-section normalization of the ordinary specialization
construction; the preceding equation verifies that the present deformation,
coordinate, conormal, and Thom conventions are exactly those of
\cite[Section~3.2]{DegliseJinKhanFundamentalClasses}.  Hence the class
constructed in Theorem~\ref{thm:specialization-deformation-datum} is the
inverse Thom class.
\end{proof}

\subsection{The virtual fundamental class}

Let
\[
    \iota_Z:Z_{\mathrm{cl}}\longrightarrow Z
\]
denote classical truncation and put
\[
    N_i^0
    :=
    \mathbb V_{Z_{\mathrm{cl}}}
    \bigl(
      \iota_Z^*N_i^*
    \bigr).
\]
Classical truncation of
Theorem~\ref{thm:spectral-deformation-normal-bundle} gives an ordinary
deformation datum
\[
    X_{\mathrm{cl}}\times\mathbb G_m
    \hookrightarrow
    (D_{Z/X})_{\mathrm{cl}}
    \hookleftarrow
    N_i^0.
\]
The vector bundle \(N_i^0\) records the virtual normal bundle supplied by
the spectral enhancement.  It need not be the ordinary normal bundle of the
closed immersion \(i_{\mathrm{cl}}\).

\begin{construction}[The virtual fundamental class]
\label{constr:virtual-closed-fundamental-class}
Apply Theorem~\ref{thm:specialization-deformation-datum} to the ordinary
deformation datum \((D_{Z/X})_{\mathrm{cl}}\).  It gives a canonical class
\[
    \eta^{\mathrm{vir}}_{i,\mathrm{cl}}:
    \operatorname{th}_{Z_{\mathrm{cl}}}
    \bigl(
      -\iota_Z^*\nu_i
    \bigr)
    \longrightarrow
    i_{\mathrm{cl}}^!\mathbf1_{X_{\mathrm{cl}}}.
\]
Stable comparison identifies the source and target with the images of
\[
    \operatorname{th}_Z(-\nu_i)
    \qquad\text{and}\qquad
    i^!\mathbf1_X.
\]
Define
\[
    \eta_i:
    \operatorname{th}_Z(-\nu_i)
    \longrightarrow
    i^!\mathbf1_X
\]
to be the unique morphism carried to
\(\eta^{\mathrm{vir}}_{i,\mathrm{cl}}\).
\end{construction}

\begin{construction}[Closed-purity transformation]
\label{constr:absolute-closed-purity-transformation}
The closed projection formula and the counit
\[
    i_*i^!B\longrightarrow B
\]
induce, by adjunction, a canonical module transformation
\[
    i^!B\otimes_Zi^*A
    \longrightarrow
    i^!(B\otimes_XA).
\]
Taking \(B=\mathbf1_X\), define
\[
    \mathfrak q_i(A):
    \Sigma^{-\nu_i}i^*A
    \longrightarrow
    i^!A
\]
as the composite
\[
\begin{aligned}
    \operatorname{th}_Z(-\nu_i)
    \otimes_Zi^*A
    &\xrightarrow{\eta_i\otimes\id}
    i^!\mathbf1_X\otimes_Zi^*A \\
    &\longrightarrow
    i^!A.
\end{aligned}
\]
This is natural in \(A\).
\end{construction}

\begin{definition}[Purity for an object]
\label{def:i-pure-object}
Let \(i:Z\hookrightarrow X\) be as above.  An object
\[
    A\in\SH(X)
\]
is \emph{pure for \(i\)} if the evaluation of the closed-purity
transformation at \(A\),
\[
    \mathfrak q_i(A):
    \Sigma^{-\nu_i}i^*A
    \longrightarrow
    i^!A,
\]
is an equivalence.

A full tensor ideal \(\mathcal C\subseteq\SH(X)\) is \emph{absolutely pure
for \(i\)} if every object of \(\mathcal C\) is pure for \(i\).
\end{definition}

\begin{lemma}[Quasi-smoothness of the successive deformation centres]
\label{lem:iterated-deformation-centres}
Let
\[
    Z_r\xrightarrow{i_r}Z_{r-1}
    \longrightarrow\cdots\longrightarrow
    Z_1\xrightarrow{i_1}Z_0
\]
be a finite chain of quasi-smooth closed immersions with finite locally free
conormal modules.

Suppose that the \((r-1)\)-fold deformation
\[
    D^{(r-1)}
    \longrightarrow
    Z_0\times
    (\mathbb A^{1,\mathrm{fl}})^{r-1}
\]
has been constructed, together with its canonical centre
\[
    c_{r-1}:
    Z_{r-1}\times
    (\mathbb A^{1,\mathrm{fl}})^{r-1}
    \longrightarrow
    D^{(r-1)}.
\]
Then the composite
\[
\begin{aligned}
    C_r
    &:=
    Z_r\times
    (\mathbb A^{1,\mathrm{fl}})^{r-1}
    \\
    &\xrightarrow{i_r\times\id}
    Z_{r-1}\times
    (\mathbb A^{1,\mathrm{fl}})^{r-1}
    \xrightarrow{c_{r-1}}
    D^{(r-1)}
\end{aligned}
\]
is a quasi-smooth closed immersion with finite locally free conormal
module.

Its conormal module fits into a canonical fibre sequence
\[
    (i_r\times\id)^*N^*_{c_{r-1}}
    \longrightarrow
    N^*_{C_r/D^{(r-1)}}
    \longrightarrow
    \operatorname{pr}_{Z_r}^*N^*_{i_r}.
\]
Consequently, in \(K(C_r)\),
\[
    [N^*_{C_r/D^{(r-1)}}]
    =
    (i_r\times\id)^*[N^*_{c_{r-1}}]
    +
    \operatorname{pr}_{Z_r}^*\nu_{i_r}.
\]
\end{lemma}

\begin{proof}
The morphism
\[
    i_r\times\id
\]
is a base change of \(i_r\), hence is a quasi-smooth closed immersion with
finite locally free conormal module
\[
    \operatorname{pr}_{Z_r}^*N^*_{i_r}.
\]
The morphism \(c_{r-1}\) is quasi-smooth with finite locally free conormal
module by
Theorem~\ref{thm:spectral-deformation-normal-bundle} and induction.

Quasi-smooth closed immersions are stable under composition.  Cotangent
transitivity for the composite gives a fibre sequence
\[
    (i_r\times\id)^*
    L_{Z_{r-1}\times(\mathbb A^{1,\mathrm{fl}})^{r-1}/D^{(r-1)}}
    \longrightarrow
    L_{C_r/D^{(r-1)}}
    \longrightarrow
    L_{C_r/
      Z_{r-1}\times(\mathbb A^{1,\mathrm{fl}})^{r-1}}.
\]
Shifting by \([-1]\) gives the asserted conormal fibre sequence.

All three terms are finite locally free.  Locally the fibre sequence splits,
so the middle term is finite locally free.  Passing to \(K\)-theory gives
the displayed additivity formula.
\end{proof}

\begin{construction}[Iterated deformation]
\label{constr:iterated-closed-deformation}
Let
\[
    Z_n\xrightarrow{i_n}Z_{n-1}
    \longrightarrow\cdots\longrightarrow
    Z_1\xrightarrow{i_1}Z_0
\]
be a finite chain of quasi-smooth closed immersions with finite locally free
conormal modules.

Define the iterated deformation inductively.  Put
\[
    D^{(1)}
    :=
    D_{Z_1/Z_0}.
\]
Suppose that
\[
    D^{(r-1)}
    \longrightarrow
    Z_0\times
    (\mathbb A^{1,\mathrm{fl}})^{r-1}
\]
has been constructed.  Let
\[
    C_r
    :=
    Z_r\times
    (\mathbb A^{1,\mathrm{fl}})^{r-1}
    \longrightarrow
    D^{(r-1)}
\]
be the quasi-smooth closed immersion of
Lemma~\ref{lem:iterated-deformation-centres}.  Define
\[
    D^{(r)}
    :=
    D_{C_r/D^{(r-1)}}.
\]
Composing its structural morphism with that of \(D^{(r-1)}\) gives
\[
    D^{(r)}
    \longrightarrow
    Z_0\times
    (\mathbb A^{1,\mathrm{fl}})^r.
\]

Write
\[
    t_1,\ldots,t_r
\]
for the deformation coordinates.  The canonical centre of the final
deformation step gives a closed immersion
\[
    Z_r\times
    (\mathbb A^{1,\mathrm{fl}})^r
    \longrightarrow
    D^{(r)}.
\]
\end{construction}

\begin{lemma}[Coherent interchange of iterated boundaries]
\label{lem:coherent-iterated-boundaries}
Let
\[
    F:
    \bigl(
      (\Delta^2)^{\mathrm{op}}
    \bigr)^r
    \longrightarrow
    \Sp
\]
be a functor whose restriction to every coordinate, with all other
coordinates fixed, is a fibre sequence
\[
    F(\ldots,2,\ldots)
    \longrightarrow
    F(\ldots,1,\ldots)
    \longrightarrow
    F(\ldots,0,\ldots).
\]

For every ordering
\[
    \sigma=(\sigma_1,\ldots,\sigma_r)
\]
of the coordinates, successive application of the connecting morphisms
defines an iterated boundary
\[
    \partial_\sigma:
    \Omega^rF(0,\ldots,0)
    \longrightarrow
    F(2,\ldots,2).
\]

For every permutation \(\tau\in\Sigma_r\), there is a canonical homotopy
between
\[
    \partial_{\sigma\tau}
\]
and the composite of \(\partial_\sigma\) with the symmetry equivalence of
the \(r\) loop coordinates induced by \(\tau\).  These homotopies satisfy
the unit, braid, and all higher permutation coherences.

The construction is natural in \(F\) and is compatible with restriction,
insertion of constant coordinates, and permutation of coordinates.
\end{lemma}

\begin{proof}
Let
\[
    \operatorname{FibSeq}(\Sp)
    \subseteq
    \operatorname{Fun}
    \bigl(
      (\Delta^2)^{\mathrm{op}},
      \Sp
    \bigr)
\]
denote the full \(\infty\)-category of fibre sequences.  There is a
canonical natural connecting transformation
\[
    \partial:
    \Omega\operatorname{ev}_0
    \longrightarrow
    \operatorname{ev}_2
\]
on \(\operatorname{FibSeq}(\Sp)\).

Fixing all coordinates except the \(a\)-th makes \(F\) a functor from the
remaining coordinates into \(\operatorname{FibSeq}(\Sp)\).  Applying
\(\partial\) in the \(a\)-th coordinate therefore gives a natural
transformation
\[
    \partial_a:
    \Omega F(\ldots,0,\ldots)
    \longrightarrow
    F(\ldots,2,\ldots).
\]
Composing these transformations in an ordering
\(\sigma=(\sigma_1,\ldots,\sigma_r)\) gives
\[
    \partial_\sigma:
    \Omega^rF(0,\ldots,0)
    \longrightarrow
    F(2,\ldots,2).
\]

It remains to compare two adjacent orders.  Restrict \(F\) to the two
coordinates being interchanged.  This gives a biexact diagram
\[
    \bigl(
      (\Delta^2)^{\mathrm{op}}
    \bigr)^2
    \longrightarrow
    \Sp.
\]
Naturality of the connecting transformation in each variable gives a
canonical comparison between the composites
\[
    \partial_a\partial_b
    \quad\text{and}\quad
    \partial_b\partial_a.
\]
Under the canonical identification
\[
    \Omega^2E
    \simeq
    S^{-1}\otimes S^{-1}\otimes E,
\]
this comparison is exactly the symmetry interchanging the two copies of
\(S^{-1}\).  Thus the two orders agree after applying the canonical
symmetry of the corresponding loop coordinates.

Every permutation is generated by adjacent transpositions.  The symmetric
monoidal coherence of \(\Sp\) supplies the braid, involutivity, and all
higher coherences among the associated symmetry equivalences.  Naturality
of the connecting transformation supplies the corresponding coherent
homotopies among the iterated boundaries.

Restriction, insertion of a constant coordinate, and permutation of
coordinates are functorial operations on the product indexing category.
The construction of the transformations \(\partial_a\) is natural under
each of them.  This proves the final compatibility assertion.
\end{proof}

\begin{proposition}[Faces, composite deformation, and the localization cube]
\label{prop:iterated-deformation-faces}
The iterated deformation of
Construction~\ref{constr:iterated-closed-deformation} has the following
properties.

\begin{enumerate}
    \item For every \(r\), the structural morphism
    \[
        D^{(r)}
        \longrightarrow
        Z_0\times
        (\mathbb A^{1,\mathrm{fl}})^r
    \]
    is separated and of finite presentation, and \(D^{(r)}\) is
    quasi-compact and quasi-separated.  The construction commutes
    coherently with arbitrary base change of the chain.

    \item Inverting the last coordinate removes the last deformation step:
    \[
        D^{(r)}|_{t_r\neq0}
        \simeq
        D^{(r-1)}
        \times
        \mathbb G_m^{\mathrm{fl}}.
    \]
    Restricting to \(t_r=0\) gives
    \[
        D^{(r)}|_{t_r=0}
        \simeq
        \mathbb V_{C_r}
        \bigl(
          N^*_{C_r/D^{(r-1)}}
        \bigr).
    \]

    \item On the locus
    \[
        t_1\neq0,\ldots,t_{r-1}\neq0,
    \]
    there is a canonical identification
    \[
        D^{(r)}
        \simeq
        D_{Z_r/Z_0}
        \times
        (\mathbb G_m^{\mathrm{fl}})^{r-1}.
    \]
    Under this identification, the remaining coordinate \(t_r\) is the
    deformation coordinate for the composite immersion
    \[
        Z_r\longrightarrow Z_0.
    \]

    \item The fibre of \(D^{(r)}\) at the origin is an iterated tower of
    vector bundles over \(Z_r\).  Its successive coordinate modules are
    the terms of the conormal filtration obtained from cotangent
    transitivity for
    \[
        Z_r\longrightarrow Z_{r-1}
        \longrightarrow\cdots\longrightarrow Z_0.
    \]
    Coherent Thom additivity gives a canonical equivalence between the
    Thom object of this tower and
    \[
        \operatorname{th}_{Z_r}
        \bigl(
          \nu_{Z_r/Z_0}
        \bigr),
    \]
    where
    \[
        \nu_{Z_r/Z_0}
        =
        \nu_{i_r}
        +
        i_r^*\nu_{i_{r-1}}
        +\cdots+
        (i_2\cdots i_r)^*\nu_{i_1}.
    \]

    \item Put
    \[
        D_r:=(D^{(r)})_{\mathrm{cl}}
    \]
    and let
    \[
        q_r:D_r\longrightarrow (Z_0)_{\mathrm{cl}}
    \]
    be its structural morphism.  For each coordinate \(t_a\), let
    \[
        D_{r,a}^0\xrightarrow{\iota_a}D_r
        \xleftarrow{\jmath_a}D_{r,a}^{\times}
    \]
    be the complementary closed--open pair defined by
    \(t_a=0\) and \(t_a\neq0\).

    For \(\epsilon\in\{0,1,2\}\), define endofunctors of
    \(\SH_{\mathrm{cl}}(D_r)\) by
    \[
        \mathcal F_{a,0}
        :=
        \jmath_{a!}\jmath_a^*,
        \qquad
        \mathcal F_{a,1}
        :=
        \id,
        \qquad
        \mathcal F_{a,2}
        :=
        \iota_{a*}\iota_a^*.
    \]
    The localization transformations
    \[
        \mathcal F_{a,0}
        \longrightarrow
        \mathcal F_{a,1}
        \longrightarrow
        \mathcal F_{a,2}
    \]
    and the Beck--Chevalley equivalences for distinct coordinates assemble
    into a canonical functor
    \[
        \mathcal Q_r:
        (\Delta^2)^r
        \longrightarrow
        \operatorname{Fun}^{\mathrm{ex}}
        \bigl(
          \SH_{\mathrm{cl}}(D_r),
          \SH_{\mathrm{cl}}(D_r)
        \bigr)
    \]
    whose restriction to every coordinate is a cofiber sequence of exact
    endofunctors.

    Applying \(\mathcal Q_r\) to \(\mathbf1_{D_r}\), then applying
    \((q_r)_!\), and finally applying
    \[
        \operatorname{map}_{\SH_{\mathrm{cl}}((Z_0)_{\mathrm{cl}})}
        \bigl(
          -,
          \mathbf1_{(Z_0)_{\mathrm{cl}}}
        \bigr)
    \]
    gives a canonical multiexact diagram
    \[
        \mathcal L_r:
        \bigl(
          (\Delta^2)^{\mathrm{op}}
        \bigr)^r
        \longrightarrow
        \Sp.
    \]
    Its restriction to every coordinate is a fibre sequence of spectra,
    and its extreme vertices are the bivariant mapping spectra of the
    corresponding coordinate strata.

    The coordinate sections
    \[
        \gamma_{t_1},\ldots,\gamma_{t_r}
    \]
    define a compatible system of points in the looped open vertices of
    \(\mathcal L_r\).  Iterated application of the boundary maps along any
    coordinate flag gives the corresponding iterated specialization class.

    Every two orderings of the coordinate boundary operators are connected
    by the canonical coherent homotopies supplied by the higher simplices of
    \(\mathcal L_r\).  For adjacent coordinates, interchange of the two
    localization boundaries contributes the stable symmetry of the two
    suspension coordinates, and interchange of the two coordinate sections
    contributes the same symmetry.  The two symmetries cancel in the
    specialization operators.

    Consequently, every ordering of the coordinate specializations in this
    fixed multiexact cube gives the same specialization class after the
    canonical symmetry and Thom-additivity identifications.  This assertion
    concerns interchange of boundary operators inside \(\mathcal L_r\); it
    does not identify coordinate panels with simplicial face or degeneracy
    constructions.
\end{enumerate}
\end{proposition}

\begin{proof}
Part~(1) follows by induction from
Theorem~\ref{thm:spectral-deformation-normal-bundle} and
Lemma~\ref{lem:iterated-deformation-centres}.  At every stage, the new
deformation morphism is separated and of finite presentation over the
preceding deformation space times the flat affine line.  Composition
preserves these properties.  Since the base is quasi-compact and
quasi-separated, each \(D^{(r)}\) is quasi-compact and quasi-separated.
Base-change compatibility is inherited inductively from
Theorem~\ref{thm:spectral-deformation-normal-bundle}.

Part~(2) is the generic- and special-fibre description of
Theorem~\ref{thm:spectral-deformation-normal-bundle}, applied to the
immersion
\[
    C_r\longrightarrow D^{(r-1)}.
\]

We prove part~(3) by induction on \(r\).  For \(r=1\), it is the defining
generic-fibre description of \(D_{Z_1/Z_0}\).  Suppose the assertion holds
through stage \(r-1\).  On the locus
\[
    t_1\neq0,\ldots,t_{r-1}\neq0,
\]
all preceding deformation steps are on their generic fibres, so the
preceding deformation space is canonically
\[
    Z_0\times
    (\mathbb G_m^{\mathrm{fl}})^{r-1}.
\]
Under this identification, the centre \(C_r\) becomes
\[
    Z_r\times
    (\mathbb G_m^{\mathrm{fl}})^{r-1}
    \longrightarrow
    Z_0\times
    (\mathbb G_m^{\mathrm{fl}})^{r-1}.
\]
Base-change compatibility of deformation therefore gives
\[
    D^{(r)}
    \simeq
    D_{Z_r/Z_0}
    \times
    (\mathbb G_m^{\mathrm{fl}})^{r-1}
\]
on this locus.  Under this equivalence, the final coordinate \(t_r\) is the
deformation coordinate for the composite immersion
\[
    Z_r\longrightarrow Z_0.
\]
This proves part~(3).

At the origin, part~(2) replaces the final centre by its conormal vector
bundle.  Iterating this description gives a tower of vector bundles over
\(Z_r\).  Lemma~\ref{lem:iterated-deformation-centres} identifies the
successive coordinate modules with the successive extensions in the
cotangent-transitivity filtration.  Coherent Thom additivity identifies the
Thom object of the tower with
\[
    \operatorname{th}_{Z_r}
    \bigl(
      \nu_{Z_r/Z_0}
    \bigr).
\]
This proves part~(4).

We now construct the multiexact diagram of part~(5).  For every coordinate
\(a\), stable localization gives a cofiber sequence of exact endofunctors
of \(\SH_{\mathrm{cl}}(D_r)\):
\[
    \jmath_{a!}\jmath_a^*
    \longrightarrow
    \id
    \longrightarrow
    \iota_{a*}\iota_a^*.
\]
For two distinct coordinates \(a\) and \(b\), all squares formed from the
corresponding closed and open coordinate strata are Cartesian.  Open base
change, closed base change, and open--closed exchange therefore give
canonical equivalences
\[
    \mathcal F_{a,\epsilon}
    \mathcal F_{b,\delta}
    \xrightarrow{\simeq}
    \mathcal F_{b,\delta}
    \mathcal F_{a,\epsilon},
    \qquad
    \epsilon,\delta\in\{0,1,2\}.
\]
The Beck--Chevalley pasting laws imply that these equivalences satisfy the
unit, braid, and all higher interchange coherences.

Consequently, the \(r\) localization diagrams assemble into a functor
\[
    \mathcal Q_r:
    (\Delta^2)^r
    \longrightarrow
    \operatorname{Fun}^{\mathrm{ex}}
    \bigl(
      \SH_{\mathrm{cl}}(D_r),
      \SH_{\mathrm{cl}}(D_r)
    \bigr)
\]
whose restriction to every coordinate is the corresponding localization
cofiber sequence.

Apply \(\mathcal Q_r\) to \(\mathbf1_{D_r}\), apply the exact functor
\[
    (q_r)_!,
\]
and then apply the contravariant mapping-spectrum functor
\[
    \operatorname{map}_{\SH_{\mathrm{cl}}((Z_0)_{\mathrm{cl}})}
    \bigl(
      -,
      \mathbf1_{(Z_0)_{\mathrm{cl}}}
    \bigr).
\]
The result is a functor
\[
    \mathcal L_r:
    \bigl(
      (\Delta^2)^{\mathrm{op}}
    \bigr)^r
    \longrightarrow
    \Sp
\]
whose restriction to every coordinate is a fibre sequence.

Exceptional composition and the open and closed normalizations identify
the extreme vertices of \(\mathcal L_r\) with the bivariant mapping spectra
of the corresponding coordinate strata.  Under these identifications, the
connecting morphism in the \(a\)-th coordinate is precisely the
localization boundary used in specialization with respect to \(t_a\).

For each coordinate, the universal invertible function \(t_a\) determines
the coordinate section
\[
    \gamma_{t_a}
\]
at the corresponding open vertex.  Functoriality of invertible functions
makes these sections compatible with restriction, base change, and
permutation of coordinates.

Choose an ordering
\[
    \sigma=(\sigma_1,\ldots,\sigma_r)
\]
of the coordinates.  The ordered collection of coordinate sections gives
a point in the iterated loop spectrum of the all-open vertex.  Applying the
iterated boundary
\[
    \partial_\sigma:
    \Omega^r
    \mathcal L_r(0,\ldots,0)
    \longrightarrow
    \mathcal L_r(2,\ldots,2)
\]
produces the specialization class associated with the coordinate flag
\(\sigma\).

Lemma~\ref{lem:coherent-iterated-boundaries} gives canonical coherent
comparisons among the iterated boundaries associated with different
orders.  For an adjacent interchange of coordinates \(a\) and \(b\), the
comparison between the two boundary composites is the symmetry
interchanging the corresponding two loop coordinates.

Interchanging the coordinate sections
\[
    \gamma_{t_a}
    \quad\text{and}\quad
    \gamma_{t_b}
\]
contributes the same stable symmetry.  Consequently, the symmetry arising
from interchange of the two boundary operators is cancelled by the
symmetry arising from interchange of the two coordinate sections.  The
resulting specialization class is therefore independent of the ordering
after the canonical symmetry and Thom-additivity identifications.

The unit, braid, and all higher permutation coherences for these boundary
interchanges follow from the corresponding coherences in
Lemma~\ref{lem:coherent-iterated-boundaries}
and from the symmetric monoidal coherence of the coordinate-section
product.  No comparison with simplicial face or degeneracy constructions is
used here.
\end{proof}


\begin{proposition}[Ordered-partition open faces]
\label{prop:iterated-deformation-ordered-partitions}
Let
\[
    Z_n\xrightarrow{i_n}Z_{n-1}
    \longrightarrow\cdots\longrightarrow
    Z_1\xrightarrow{i_1}Z_0
\]
be a finite chain of quasi-smooth closed immersions with finite locally free
conormal modules, and let
\[
    D^{(n)}
    \longrightarrow
    Z_0\times(\mathbb A^{1,\mathrm{fl}})^n
\]
be its iterated deformation.  Let
\[
    P:
    0=a_0<a_1<\cdots<a_m=n
\]
be an ordered partition.  For \(1\leq b\leq m\), write
\[
    h_b:
    Z_{a_b}\longrightarrow Z_{a_{b-1}}
\]
for the composite of the immersions in the \(b\)-th consecutive block, and
let \(D_P^{(m)}\) be the iterated deformation of the block-composite chain
\[
    Z_n\xrightarrow{h_m}Z_{a_{m-1}}
    \longrightarrow\cdots\longrightarrow
    Z_{a_1}\xrightarrow{h_1}Z_0.
\]

Let \(U_P\) be the locus in
\((\mathbb A^{1,\mathrm{fl}})^n\) on which
\[
    t_j\neq0
    \qquad
    \text{for every }
    j\notin\{a_1,\ldots,a_m\}.
\]
Then there is a canonical equivalence over \(U_P\)
\[
    D^{(n)}|_{U_P}
    \xrightarrow{\simeq}
    D_P^{(m)}
    \times
    (\mathbb G_m^{\mathrm{fl}})^{n-m},
\]
where the remaining deformation coordinates of \(D_P^{(m)}\) are
\[
    t_{a_1},\ldots,t_{a_m}.
\]
These equivalences commute coherently with arbitrary base change of the
chain and with refinement of ordered partitions.

\end{proposition}

\begin{proof}
We first prove the geometric equivalence.  Within the first block, invert
\[
    t_1,\ldots,t_{a_1-1}.
\]
Repeated application of
Proposition~\ref{prop:iterated-deformation-faces}(2), followed by the
composite-deformation identification of
Proposition~\ref{prop:iterated-deformation-faces}(3), identifies the
restriction of the first \(a_1\) deformation stages with
\[
    D_{Z_{a_1}/Z_0}
    \times
    (\mathbb G_m^{\mathrm{fl}})^{a_1-1},
\]
with \(t_{a_1}\) as the deformation coordinate for the composite
\(h_1\).

The next deformation stage is formed from the canonical centre in the
preceding deformation space.  The deformation construction and its
canonical centre commute with arbitrary base change.  Consequently, after
restriction to the preceding open locus, the stages in the second block are
the base changes of the deformation stages for the composite
\[
    h_2:Z_{a_2}\longrightarrow Z_{a_1}.
\]
Inverting the coordinates internal to that block identifies those stages
with
\[
    D_{Z_{a_2}/Z_{a_1}}
    \times
    (\mathbb G_m^{\mathrm{fl}})^{a_2-a_1-1}
\]
over the already constructed first block.  Continuing block by block gives
the canonical equivalence
\[
    D^{(n)}|_{U_P}
    \simeq
    D_P^{(m)}
    \times
    (\mathbb G_m^{\mathrm{fl}})^{n-m}.
\]

Every comparison used in this construction is a generic-fibre
identification or a base-change equivalence for the deformation and its
canonical centre.  Those comparisons are natural and satisfy their usual
pasting laws.  Therefore the resulting equivalence commutes coherently with
base change of the chain.  If \(P'\) refines \(P\), restricting the
\(P'\)-equivalence to the locus \(U_P\) performs the same generic-fibre and
base-change identifications in smaller consecutive groups.  The pasting
laws identify that restriction with the \(P\)-equivalence, giving coherent
compatibility with refinement.

\end{proof}


\begin{lemma}[Coordinate cancellation]
\label{lem:coordinate-cancellation}
Work in ordinary stable motivic homotopy theory.  Let
\[
    h:W\longrightarrow X
\]
belong to \(\mathcal E\), and let \(\xi\in K(W)\).  Let
\[
    \gamma_t^{W/W}:
    \mathbf B_{\mathrm{cl}}(W/W,0)
    \longrightarrow
    \Omega\mathbf B_{\mathrm{cl}}
    (W\times\mathbb G_m/W,0)
\]
be the coordinate section for the universal complementary pair
\[
    W\times\{0\}
    \longrightarrow
    W\times\mathbb A^1
    \longleftarrow
    W\times\mathbb G_m,
\]
and let
\[
    \{t\}_W
    :=
    \Omega^\infty\gamma_t^{W/W}(\mathbf1_W)
    \in
    \Omega^\infty\Omega\mathbf B_{\mathrm{cl}}
    (W\times\mathbb G_m/W,0)
\]
be its coordinate class.  Bivariant composition with this class defines a
coordinate-suspension morphism
\[
    \gamma_t^W:
    \mathbf B_{\mathrm{cl}}(W/X,\xi)
    \longrightarrow
    \Omega\mathbf B_{\mathrm{cl}}
    (W\times\mathbb G_m/X,\operatorname{pr}_W^*\xi).
\]
Let
\[
    \partial_t^W:
    \Omega\mathbf B_{\mathrm{cl}}
    (W\times\mathbb G_m/X,\operatorname{pr}_W^*\xi)
    \longrightarrow
    \mathbf B_{\mathrm{cl}}(W/X,\xi)
\]
be the localization boundary.  Then there is a canonical homotopy
\[
    \partial_t^W\gamma_t^W
    \simeq
    \id_{\mathbf B_{\mathrm{cl}}(W/X,\xi)}.
\]
The homotopy is natural in \(h\), \(\xi\), and Cartesian pullback.
\end{lemma}

\begin{proof}
The bivariant composition product is constructed from the tensor product,
the lax module structure on exceptional inverse image, exceptional
inverse-image composition, and Thom additivity.  Each of these operations
is exact in every stable variable.  Multiplication by a fixed bivariant
class therefore gives a morphism of the localization fibre sequences.
Naturality of connecting morphisms identifies the composite
\[
    \partial_t^W\gamma_t^W
\]
with multiplication by the boundary of the coordinate class.  By the
defining normalization of the coordinate section,
\[
    \partial_t^W(\{t\}_W)
    =
    \mathbf1_W
\]
in the underlying infinite-loop space of
\(\mathbf B_{\mathrm{cl}}(W/W,0)\).  The unit property of the bivariant
composition product therefore gives a canonical homotopy of morphisms of
spectra
\[
    \partial_t^W\gamma_t^W
    \simeq
    \id_{\mathbf B_{\mathrm{cl}}(W/X,\xi)}.
\]
All constructions used in this calculation commute with Cartesian
pullback, which gives the stated naturality.
\end{proof}


\begin{lemma}[Multiplicativity of successive specialization]
\label{lem:successive-specialization-product}
Let
\[
    W\xrightarrow{k}Z\xrightarrow{i}X
\]
be quasi-smooth closed immersions with finite locally free conormal modules,
and let \(D^{(2)}\) be the two-parameter iterated deformation of this
chain.  After classical truncation and the canonical Thom-additivity
identification, the coordinate flag that first specializes the deformation
of \(i\) and then specializes the induced deformation of \(k\) represents
the bivariant product
\[
    \eta_k\cdot\eta_i
    \in
    \mathbb B
    \bigl(
      W/X,
      -\nu_k-k^*\nu_i
    \bigr).
\]
The same assertion holds for a finite chain: successive coordinate
specialization represents the iterated bivariant product, with the
specified order and the canonical Thom-additivity identifications.
\end{lemma}

\begin{proof}
The bivariant product of
Proposition~\ref{prop:bivariant-composition-product} is constructed from
the tensor product, the canonical module transformation
\[
    a^!B\otimes a^*C
    \longrightarrow
    a^!(B\otimes C),
\]
exceptional inverse-image composition, and Thom additivity.  Each of these
operations is exact in every stable variable.

Consider the localization cofiber sequence for one deformation coordinate.
Tensoring it with a fixed coefficient, applying exceptional direct image,
and then applying the contravariant mapping-spectrum functor into the unit
produces a morphism of fibre sequences.  Naturality of connecting morphisms
therefore gives compatibility of the localization boundary with the
bivariant product.  Explicitly, if \(\partial_a\) is the boundary in the
\(a\)-th coordinate, \(\alpha\) is a bivariant class independent of that
coordinate, and \(\beta\) is a class on the corresponding open stratum,
then
\[
    \partial_a(\beta\cdot\alpha)
    \simeq
    \partial_a(\beta)\cdot\alpha,
\]
after the canonical symmetry moving the suspension coordinate of the
boundary past the suspension degree of \(\alpha\).

Apply this to the two coordinate fibre sequences in
\[
    \mathcal L_2:
    \bigl((\Delta^2)^{\mathrm{op}}\bigr)^2
    \longrightarrow
    \Sp.
\]
Retain both coordinate sections at the all-open vertex.  Taking the first
boundary, corresponding to the deformation of \(i\), and then applying the
first inverse Thom equivalence gives
\[
    \eta_i.
\]
The second coordinate section remains as the coordinate class for the
induced deformation of \(k\).  Taking the second boundary in the bivariant
theory already carrying \(\eta_i\), and applying the second inverse Thom
equivalence, gives multiplication by \(\eta_k\).  Thus the ordered
successive specialization is
\[
    \eta_k\cdot\eta_i.
\]
The two inverse Thom equivalences combine by coherent Thom additivity to
identify the source with
\[
    \operatorname{th}_W(-\nu_k-k^*\nu_i).
\]

Interchanging the two connecting morphisms contributes the stable symmetry
of their suspension coordinates.  Interchanging the two coordinate
sections contributes the same symmetry.  The two contributions therefore
cancel in the specialization operator, exactly as encoded by the
two-dimensional faces of \(\mathcal L_2\).

For a finite chain, retain all coordinate sections at the all-open vertex
and apply the boundaries in the specified order.  Boundary compatibility
with the bivariant product identifies each new specialization with
multiplication by the fundamental class of the corresponding immersion.
Associativity of the module transformation, exceptional inverse-image
composition, and coherent Thom additivity identify the resulting class
with the specified iterated bivariant product.
\end{proof}

\begin{theorem}[Composition of virtual fundamental classes]
\label{thm:virtual-class-composition-coherence}
For composable quasi-smooth closed immersions
\[
    W\xrightarrow{k}Z\xrightarrow{i}X
\]
with finite locally free conormal modules, cotangent transitivity gives a
canonical fibre sequence
\[
    k^*N_i^*
    \longrightarrow
    N_{ik}^*
    \longrightarrow
    N_k^*
\]
and hence a canonical path
\[
    \nu_{ik}
    \simeq
    \nu_k+k^*\nu_i.
\]
Under this path and exceptional inverse-image composition, the double
deformation construction gives a canonical path
\[
    \eta_{ik}
    \simeq
    \eta_k\cdot\eta_i
\]
in
\[
    \mathbb B
    \bigl(
      W/X,
      -\nu_{ik}
    \bigr).
\]

For every finite chain
\[
    Z_n\xrightarrow{i_n}Z_{n-1}
    \longrightarrow\cdots\longrightarrow
    Z_1\xrightarrow{i_1}Z_0
\]
of quasi-smooth closed immersions with finite locally free conormal
modules, the virtual fundamental class of the composite and every
parenthesized iterated bivariant product of the classes
\[
    \eta_{i_n},\ldots,\eta_{i_1}
\]
represent the same element of
\[
    \pi_0\mathbf B
    \bigl(
      Z_n/Z_0,
      -\nu_{Z_n/Z_0}
    \bigr)
\]
under the canonical cotangent-transitivity and Thom-additivity
identifications.  For an identity immersion, the virtual fundamental class
is the bivariant unit.  No fully simplicial or homotopy-coherent Gysin
refinement is asserted.
\end{theorem}

\begin{proof}
Apply
Construction~\ref{constr:iterated-closed-deformation}
to the chain
\[
    W\xrightarrow{k}Z\xrightarrow{i}X.
\]
The resulting two-parameter deformation
\[
    D^{(2)}
    \longrightarrow
    X\times
    (\mathbb A^{1,\mathrm{fl}})^2
\]
is separated and of finite presentation by
Proposition~\ref{prop:iterated-deformation-faces}.

After classical truncation,
Proposition~\ref{prop:iterated-deformation-faces}(5)
gives a biexact diagram
\[
    \mathcal L_2:
    \bigl(
      (\Delta^2)^{\mathrm{op}}
    \bigr)^2
    \longrightarrow
    \Sp.
\]
Let \(\Gamma_{12}\) denote the product of the two coordinate sections at
the all-open vertex.  By
Lemma~\ref{lem:successive-specialization-product}, applying first the
boundary for the deformation of \(i\) and then the boundary for the induced
deformation of \(k\), followed by the two inverse Thom equivalences, gives
\[
    \eta^{\mathrm{vir}}_{k,\mathrm{cl}}
    \cdot
    \eta^{\mathrm{vir}}_{i,\mathrm{cl}}.
\]

On the open panel where the first deformation coordinate \(t_1\) is
invertible,
Proposition~\ref{prop:iterated-deformation-faces}(3)
identifies the remaining family with
\[
    D_{W/X}\times\mathbb G_{m,t_1}^{\mathrm{fl}}.
\]
Under this identification, \(t_2\) is the deformation coordinate for the
composite immersion
\[
    ik:W\longrightarrow X,
\]
while \(t_1\) is the coordinate on the residual trivial
\(\mathbb G_m\)-factor.  Let
\[
    \sigma_{ik}
    \in
    \mathbb B_{\mathrm{cl}}
    (N_{ik}^0/X_{\mathrm{cl}},0)
\]
be the specialization class for the one-parameter deformation of the
composite before applying the inverse Thom equivalence.  Naturality of the
localization boundary and of the coordinate sections under the panel
identification gives
\[
    \partial_{t_2}(\Gamma_{12})
    \simeq
    \gamma_{t_1}^{N_{ik}^0}(\sigma_{ik}).
\]
Thus specialization in \(t_2\) does not discard the residual
\(\mathbb G_m\)-factor: it gives the coordinate suspension of
\(\sigma_{ik}\).  Applying the remaining boundary and using
Lemma~\ref{lem:coordinate-cancellation} gives
\[
\begin{aligned}
    \partial_{t_1}\partial_{t_2}(\Gamma_{12})
    &\simeq
    \partial_{t_1}
    \gamma_{t_1}^{N_{ik}^0}(\sigma_{ik}) \\
    &\simeq
    \sigma_{ik}.
\end{aligned}
\]
After the inverse Thom equivalence, this is precisely
\[
    \eta^{\mathrm{vir}}_{ik,\mathrm{cl}}.
\]

The two constructions are obtained from the same biexact localization
diagram.  The comparison of the two boundary factorizations contributes
the symmetry of the two loop coordinates.  Interchanging the two
coordinate sections contributes the same symmetry, so the two symmetries
cancel.  The resulting two-dimensional comparison gives a canonical path
\[
    \eta^{\mathrm{vir}}_{ik,\mathrm{cl}}
    \simeq
    \eta^{\mathrm{vir}}_{k,\mathrm{cl}}
    \cdot
    \eta^{\mathrm{vir}}_{i,\mathrm{cl}}
\]
in
\[
    \mathbb B_{\mathrm{cl}}
    \bigl(
      W_{\mathrm{cl}}/X_{\mathrm{cl}},
      -\iota_W^*\nu_{ik}
    \bigr).
\]

At the origin, the conormal filtration is
\[
    k^*N_i^*
    \longrightarrow
    N_{ik}^*
    \longrightarrow
    N_k^*.
\]
Coherent Thom additivity identifies the source of the common specialization
map by
\[
    \operatorname{th}_{W_{\mathrm{cl}}}
    \bigl(
      -\iota_W^*\nu_k
      -
      \iota_W^*k^*\nu_i
    \bigr)
    \xrightarrow{\simeq}
    \operatorname{th}_{W_{\mathrm{cl}}}
    \bigl(
      -\iota_W^*\nu_{ik}
    \bigr).
\]

Stable comparison identifies the spectral and classical bivariant mapping
spectra, virtual fundamental classes, bivariant products, exceptional
inverse-image composition equivalences, and Thom-additivity paths.
Transporting the preceding path therefore gives
\[
    \eta_{ik}
    \simeq
    \eta_k\cdot\eta_i
\]
in
\[
    \mathbb B
    \bigl(
      W/X,
      -\nu_{ik}
    \bigr).
\]

For a finite chain, the ordered successive specializations give the
corresponding iterated bivariant product by
Lemma~\ref{lem:successive-specialization-product}.  On the panel on which
all but the final deformation coordinate are invertible,
Proposition~\ref{prop:iterated-deformation-faces}(3) identifies the family
with the deformation of the full composite times
\((\mathbb G_m^{\mathrm{fl}})^{n-1}\).  Specialization in the final
coordinate gives the specialization class of the composite with the
remaining \(n-1\) coordinate suspensions still present.  Repeated
application of Lemma~\ref{lem:coordinate-cancellation} removes those
suspensions one at a time.  The coherent interchange of boundaries and
coordinate sections in \(\mathcal L_n\) identifies this class with every
ordered successive specialization.  Coherent associativity of the
bivariant product and coherent Thom additivity then identify every
parenthesization with the same element of \(\pi_0\).  For the identity
immersion, the deformation is the trivial deformation, its conormal class
is zero, and coordinate cancellation gives the bivariant unit.
\end{proof}

\begin{definition}[Excess square and refined pullback]
\label{def:closed-excess-square}
A commutative square
\[
\begin{tikzcd}[column sep=large,row sep=large]
    Z' \arrow[r,"g'"] \arrow[d,"i'"']
        & Z \arrow[d,"i"] \\
    X' \arrow[r,"g"']
        & X
\end{tikzcd}
\]
whose vertical morphisms \(i\) and \(i'\) are quasi-smooth closed
immersions is an \emph{excess square} if it is Cartesian on classical
truncations and the fibre \(\xi\) of the canonical
conormal comparison
\[
    (g')^*N_i^*
    \longrightarrow
    N_{i'}^*
\]
is finite locally free.  Thus there is a fibre sequence
\[
    \xi
    \longrightarrow
    (g')^*N_i^*
    \longrightarrow
    N_{i'}^*
\]
and a coherent path
\[
    (g')^*\nu_i
    \simeq
    \xi+\nu_{i'}.
\]

The ordinary change-of-base morphism for the classically Cartesian square,
transported through stable comparison, gives a canonical morphism
\[
    \operatorname{Ex}^{!*}_\Delta(\mathbf1_X):
    (g')^*i^!\mathbf1_X
    \longrightarrow
    i'^!\mathbf1_{X'}.
\]
Define the refined pullback of the fundamental class by
\[
    \Delta^\star\eta_i
    :=
    \operatorname{Ex}^{!*}_\Delta(\mathbf1_X)
    \circ
    (g')^*\eta_i.
\]
For a derived Cartesian square, this morphism agrees with the canonical
exchange transformation of
Construction~\ref{constr:pullback-exceptional-exchange}.
\end{definition}

\begin{lemma}[Affine extended-Rees calculation]
\label{lem:affine-extended-rees-deformation}
Let
\[
    A\longrightarrow B
\]
be a morphism of connective commutative ring spectra which presents a
quasi-smooth closed immersion with finite locally free conormal module, and
put
\[
    I_{B/A}:=\operatorname{fib}(A\longrightarrow B).
\]
Write
\[
    D_{B/A}
    :=
    D_{\operatorname{Spec}(B)/\operatorname{Spec}(A)}
    =
    \operatorname{Spec}
    \bigl(
      R^{\mathrm{ext}}_{B/A}
    \bigr).
\]
The scaling action on the deformation coordinate makes
\(R^{\mathrm{ext}}_{B/A}\) a \(\mathbb Z\)-graded
\(A[t]\)-algebra, with \(t\) in degree \(-1\).  Its nonpositive and
nonnegative parts satisfy
\[
    \bigl(
      R^{\mathrm{ext}}_{B/A}
    \bigr)_{\leq0}
    \simeq
    A[t],
    \qquad
    \bigl(
      R^{\mathrm{ext}}_{B/A}
    \bigr)_{\geq0}
    =:
    R_{B/A}.
\]
The degree-one component of the unextended Rees algebra is canonically
\[
    (R_{B/A})_1
    \simeq
    I_{B/A},
\]
and the ordinary graded ring \(\pi_0(R_{B/A})\) is generated over
\(\pi_0(A)\) by its degree-one component.  These identifications are
natural in commutative triangles of connective commutative ring spectra
which are surjective on \(\pi_0\).
\end{lemma}

\begin{proof}
The assertions are Zariski-local on \(\operatorname{Spec}(A)\).  By the
local zero-locus theorem, after a Zariski localization there are a finite
free \(A\)-module \(E\) and a morphism
\[
    s:E\longrightarrow A
\]
such that, writing \(A_s\) and \(A_0\) for the two
\(\operatorname{Sym}_A(E)\)-algebra structures classified by \(s\) and by
the zero morphism, respectively,
\[
    B
    \simeq
    A_s\otimes_{\operatorname{Sym}_A(E)}A_0.
\]
Write \(E(1)\) for \(E\) placed in weight \(1\), and \(E(0)\) for \(E\)
placed in weight \(0\).  Put
\[
    C
    :=
    A[t]\otimes_A\operatorname{Sym}_A(E(1)),
    \qquad
    W
    :=
    \operatorname{Sym}_A(E(0)).
\]
Let
\[
    \alpha:W\longrightarrow C
\]
be the graded algebra morphism whose restriction to \(E\) is
\[
    e\longmapsto t v(e)-s(e),
\]
and let
\[
    \epsilon_0:W\longrightarrow A
\]
be the zero augmentation.  The local equation
\[
    t v_{\mathrm{univ}}-s=0
\]
of Theorem~\ref{thm:spectral-deformation-normal-bundle} gives a canonical
equivalence of graded \(A[t]\)-algebras
\[
    R^{\mathrm{ext}}_{B/A}
    \simeq
    C\otimes_{W,\alpha}A_0.
\]
The relative tensor product is computed by the two-sided bar construction,
and \(W\) is concentrated in weight zero.  It may therefore be computed
weight by weight.

For every \(r\geq0\), the weight-\((-r)\) component of \(C\) is
\[
    C_{-r}
    \simeq
    \bigoplus_{q\geq0}
    t^{q+r}\operatorname{Sym}_A^q(E).
\]
There is a canonical equivalence of \(A\)-modules
\[
    \chi_r:
    t^rW
    \xrightarrow{\simeq}
    C_{-r}
\]
which on the degree-\(q\) symmetric-power summand sends
\[
    t^r x
    \longmapsto
    t^{q+r}v^{(q)}(x).
\]
Under \(\chi_0\), the morphism \(\alpha:W\to C_0\) is the translation
automorphism
\[
    \tau_{-s}:W\xrightarrow{\simeq}W,
    \qquad
    e\longmapsto e-s(e),
\]
whose inverse is \(e\mapsto e+s(e)\).  Consequently, as a right
\(W\)-module through \(\alpha\),
\[
    C_{-r}
    \simeq
    t^rW_{\tau_{-s}}.
\]
Taking the weight-\((-r)\) component of the two-sided bar construction now
gives
\[
\begin{aligned}
    \bigl(R^{\mathrm{ext}}_{B/A}\bigr)_{-r}
    &\simeq
    C_{-r}\otimes_WA_0 \\
    &\simeq
    t^rW_{\tau_{-s}}\otimes_WA_0 \\
    &\simeq
    A\,t^r.
\end{aligned}
\]
These equivalences are multiplicative in \(r\), and the resulting map is
the structural morphism from \(A[t]\).  Hence
\[
    \bigl(R^{\mathrm{ext}}_{B/A}\bigr)_{\leq0}
    \simeq
    A[t].
\]

Base change along the zero evaluation
\[
    A[t]\longrightarrow A,
    \qquad
    t\longmapsto0,
\]
gives the special fibre.  From the displayed pushout presentation,
\[
    R^{\mathrm{ext}}_{B/A}\otimes_{A[t]}A
    \simeq
    \operatorname{Sym}_B(E\otimes_AB),
\]
with \(E\otimes_AB\) in weight \(1\).  In particular, the weight-zero
component of the special fibre is \(B\).  The special fibre is the graded cofiber of multiplication by \(t\).
Taking its weight-zero component therefore gives a cofiber sequence
\[
    \bigl(R^{\mathrm{ext}}_{B/A}\bigr)_1
    \xrightarrow{t}
    \bigl(R^{\mathrm{ext}}_{B/A}\bigr)_0
    \longrightarrow
    B.
\]
The nonpositive calculation identifies the middle term with \(A\), so
\[
    \bigl(R^{\mathrm{ext}}_{B/A}\bigr)_1
    \simeq
    \operatorname{fib}(A\longrightarrow B)
    =
    I_{B/A}.
\]
Since the degree-one component lies in the nonnegative part, this is the
canonical equivalence
\[
    (R_{B/A})_1
    \simeq
    I_{B/A}.
\]

It remains to prove generation on \(\pi_0\).  Trivialize \(E\) with basis
\(e_1,\ldots,e_m\), and write
\[
    x_a:=s(e_a)\in\pi_0(A).
\]
The functor \(\pi_0\) preserves pushouts of connective commutative ring
spectra.  Applied in the graded category to the preceding pushout, it gives
an isomorphism of ordinary graded rings
\[
    \pi_0R^{\mathrm{ext}}_{B/A}
    \cong
    \pi_0(A)[t,v_1,\ldots,v_m]
    /
    (tv_1-x_1,\ldots,tv_m-x_m),
\]
where \(|t|=-1\) and \(|v_a|=1\).  Let
\[
    t^a v_1^{b_1}\cdots v_m^{b_m}
\]
be a homogeneous monomial of nonnegative weight
\[
    n=b_1+\cdots+b_m-a\geq0.
\]
Pair each of the \(a\) copies of \(t\) with one of the positive-weight
factors and use the relations \(tv_j=x_j\).  The monomial becomes an
element of \(\pi_0(A)\) times a monomial of total \(v\)-degree \(n\).
Thus every nonnegative homogeneous component is generated by products of
weight-one elements.  Hence \(\pi_0(R_{B/A})\) is generated over
\(\pi_0(A)\) by its degree-one component.

The nonpositive equivalence is induced by the structural morphism
\(A[t]\to R^{\mathrm{ext}}_{B/A}\), and the degree-one equivalence is the
fibre sequence obtained from multiplication by \(t\) and the special
fibre.  They are therefore intrinsic and natural under commutative
triangles.  The local calculations agree on overlaps, and the generation
statement is Zariski-local.  This gives the asserted global identifications
and their naturality for commutative triangles which are surjective on
\(\pi_0\).
\end{proof}

\begin{lemma}[Proper covariance for the excess deformation]
\label{lem:excess-deformation-proper-covariance}
Let
\[
\begin{tikzcd}[column sep=large,row sep=large]
    Z' \arrow[r,"g'"] \arrow[d,"i'"']
        & Z \arrow[d,"i"] \\
    X' \arrow[r,"g"']
        & X
\end{tikzcd}
\]
be an excess square in the sense of
Definition~\ref{def:closed-excess-square}.  Put
\[
    \widetilde D
    :=
    D_{Z/X}\times_XX'
\]
and let
\[
    \rho:
    D_{Z'/X'}
    \longrightarrow
    \widetilde D
\]
be the canonical morphism induced by the functor-of-points definition of
the deformation prestacks.

The morphism
\[
    \rho:
    D_{Z'/X'}
    \longrightarrow
    \widetilde D
\]
is a spectral closed immersion.  In particular, after classical
truncation, the morphism
\[
    \rho_{\mathrm{cl}}:
    (D_{Z'/X'})_{\mathrm{cl}}
    \longrightarrow
    \widetilde D_{\mathrm{cl}}
\]
is proper.  Its restriction to the generic fibre is the identity
\[
    X'_{\mathrm{cl}}\times\mathbb G_m
    \longrightarrow
    X'_{\mathrm{cl}}\times\mathbb G_m,
\]
and its restriction to the special fibre is the proper morphism
\[
    \rho_0:
    \mathbb V_{Z'_{\mathrm{cl}}}
    \bigl(
      \iota_{Z'}^*N_{i'}^*
    \bigr)
    \longrightarrow
    \mathbb V_{Z'_{\mathrm{cl}}}
    \bigl(
      \iota_{Z'}^*(g')^*N_i^*
    \bigr)
\]
induced by the conormal comparison
\[
    (g')^*N_i^*
    \longrightarrow
    N_{i'}^*.
\]

Write
\[
    D'
    :=
    (D_{Z'/X'})_{\mathrm{cl}},
    \qquad
    \widetilde D'
    :=
    \widetilde D_{\mathrm{cl}},
\]
and denote their special fibres by
\[
    N'
    :=
    \mathbb V_{Z'_{\mathrm{cl}}}
    \bigl(
      \iota_{Z'}^*N_{i'}^*
    \bigr),
    \qquad
    \widetilde N
    :=
    \mathbb V_{Z'_{\mathrm{cl}}}
    \bigl(
      \iota_{Z'}^*(g')^*N_i^*
    \bigr).
\]
Proper covariance for \(\rho_{\mathrm{cl}}\) and \(\rho_0\) gives a
commutative square of localization boundaries
\[
\begin{tikzcd}[column sep=huge,row sep=large]
    \Omega
    \mathbf B_{\mathrm{cl}}
    (D'_\eta/X'_{\mathrm{cl}},0)
        \arrow[r,"\partial_{D'}"]
        \arrow[d,"\simeq"']
        &
    \mathbf B_{\mathrm{cl}}
    (N'/X'_{\mathrm{cl}},0)
        \arrow[d,"\rho_{0*}"]
    \\
    \Omega
    \mathbf B_{\mathrm{cl}}
    (\widetilde D'_\eta/X'_{\mathrm{cl}},0)
        \arrow[r,"\partial_{\widetilde D'}"']
        &
    \mathbf B_{\mathrm{cl}}
    (\widetilde N/X'_{\mathrm{cl}},0).
\end{tikzcd}
\]
The left vertical equivalence is induced by the common generic-fibre
identification and preserves the coordinate section \(\gamma_t\).

Let
\[
    \xi
    \longrightarrow
    (g')^*N_i^*
    \longrightarrow
    N_{i'}^*
\]
be the excess fibre sequence.  Under the inverse Thom equivalences for the
two special-fibre vector-bundle projections, the proper-covariance map
\(\rho_{0*}\) is the map
\[
\begin{aligned}
    \mathbf B_{\mathrm{cl}}
    \bigl(
      Z'_{\mathrm{cl}}/X'_{\mathrm{cl}},
      -\iota_{Z'}^*\nu_{i'}
    \bigr)
    &\longrightarrow
    \mathbf B_{\mathrm{cl}}
    \bigl(
      Z'_{\mathrm{cl}}/X'_{\mathrm{cl}},
      -\iota_{Z'}^*(g')^*\nu_i
    \bigr)
\end{aligned}
\]
obtained by precomposition with
\[
\begin{aligned}
    &
    \operatorname{th}_{Z'_{\mathrm{cl}}}
    \bigl(
      -\iota_{Z'}^*(g')^*\nu_i
    \bigr)
    \\
    &\qquad\xrightarrow{\simeq}
    \operatorname{th}_{Z'_{\mathrm{cl}}}
    \bigl(
      -\iota_{Z'}^*\xi
    \bigr)
    \otimes
    \operatorname{th}_{Z'_{\mathrm{cl}}}
    \bigl(
      -\iota_{Z'}^*\nu_{i'}
    \bigr)
    \\
    &\qquad\xrightarrow{
      e^\vee(\iota_{Z'}^*\xi)\otimes\id
    }
    \operatorname{th}_{Z'_{\mathrm{cl}}}
    \bigl(
      -\iota_{Z'}^*\nu_{i'}
    \bigr).
\end{aligned}
\]
\end{lemma}

\begin{proof}
The functor-of-points construction of
Construction~\ref{constr:spectral-deformation-prestack}
gives the canonical morphism
\[
    \rho:
    D_{Z'/X'}
    \longrightarrow
    D_{Z/X}\times_XX'.
\]
Its generic fibre is the identity because both deformation spaces restrict
canonically to
\[
    X'\times\mathbb G_m^{\mathrm{fl}}.
\]
Its special fibre is induced, by contravariance of relative linear spaces,
from the conormal comparison
\[
    (g')^*N_i^*
    \longrightarrow
    N_{i'}^*.
\]

We prove that \(\rho\) is a spectral closed immersion.  Put
\[
    W:=Z\times_XX'
\]
and let
\[
    e:Z'\longrightarrow W
\]
be the induced morphism.  The assumption that the square is Cartesian on
classical truncations identifies \(e_{\mathrm{cl}}\) with an isomorphism.
Since \(Z'\) and \(W\) are closed over \(X'\), the morphism \(e\) is
affine.  The classical-truncation identification says that its map on
\(\pi_0\) is an isomorphism.  Thus \(e\) is a spectral closed immersion.

Cotangent transitivity for
\[
    Z'\xrightarrow{e}W\longrightarrow X'
\]
gives a fibre sequence
\[
    e^*N^*_{W/X'}
    \longrightarrow
    N^*_{Z'/X'}
    \longrightarrow
    N^*_{Z'/W}.
\]
The first arrow is the conormal comparison in the definition of the excess
square.  Since its fibre is the finite locally free module \(\xi\), the
third term is canonically
\[
    N^*_{Z'/W}
    \simeq
    \xi[1].
\]
It is therefore \(1\)-connective.

Let
\[
    I_e
    :=
    \operatorname{fib}
    \bigl(
      \mathcal O_W
      \longrightarrow
      e_*\mathcal O_{Z'}
    \bigr).
\]
The spectral Hurewicz comparison for closed immersions identifies
\(\pi_0(e^*I_e)\) with \(\pi_0(N^*_{Z'/W})\), and the latter vanishes.
Because \(e\) is a closed immersion, \(I_e\) is connective.  Since
\(\pi_0(\mathcal O_W)\to\pi_0(e_*\mathcal O_{Z'})\) is an isomorphism, the
connective tensor-product formula gives
\[
    \pi_0(e^*I_e)
    \simeq
    \pi_0(I_e).
\]
Consequently,
\[
    \pi_0(I_e)=0,
\]
so \(I_e\) is \(1\)-connective.

The assertion is affine-local on \(X'\).  Write
\[
    X'=\operatorname{Spec}(A),
    \qquad
    W=\operatorname{Spec}(C),
    \qquad
    Z'=\operatorname{Spec}(B),
\]
and put
\[
    I_C:=\operatorname{fib}(A\longrightarrow C),
    \qquad
    I_B:=\operatorname{fib}(A\longrightarrow B).
\]
The stable \(3\times3\) lemma applied to
\[
\begin{tikzcd}[column sep=large,row sep=large]
    I_C \arrow[r] \arrow[d]
        & A \arrow[r] \arrow[d,"\id"]
        & C \arrow[d] \\
    I_B \arrow[r]
        & A \arrow[r]
        & B
\end{tikzcd}
\]
gives
\[
    \operatorname{fib}(I_C\longrightarrow I_B)
    \simeq
    \operatorname{fib}(C\longrightarrow B)[-1]
    \simeq
    I_e[-1].
\]
Since \(I_e\) is \(1\)-connective, the fibre on the left is connective.
Consequently,
\[
    \pi_0(I_C)
    \longrightarrow
    \pi_0(I_B)
\]
is surjective.

Base-change compatibility of deformation identifies
\[
    \widetilde D
    \simeq
    D_{W/X'}.
\]
By
Lemma~\ref{lem:affine-extended-rees-deformation}, the affine coordinate map
of \(\rho\) is the graded morphism
\[
    R^{\mathrm{ext}}_{C/A}
    \longrightarrow
    R^{\mathrm{ext}}_{B/A}.
\]
It is the identity on the nonpositive part \(A[t]\).  On degree one it
is the map
\[
    I_C\longrightarrow I_B,
\]
which is surjective on \(\pi_0\).  Since
\(\pi_0(R_{B/A})\) is generated in degree one, the morphism
\[
    \pi_0
    \bigl(
      R^{\mathrm{ext}}_{C/A}
    \bigr)
    \longrightarrow
    \pi_0
    \bigl(
      R^{\mathrm{ext}}_{B/A}
    \bigr)
\]
is surjective in every degree.  Hence \(\rho\) is a spectral closed
immersion.  In particular, \(\rho_{\mathrm{cl}}\) is a closed immersion and
therefore proper.

On a Zariski-local splitting
\[
    (g')^*N_i^*
    \simeq
    \xi\oplus N_{i'}^*,
\]
the special-fibre map is
\[
    \mathbb V(N_{i'}^*)
    \longrightarrow
    \mathbb V(\xi)
    \times_{Z'}
    \mathbb V(N_{i'}^*),
\]
given by the zero section in the \(\xi\)-direction and the identity in the
\(N_{i'}^*\)-direction.  This description is independent of the chosen
local splitting because it is induced by the canonical conormal morphism.

The morphism \(\rho_{\mathrm{cl}}\) preserves the complementary open and
closed fibres.  Proper covariance for the vertical morphisms and
functoriality of stable localization therefore give a morphism between the
two localization fibre sequences.  This is the displayed commutative
square of boundary maps.  Since the generic-fibre morphism is the identity
over \(X'_{\mathrm{cl}}\times\mathbb G_m\), it preserves the coordinate
\(t\) and hence the coordinate section \(\gamma_t\).

It remains to identify the special-fibre proper-covariance map.  Under the
local splitting above, \(\rho_0\) is the product of the zero section of
\(\mathbb V(\xi)\) with the identity of
\(\mathbb V(N_{i'}^*)\).  The zero-section normalization of
Lemma~\ref{lem:specialization-zero-section-normalization}
identifies proper covariance in the first factor with precomposition by
\[
    e^\vee(\xi):
    \operatorname{th}(-\xi)
    \longrightarrow
    \mathbf1.
\]
Compatibility of inverse Thom equivalences with exterior products
therefore identifies \(\rho_{0*}\) with precomposition by
\[
    e^\vee(\xi)
    \otimes
    \id_{\operatorname{th}(-\nu_{i'})}.
\]

The conormal fibre sequence, its Thom-additivity equivalence, the morphism
\(\rho_0\), and the Euler class are canonical.  The local identifications
therefore agree on overlaps and give the asserted global identification.
\end{proof}

\begin{proposition}[Base change and excess]
\label{prop:virtual-class-base-change-excess}
The virtual fundamental classes satisfy the following formulas.

\begin{enumerate}
    \item For every derived Cartesian square
    \[
    \begin{tikzcd}[column sep=large,row sep=large]
        Z' \arrow[r,"g'"] \arrow[d,"i'"']
            & Z \arrow[d,"i"] \\
        X' \arrow[r,"g"']
            & X,
    \end{tikzcd}
    \]
    cotangent base change gives
    \[
        \nu_{i'}
        \simeq
        (g')^*\nu_i,
    \]
    and
    \[
        \operatorname{Ex}^{!*}_{i,g}(\mathbf1_X)
        \circ
        (g')^*\eta_i
        =
        \eta_{i'}.
    \]

    \item For every excess square with excess bundle \(\xi\), let
    \[
        e(\xi):
        \mathbf1_{Z'}
        \longrightarrow
        \operatorname{th}_{Z'}(\xi)
    \]
    be the Euler class of the zero section, and let
    \[
        e^\vee(\xi):
        \operatorname{th}_{Z'}(-\xi)
        \longrightarrow
        \mathbf1_{Z'}
    \]
    be the morphism obtained by tensoring \(e(\xi)\) with
    \(\operatorname{th}_{Z'}(-\xi)\) and applying Thom cancellation.

    Under the Thom-additivity equivalence
    \[
        \operatorname{th}_{Z'}
        \bigl(
          -(g')^*\nu_i
        \bigr)
        \simeq
        \operatorname{th}_{Z'}(-\xi)
        \otimes
        \operatorname{th}_{Z'}(-\nu_{i'}),
    \]
    one has
    \[
        \Delta^\star\eta_i
        =
        \eta_{i'}
        \circ
        \bigl(
          e^\vee(\xi)
          \otimes
          \id_{\operatorname{th}(-\nu_{i'})}
        \bigr).
    \]
\end{enumerate}

Both formulas are compatible with composition of comparison squares.
\end{proposition}

\begin{proof}
For every commutative square
\[
\begin{tikzcd}[column sep=large,row sep=large]
    Z' \arrow[r,"g'"] \arrow[d,"i'"']
        & Z \arrow[d,"i"] \\
    X' \arrow[r,"g"']
        & X,
\end{tikzcd}
\]
the functor-of-points definition of the deformation prestack gives a
canonical morphism
\[
    \rho:
    D_{Z'/X'}
    \longrightarrow
    D_{Z/X}\times_XX'.
\]
Indeed, for every test spectral scheme
\[
    T\longrightarrow X'\times\mathbb A^{1,\mathrm{fl}},
\]
a map
\[
    T_0\longrightarrow Z'
\]
over \(X'\) composes with \(g':Z'\to Z\) to give a map
\[
    T_0\longrightarrow Z
\]
over \(X\).  This construction is natural in \(T\).

Suppose first that the square is derived Cartesian.  Then
\[
    Z'
    \simeq
    Z\times_XX',
\]
so the preceding map of functors of points is an equivalence:
\[
    D_{Z'/X'}
    \xrightarrow{\simeq}
    D_{Z/X}\times_XX'.
\]
It preserves the canonical centres, generic fibres, special fibres, and
conormal identifications.

After classical truncation, exceptional base change identifies the
localization boundary for \(D_{Z'/X'}\) with the pullback of the
localization boundary for \(D_{Z/X}\).  The coordinate section
\(\gamma_t\) and the vector-bundle inverse Thom equivalence also commute
with pullback.  Therefore
\[
    \operatorname{Ex}^{!*}_{i_{\mathrm{cl}},g_{\mathrm{cl}}}
    (\mathbf1_{X_{\mathrm{cl}}})
    \circ
    (g'_{\mathrm{cl}})^*
    \eta^{\mathrm{vir}}_{i,\mathrm{cl}}
    =
    \eta^{\mathrm{vir}}_{i',\mathrm{cl}}.
\]
Stable comparison transports this equality to
\[
    \operatorname{Ex}^{!*}_{i,g}(\mathbf1_X)
    \circ
    (g')^*\eta_i
    =
    \eta_{i'}.
\]
Cotangent base change gives
\[
    \nu_{i'}
    \simeq
    (g')^*\nu_i,
\]
so this proves part~(1).

Now suppose that the square is an excess square.  Put
\[
    \widetilde D
    :=
    D_{Z/X}\times_XX'
\]
and write
\[
    \widetilde N
    :=
    \mathbb V_{Z'}
    \bigl(
      (g')^*N_i^*
    \bigr)
\]
for its special fibre.  Pullback of the original deformation datum to
\(X'\) gives, after classical truncation, a specialization class
\[
    \sigma_{\widetilde D}
    \in
    \mathbb B_{\mathrm{cl}}
    \bigl(
      \widetilde N_{\mathrm{cl}}/X'_{\mathrm{cl}},
      0
    \bigr).
\]
Under the inverse Thom equivalence for
\[
    \widetilde N_{\mathrm{cl}}
    \longrightarrow
    Z'_{\mathrm{cl}},
\]
this class is precisely the refined pullback
\[
    \Delta^\star
    \eta^{\mathrm{vir}}_{i,\mathrm{cl}}
    \in
    \mathbb B_{\mathrm{cl}}
    \bigl(
      Z'_{\mathrm{cl}}/X'_{\mathrm{cl}},
      -\iota_{Z'}^*(g')^*\nu_i
    \bigr).
\]

The canonical comparison
\[
    \rho:
    D_{Z'/X'}
    \longrightarrow
    \widetilde D
\]
satisfies
Lemma~\ref{lem:excess-deformation-proper-covariance}.
Its generic-fibre morphism is the identity and preserves the coordinate
section \(\gamma_t\).  Naturality of the localization boundary under proper
covariance therefore gives
\[
    \sigma_{\widetilde D}
    =
    \rho_{0*}
    \bigl(
      \sigma_{D_{Z'/X'}}
    \bigr)
\]
in
\[
    \mathbb B_{\mathrm{cl}}
    \bigl(
      \widetilde N_{\mathrm{cl}}/X'_{\mathrm{cl}},
      0
    \bigr).
\]

Under the inverse Thom equivalence for
\[
    \mathbb V_{Z'_{\mathrm{cl}}}
    \bigl(
      \iota_{Z'}^*N_{i'}^*
    \bigr)
    \longrightarrow
    Z'_{\mathrm{cl}},
\]
the class
\[
    \sigma_{D_{Z'/X'}}
\]
corresponds to
\[
    \eta^{\mathrm{vir}}_{i',\mathrm{cl}}.
\]
By the final assertion of
Lemma~\ref{lem:excess-deformation-proper-covariance},
proper covariance along \(\rho_0\) corresponds, under the two inverse Thom
equivalences, to precomposition with
\[
\begin{aligned}
    &
    \operatorname{th}_{Z'_{\mathrm{cl}}}
    \bigl(
      -\iota_{Z'}^*(g')^*\nu_i
    \bigr)
    \\
    &\qquad\xrightarrow{\simeq}
    \operatorname{th}_{Z'_{\mathrm{cl}}}
    \bigl(
      -\iota_{Z'}^*\xi
    \bigr)
    \otimes
    \operatorname{th}_{Z'_{\mathrm{cl}}}
    \bigl(
      -\iota_{Z'}^*\nu_{i'}
    \bigr)
    \\
    &\qquad\xrightarrow{
      e^\vee(\iota_{Z'}^*\xi)\otimes\id
    }
    \operatorname{th}_{Z'_{\mathrm{cl}}}
    \bigl(
      -\iota_{Z'}^*\nu_{i'}
    \bigr).
\end{aligned}
\]
Consequently,
\[
    \Delta^\star
    \eta^{\mathrm{vir}}_{i,\mathrm{cl}}
    =
    \eta^{\mathrm{vir}}_{i',\mathrm{cl}}
    \circ
    \bigl(
      e^\vee(\iota_{Z'}^*\xi)
      \otimes
      \id_{\operatorname{th}(-\iota_{Z'}^*\nu_{i'})}
    \bigr).
\]

Stable comparison identifies the Thom-additivity path
\[
    (g')^*\nu_i
    \simeq
    \xi+\nu_{i'},
\]
the Euler class, the refined pullback transformation, and the two virtual
fundamental classes with their classical counterparts.  Transporting the
preceding equality gives
\[
    \Delta^\star\eta_i
    =
    \eta_{i'}
    \circ
    \bigl(
      e^\vee(\xi)
      \otimes
      \id_{\operatorname{th}(-\nu_{i'})}
    \bigr).
\]
This proves part~(2).

For composable comparison squares, the base-changed deformation data and
the canonical excess-deformation morphisms compose.  Proper covariance and
localization boundaries are functorial under composition.  The successive
excess bundles fit into the corresponding fibre sequences, Euler classes
are multiplicative under direct sums and extensions, and Thom additivity is
coherently associative.  Exceptional exchange and stable comparison
preserve these composition diagrams.  Hence both formulas are compatible
with composition of comparison squares.
\end{proof}

\begin{proposition}[Formal properties of closed purity]
\label{prop:closed-purity-formal-properties}
The transformations \(\mathfrak q_i\) have the following properties.
\begin{enumerate}
    \item They are natural in the coefficient object.

    \item They satisfy derived base change and the type-correct excess
    formula of
    Proposition~\ref{prop:virtual-class-base-change-excess}.

    \item For composable quasi-smooth closed immersions, the purity
    transformation for the composite is canonically homotopic to the pasted
    composite of the two purity transformations under the
    conormal-additivity path.  For a finite chain, every parenthesized pasted
    composite represents the same component of the corresponding
    natural-transformation space.

    \item If \(L\in\Pic(\SH(X))\), then purity for \(A\) implies purity
    for \(A\otimes L\), and the resulting equivalence is the tensor product
    of \(\mathfrak q_i(A)\) with \(i^*L\).

    \item For the zero section
    \[
        s_E:X\hookrightarrow\mathbb V_X(E),
    \]
    the transformation is the inverse Thom transformation.

    \item On objects pulled back from a common smooth base, the
    transformation agrees with the relative closed-purity transformation of
    Construction~\ref{constr:closed-purity-transformation}.
\end{enumerate}
\end{proposition}

\begin{proof}
Naturality in coefficients follows from the module transformation used in
Construction~\ref{constr:absolute-closed-purity-transformation}.  Derived
base change and excess follow by tensoring the formulas of
Proposition~\ref{prop:virtual-class-base-change-excess} with the pulled-back
coefficient and using naturality of the module transformation.

The composition assertion and the finite-chain statement on components
follow from
Theorem~\ref{thm:virtual-class-composition-coherence} and coherent
associativity of the bivariant product.  If \(L\) is invertible, the
dualizable coefficient formula gives
\[
    i^!(A\otimes_XL)
    \simeq
    i^!A\otimes_Zi^*L,
\]
and under this equivalence
\[
    \mathfrak q_i(A\otimes L)
    \simeq
    \mathfrak q_i(A)\otimes i^*L.
\]

For a zero section, the deformation datum is locally given by the equation
\[
    x=tv.
\]
Theorem~\ref{thm:specialization-deformation-datum} and
Lemma~\ref{lem:specialization-zero-section-normalization}
identify its specialization class with the inverse Thom class.  The
coefficientwise transformation obtained from this class is therefore the
inverse Thom transformation.

Suppose now that the closed pair is smooth over a common base.  Under
stable classical comparison, the relative closed-purity transformation of
Construction~\ref{constr:closed-purity-transformation} is carried to the
ordinary relative-purity transformation for the corresponding ordinary
smooth closed pair.

On the other hand, the transformation
\[
    \mathfrak q_i:
    \Sigma^{-\nu_i}i^*
    \longrightarrow
    i^!
\]
constructed in
Construction~\ref{constr:absolute-closed-purity-transformation}
is carried to the purity transformation associated with the ordinary
deformation fundamental class
\[
    \eta^{\mathrm{vir}}_{i,\mathrm{cl}}.
\]
In ordinary motivic homotopy theory, the deformation fundamental class of a
regular closed immersion is compatible with the relative-purity
transformation for a smooth closed pair.  Consequently, the ordinary images
of the two transformations agree.

Stable classical comparison is an equivalence and therefore induces an
equivalence on spaces of natural transformations.  It follows that the two
spectral transformations agree.
\end{proof}

\begin{corollary}[Closed purity on pure coefficients]
\label{cor:closed-purity-pure-coefficients}
If \(A\in\SH(X)\) is pure for \(i\), then there is a canonical equivalence
\[
    i^!A
    \simeq
    \Sigma^{-\nu_i}i^*A.
\]
The equivalence is compatible with base change, composition, and Thom
twists whenever all coefficients involved are pure for the corresponding
immersions.
\end{corollary}

\begin{proof}
This is Definition~\ref{def:i-pure-object} together with
Proposition~\ref{prop:closed-purity-formal-properties}.
\end{proof}

\begin{remark}[No unconditional integral absolute-purity theorem]
\label{rem:no-integral-absolute-purity}
The chapter does not assert that every object of \(\SH(X)\), or even the
sphere object, is pure for every regular closed immersion.  Such a statement
would be an absolute-purity theorem beyond the formal six-operation
construction.  Whenever a coefficient theory is known independently to
make the virtual specialization class of
Construction~\ref{constr:virtual-closed-fundamental-class} invertible, the
corresponding spectral absolute-purity equivalence follows from stable
comparison and the preceding construction.
\end{remark}


\section{Localization and recollement in six-functor notation}
\label{sec:six-localization}

Let
\[
    Z\xrightarrow{i}X\xleftarrow{j}U
\]
be a complementary closed--open pair, where \(j\) is quasi-compact.

\begin{theorem}[Six-functor localization]
\label{thm:six-functor-localization}
For every \(A\in\SH(X)\), there are canonical exact triangles
\[
    j_!j^*A
    \longrightarrow
    A
    \longrightarrow
    i_*i^*A
\]
and
\[
    i_*i^!A
    \longrightarrow
    A
    \longrightarrow
    j_*j^*A.
\]
The functors \(j_!,j_*,i_*\) are fully faithful, and
\[
    i^*j_!\simeq0,
    \qquad
    i^!j_*\simeq0,
    \qquad
    j^*i_*\simeq0.
\]

The first localization triangle is compatible with arbitrary pullback and
arbitrary tensor products.

For the co-localization triangle, every morphism
\[
    g:X'\longrightarrow X
\]
gives a canonical comparison from the pullback of the displayed triangle
to the co-localization triangle for the pulled-back complementary pair.
This comparison is an equivalence when \(g\) is smooth.

Suppose additionally that \(i\) and its pullback
\[
    i':Z'\hookrightarrow X'
\]
are quasi-smooth with finite locally free conormal modules.  If \(A\) is
pure for \(i\), \(g^*A\) is pure for \(i'\), and the two purity
transformations are compatible with the base-change square, then the
co-localization comparison is also an equivalence on \(A\).

For \(B\in\SH(X)\), there is a canonical transformation
\[
    i^!A\otimes_Zi^*B
    \longrightarrow
    i^!(A\otimes_XB).
\]
It is an equivalence when \(B\) is strongly dualizable.  Internal Homs
satisfy the unconditional support identity
\[
    \underline{\operatorname{Hom}}_X
    \bigl(
      B,
      i_*i^!A
    \bigr)
    \simeq
    i_*i^!
    \underline{\operatorname{Hom}}_X(B,A).
\]
\end{theorem}

\begin{proof}
Under the open normalization
\[
    j_!\simeq j_\sharp,
\]
the first triangle is stable localization from
Theorem~\ref{thm:stable-localization}.  The second is stable
co-localization from
Theorem~\ref{thm:stable-colocalization}.  Full faithfulness and
orthogonality are the stable recollement package of
Proposition~\ref{prop:stable-recollement}.

Arbitrary pullback compatibility of the first triangle follows from open
and closed base change.  Its tensor compatibility follows from the open and
closed projection formulas.

For the second triangle, closed base change gives
\[
    g^*i_*
    \simeq
    i'_*g_Z^*,
\]
and the canonical closed exchange map
\[
    g_Z^*i^!
    \longrightarrow
    i'^!g^*
\]
gives the comparison of co-localization triangles.  Smooth--closed exchange
makes this map an equivalence when \(g\) is smooth.

Now suppose that \(i\) and \(i'\) are quasi-smooth with finite locally free
conormal modules and that the stated purity conditions hold.  Then
\[
    i^!A
    \simeq
    \Sigma^{-\nu_i}i^*A
\]
and
\[
    i'^!g^*A
    \simeq
    \Sigma^{-\nu_{i'}}i'^*g^*A.
\]
Compatibility of the purity transformations with the square, together with
base change of the conormal class, identifies
\[
    g_Z^*i^!A
    \longrightarrow
    i'^!g^*A
\]
with an equivalence.  Hence the comparison of co-localization triangles is
an equivalence on \(A\).

The tensor transformation is adjoint to
\[
\begin{aligned}
    i_*
    \bigl(
      i^!A\otimes_Zi^*B
    \bigr)
    &\simeq
    i_*i^!A\otimes_XB \\
    &\longrightarrow
    A\otimes_XB.
\end{aligned}
\]
When \(B\) is strongly dualizable,
Proposition~\ref{prop:closed-internal-hom} gives
\[
\begin{aligned}
    i^!(A\otimes_XB)
    &\simeq
    i^!
    \underline{\operatorname{Hom}}_X(B^\vee,A) \\
    &\simeq
    \underline{\operatorname{Hom}}_Z
    \bigl(
      i^*B^\vee,
      i^!A
    \bigr) \\
    &\simeq
    i^!A\otimes_Zi^*B.
\end{aligned}
\]
Under this equivalence the displayed tensor transformation is the identity.

Finally,
Proposition~\ref{prop:closed-internal-hom}
and full faithfulness of \(i_*\) give
\[
\begin{aligned}
    \underline{\operatorname{Hom}}_X
    \bigl(
      B,
      i_*i^!A
    \bigr)
    &\simeq
    i_*
    \underline{\operatorname{Hom}}_Z
    \bigl(
      i^*B,
      i^!A
    \bigr) \\
    &\simeq
    i_*i^!
    \underline{\operatorname{Hom}}_X(B,A).
\end{aligned}
\]
\end{proof}

\begin{proposition}[Exceptional direct image and localization]
\label{prop:exceptional-direct-image-localization}
Let \(f:X\to Y\) belong to \(\mathcal E\), and let
\[
    Z\xrightarrow{i}X\xleftarrow{j}U
\]
be a complementary closed--open pair with \(j\) quasi-compact.  Put
\[
    f_U:=fj.
\]
For every \(A\in\SH(X)\), applying \(f_!\) to localization gives a
canonical exact triangle
\[
    (f_U)_!j^*A
    \longrightarrow
    f_!A
    \longrightarrow
    f_!i_*i^*A.
\]
It is compatible with base change and composition in \(f\).

When \(i\) is of finite presentation, put \(f_Z:=fi\).  Then
\(f_Z\in\mathcal E\), and exceptional composition and closed normalization
identify the third term with
\[
    f_!i_*i^*A
    \simeq
    (f_Z)_!i^*A.
\]
\end{proposition}

\begin{proof}
The functor \(f_!\) preserves colimits and is exact.  Apply it to the first
localization triangle and use exceptional composition for the open term:
\[
    f_!j_!
    \simeq
    (fj)_!.
\]
This gives the first displayed triangle without imposing a
finite-presentation condition on the closed immersion.

If \(i\) is of finite presentation, closed normalization gives
\[
    i_*\simeq i_!,
\]
and exceptional composition gives
\[
    f_!i_!
    \simeq
    (fi)_!.
\]
The compatibility assertions follow from the coherent composition and
base-change data.
\end{proof}

\begin{proposition}[Exceptional inverse image and co-localization]
\label{prop:exceptional-inverse-colocalization}
Let \(f:X\to Y\) belong to \(\mathcal E\).  For a complementary pair
\[
    Z\xrightarrow{i}Y\xleftarrow{j}U
\]
with \(j\) quasi-compact, pullback to \(X\) gives a complementary pair
\[
    Z_X\xrightarrow{i'}X\xleftarrow{j'}U_X
\]
and base-changed morphisms
\[
    f_Z:Z_X\longrightarrow Z,
    \qquad
    f_U:U_X\longrightarrow U.
\]
There are canonical equivalences
\[
    (j')^*f^!
    \simeq
    f_U^!j^*,
    \qquad
    (i')^!f^!
    \simeq
    f_Z^!i^!.
\]
Under these identifications, \(f^!\) carries the co-localization triangle
on \(Y\) to the corresponding co-localization triangle on \(X\).
\end{proposition}

\begin{proof}
Open normalization and exceptional inverse-image composition give
\[
\begin{aligned}
    (j')^*f^!
    &\simeq
    (j')^!f^! \\
    &\simeq
    (fj')^! \\
    &\simeq
    (jf_U)^! \\
    &\simeq
    f_U^!j^! \\
    &\simeq
    f_U^!j^*.
\end{aligned}
\]

For the closed term, choose a compactification
\[
    X\xrightarrow{k}\overline X\xrightarrow{p}Y
\]
of \(f\), and pull it back over \(Z\):
\[
\begin{tikzcd}[column sep=large,row sep=large]
    Z_X \arrow[r,"k_Z"] \arrow[d,"i'"']
        & \overline X_Z \arrow[r,"p_Z"] \arrow[d,"\overline i"']
        & Z \arrow[d,"i"] \\
    X \arrow[r,"k"']
        & \overline X \arrow[r,"p"']
        & Y.
\end{tikzcd}
\]
Stable smooth--closed exchange for the open immersion \(k\) gives
\[
    k_\sharp i'_*
    \simeq
    \overline i_*k_{Z\sharp}.
\]
Hence
\[
\begin{aligned}
    f_!i'_*
    &\simeq
    p_*k_\sharp i'_* \\
    &\simeq
    p_*\overline i_*k_{Z\sharp} \\
    &\simeq
    i_*p_{Z*}k_{Z\sharp} \\
    &\simeq
    i_*f_{Z!}.
\end{aligned}
\]
The construction is invariant under proper refinements, so it is canonical.
Taking right adjoints yields
\[
    (i')^!f^!
    \simeq
    f_Z^!i^!.
\]

Right-adjoint exceptional base change also gives
\[
    f^!j_*
    \simeq
    j'_*f_U^!,
    \qquad
    f^!i_*
    \simeq
    i'_*f_Z^!.
\]
Applying the exact functor \(f^!\) to the co-localization triangle on \(Y\)
and using the preceding equivalences identifies its three terms with
\[
    i'_*i'^!f^!A,
    \qquad
    f^!A,
    \qquad
    j'_*j'^*f^!A.
\]
This is the co-localization triangle on \(X\).
\end{proof}

\section{Cohomology with support and Borel--Moore objects}
\label{sec:six-borel-moore}

\subsection{Cohomology with support}

\begin{definition}[Cohomology with support]
\label{def:cohomology-with-support}
Let \(i:Z\hookrightarrow X\) be a closed immersion with quasi-compact open
complement.  Define the support functor
\[
    \Gamma_Z:
    \SH(X)
    \longrightarrow
    \SH(X)
\]
by
\[
    \Gamma_Z(A)
    :=
    i_*i^!A.
\]
\end{definition}

\begin{proposition}[Support triangle]
\label{prop:support-triangle}
For every \(A\in\SH(X)\), there is a canonical exact triangle
\[
    \Gamma_Z(A)
    \longrightarrow
    A
    \longrightarrow
    j_*j^*A.
\]
The functor \(\Gamma_Z\) is idempotent, and its essential image is the full
subcategory of objects whose restriction to \(U\) vanishes.
\end{proposition}

\begin{proof}
The triangle is the second localization triangle.  Full faithfulness of
\(i_*\) gives
\[
    \Gamma_Z\Gamma_Z
    \simeq
    i_*i^!i_*i^!
    \simeq
    i_*i^!
    =
    \Gamma_Z.
\]
The support description is stable recollement.
\end{proof}

\begin{proposition}[Base change and nested supports]
\label{prop:support-base-change-nested}
The support functor has the following properties.

\begin{enumerate}
    \item For every morphism
    \[
        g:X'\longrightarrow X,
    \]
    with pulled-back closed immersion
    \[
        i':Z'\hookrightarrow X',
    \]
    there is a canonical transformation
    \[
        g^*\Gamma_Z
        \longrightarrow
        \Gamma_{Z'}g^*.
    \]
    It is an equivalence when \(g\) is smooth.

    Suppose that \(i\) and \(i'\) are quasi-smooth with finite locally free
    conormal modules.  It is also an equivalence on \(A\) whenever
    \(A\) is pure for \(i\), \(g^*A\) is pure for \(i'\), and the purity
    transformations are compatible with the base-change square.

    \item If
    \[
        W\xrightarrow{k}Z\xrightarrow{i}X
    \]
    are closed immersions, then
    \[
        \Gamma_W
        \simeq
        (ik)_*(ik)^!
        \simeq
        i_*\,k_*k^!\,i^!.
    \]

    \item If \(i\) is quasi-smooth with finite locally free conormal module
    and \(A\) is pure for \(i\), then
    \[
        \Gamma_Z(A)
        \simeq
        i_*\Sigma^{-\nu_i}i^*A.
    \]
\end{enumerate}
\end{proposition}

\begin{proof}
Closed base change gives
\[
    g^*i_*
    \simeq
    i'_*g_Z^*.
\]
Composing with the canonical exchange transformation
\[
    g_Z^*i^!
    \longrightarrow
    i'^!g^*
\]
gives
\[
    g^*i_*i^!
    \longrightarrow
    i'_*i'^!g^*,
\]
which is the asserted transformation
\[
    g^*\Gamma_Z
    \longrightarrow
    \Gamma_{Z'}g^*.
\]
Smooth--closed exchange makes it an equivalence when \(g\) is smooth.

Suppose now that \(i\) and \(i'\) are quasi-smooth with finite locally free
conormal modules and that the stated purity conditions hold.  Then
\[
    i^!A
    \simeq
    \Sigma^{-\nu_i}i^*A
\]
and
\[
    i'^!g^*A
    \simeq
    \Sigma^{-\nu_{i'}}i'^*g^*A.
\]
Compatibility of purity with the square and base change of the conormal
class identify
\[
    g_Z^*i^!A
    \longrightarrow
    i'^!g^*A
\]
with an equivalence.  Applying \(i'_*\) proves the first assertion on \(A\).

For nested supports, exceptional inverse-image composition gives
\[
    (ik)^!
    \simeq
    k^!i^!,
\]
while composition of closed direct images gives
\[
    (ik)_*
    \simeq
    i_*k_*.
\]
Hence
\[
    (ik)_*(ik)^!
    \simeq
    i_*k_*k^!i^!.
\]

Finally, if \(A\) is pure for \(i\), then
\[
\begin{aligned}
    \Gamma_Z(A)
    &=
    i_*i^!A \\
    &\simeq
    i_*\Sigma^{-\nu_i}i^*A.
\end{aligned}
\]
\end{proof}

\subsection{Borel--Moore objects}

\begin{definition}[Borel--Moore motive]
\label{def:borel-moore-motive}
Let
\[
    f:X\longrightarrow S
\]
belong to \(\mathcal E\).  Define the Borel--Moore motive of \(X\) over
\(S\) by
\[
    \mathrm M^{\mathrm{BM}}_S(X)
    :=
    f_!\mathbf1_X
    \in
    \SH(S).
\]
More generally, for \(A\in\SH(X)\), define the Borel--Moore object with
coefficients \(A\) by
\[
    \mathrm M^{\mathrm{BM}}_S(X;A)
    :=
    f_!A.
\]
\end{definition}


\begin{proposition}[Formal properties of Borel--Moore objects]
\label{prop:borel-moore-formal-properties}
Borel--Moore objects satisfy the following properties.
\begin{enumerate}
    \item Suppose
    \[
        X\xrightarrow{p}Y\xrightarrow{g}S
    \]
    with \(p\) proper and \(p,g\in\mathcal E\).  For every
    \(A\in\SH(X)\), exceptional composition and proper normalization give
    a canonical equivalence
    \[
        \mathrm M^{\mathrm{BM}}_S(X;A)
        \xrightarrow{\simeq}
        \mathrm M^{\mathrm{BM}}_S(Y;p_*A).
    \]
    More generally, every morphism
    \[
        B\longrightarrow p_*A
    \]
    induces
    \[
        \mathrm M^{\mathrm{BM}}_S(Y;B)
        \longrightarrow
        \mathrm M^{\mathrm{BM}}_S(X;A).
    \]
    In particular, the adjunction unit
    \[
        \mathbf1_Y
        \longrightarrow
        p_*\mathbf1_X
    \]
    gives a canonical map
    \[
        \mathrm M^{\mathrm{BM}}_S(Y)
        \longrightarrow
        \mathrm M^{\mathrm{BM}}_S(X).
    \]
    Consequently, for every \(E\in\SH(S)\), the represented theory
    \[
        E^{\mathrm{BM}}(X/S)
        :=
        \operatorname{Map}_{\SH(S)}
        \bigl(
          \mathrm M^{\mathrm{BM}}_S(X),
          E
        \bigr)
    \]
    has the usual proper pushforward
    \[
        E^{\mathrm{BM}}(X/S)
        \longrightarrow
        E^{\mathrm{BM}}(Y/S).
    \]

    \item For every Cartesian square
    \[
    \begin{tikzcd}[column sep=large,row sep=large]
        X_T \arrow[r] \arrow[d,"f_T"']
            & X \arrow[d,"f"] \\
        T \arrow[r,"g"']
            & S,
    \end{tikzcd}
    \]
    exceptional base change gives
    \[
        g^*\mathrm M^{\mathrm{BM}}_S(X;A)
        \simeq
        \mathrm M^{\mathrm{BM}}_T(X_T;A|_{X_T}).
    \]

    \item For a complementary pair
    \[
        Z\xrightarrow{i}X\xleftarrow{j}U,
    \]
    suppose that the structural morphisms
    \[
        X\longrightarrow S,
        \qquad
        Z\longrightarrow S,
        \qquad
        U\longrightarrow S
    \]
    belong to \(\mathcal E\).  Then there is an exact triangle
    \[
        \mathrm M^{\mathrm{BM}}_S(U;j^*A)
        \longrightarrow
        \mathrm M^{\mathrm{BM}}_S(X;A)
        \longrightarrow
        \mathrm M^{\mathrm{BM}}_S(Z;i^*A).
    \]

    \item For exceptional morphisms \(X\to S\) and \(Y\to T\), there is a
    canonical exterior-product equivalence
    \[
        \mathrm M^{\mathrm{BM}}_S(X;A)
        \boxtimes
        \mathrm M^{\mathrm{BM}}_T(Y;B)
        \simeq
        \mathrm M^{\mathrm{BM}}_{S\times T}
        (X\times Y;A\boxtimes B).
    \]
\end{enumerate}
\end{proposition}

\begin{proof}
For the first assertion,
\[
\begin{aligned}
    \mathrm M^{\mathrm{BM}}_S(X;A)
    &=(gp)_!A \\
    &\simeq
    g_!p_!A \\
    &\simeq
    g_!p_*A \\
    &=
    \mathrm M^{\mathrm{BM}}_S(Y;p_*A).
\end{aligned}
\]
Applying \(g_!\) to a morphism \(B\to p_*A\) gives the stated map.  For
unit coefficients this morphism is the unit of \(p^*\dashv p_*\).  Applying
\(\operatorname{Map}_{\SH(S)}(-,E)\) reverses its direction and gives the
usual covariant proper pushforward on the represented Borel--Moore theory.

The base-change formula is
Theorem~\ref{thm:exceptional-base-change}.  For the localization statement,
the closed immersion \(i:Z\to X\) belongs to \(\mathcal E\) under the
stated scope.  Indeed, it factors over \(S\) as
\[
    Z
    \xrightarrow{\Gamma_i}
    Z\times_SX
    \xrightarrow{\operatorname{pr}_X}
    X.
\]
The projection is a base change of \(Z\to S\), hence is separated and of
finite presentation.  The graph \(\Gamma_i\) is a base change of the
diagonal of \(X\to S\), hence is a closed immersion of finite presentation
because \(X\to S\) is separated and of finite presentation.  Therefore
\(i\in\mathcal E\).  Proposition~\ref{prop:exceptional-direct-image-localization}
and closed normalization then identify the closed term with
\[
    (fi)_!i^*A,
\]
giving the stated Borel--Moore localization triangle.  External products
are Proposition~\ref{prop:exceptional-external-products}.
\end{proof}

\begin{proposition}[Smooth Borel--Moore formula]
\label{prop:smooth-borel-moore-formula}
Let \(p:X\to S\) be smooth, separated, and of finite presentation.  Then
\[
    \mathrm M^{\mathrm{BM}}_S(X;A)
    \simeq
    p_\sharp\Sigma^{-T_p}A.
\]
In particular,
\[
    \mathrm M^{\mathrm{BM}}_S(X)
    \simeq
    p_\sharp\operatorname{th}_X(-T_p).
\]
\end{proposition}

\begin{proof}
This is Corollary~\ref{cor:smooth-exceptional-direct-image} evaluated on
\(A\).
\end{proof}

\subsection{Relative dualizing objects}

\begin{definition}[Relative dualizing object]
\label{def:relative-dualizing-object}
For \(f:X\to Y\) in \(\mathcal E\), define
\[
    \omega_f
    :=
    f^!\mathbf1_Y
    \in
    \SH(X).
\]
\end{definition}

\begin{proposition}[Calculus of dualizing objects]
\label{prop:dualizing-object-calculus}
The relative dualizing objects satisfy:
\begin{enumerate}
    \item for composable \(X\xrightarrow{f}Y\xrightarrow{g}Z\),
    \[
        \omega_{gf}
        \simeq
        f^!\omega_g;
    \]

    \item for smooth \(f\),
    \[
        \omega_f
        \simeq
        \operatorname{th}_X(T_f);
    \]

    \item for a quasi-smooth closed immersion \(i\), if \(\mathbf1_X\) is
    pure for \(i\), then
    \[
        \omega_i
        \simeq
        \operatorname{th}_Z(-\nu_i).
    \]
\end{enumerate}
\end{proposition}

\begin{proof}
The first assertion is exceptional inverse-image composition evaluated on
\(\mathbf1_Z\).  The second is smooth purity and
\(f^*\mathbf1_Y\simeq\mathbf1_X\).  The third is closed purity evaluated
on the unit.
\end{proof}

\section{Constructible and compact objects}
\label{sec:six-constructible}

\begin{definition}[Constructible objects]
\label{def:six-constructible-objects}
Let \(S\) be a quasi-compact quasi-separated spectral scheme.  Define
\[
    \SH_c(S)
    \subseteq
    \SH(S)
\]
to be the smallest thick idempotent-complete subcategory containing every
object of the form
\[
    p_\sharp\Sigma^\xi\mathbf1_X,
\]
where
\[
    p:X\longrightarrow S
\]
is smooth of finite presentation and
\[
    \xi\in K(X).
\]
When \(p\) is separated, smooth purity identifies this generator with
\[
    p_!\Sigma^{T_p+\xi}\mathbf1_X.
\]
\end{definition}

\begin{theorem}[Constructible objects are the compact objects]
\label{thm:constructible-equals-compact}
For every quasi-compact quasi-separated spectral scheme \(S\), there is an
equality of full subcategories
\[
    \SH_c(S)
    =
    \SH(S)^\omega.
\]
\end{theorem}

\begin{proof}
Let \(p:X\to S\) be smooth of finite presentation.  The sphere
\(\mathbf1_X\) is compact by Corollary~\ref{cor:compact-sphere}, and virtual
Thom twists preserve compact objects.  The functor \(p_\sharp\) preserves
compact objects because its right adjoint \(p^*\) preserves filtered
colimits.  Hence every generator of \(\SH_c(S)\) is compact, so
\[
    \SH_c(S)
    \subseteq
    \SH(S)^\omega.
\]

Conversely, Theorem~\ref{thm:compact-generation-SH} gives compact generators
of the form
\[
    \Sigma^{n\mathcal O_S}
    \Sigma_{T,+}^{\infty}\Mot_S(X),
\]
where \(p:X\to S\) is smooth of finite presentation and \(n\in\mathbb Z\).
The stable smooth motive is
\[
    \Sigma_{T,+}^{\infty}\Mot_S(X)
    \simeq
    p_\sharp\mathbf1_X.
\]
Using the smooth projection formula,
\[
    \Sigma^{n\mathcal O_S}p_\sharp\mathbf1_X
    \simeq
    p_\sharp\Sigma^{n\mathcal O_X}\mathbf1_X.
\]
Thus every standard compact generator belongs to \(\SH_c(S)\).  When
\(p\) is separated, the same object may equivalently be written
\[
    p_!\Sigma^{T_p+n\mathcal O_X}\mathbf1_X.
\]  The compact
objects are the thick idempotent-complete subcategory generated by those
objects, proving the reverse inclusion.
\end{proof}

\begin{proposition}[Formal preservation of constructibles]
\label{prop:formal-constructible-preservation}
The following functors preserve constructible objects.
\begin{enumerate}
    \item \(f^*\) for every morphism \(f\) of qcqs spectral schemes;

    \item tensor products and virtual Thom twists;

    \item \(p_\sharp\) and \(p_!\) for every smooth separated morphism
    \(p\) of finite presentation;

    \item \(j_!=j_\sharp\) for every quasi-compact open immersion \(j\);

    \item \(p_*=p_!\) for every smooth proper morphism \(p\) of finite
    presentation.
\end{enumerate}
\end{proposition}

\begin{proof}
The proof of Theorem~\ref{thm:arbitrary-direct-image-cocontinuous} shows
that arbitrary pullback preserves compact objects.  Virtual Thom twists
preserve compactness because their Thom objects are compact and invertible.
The standard compact generators are closed under tensor product: the tensor
product of two twisted smooth motives is the corresponding twisted motive of
their fibre product over the base.  Exactness and preservation of retracts
therefore show that tensor products preserve all compact objects.

Smooth direct image preserves compact objects because its right adjoint
preserves filtered colimits.  Smooth exceptional direct image is smooth
direct image followed by a Thom twist.  Open direct image is the smooth
direct image of an open immersion.  The final assertion follows from smooth
proper ambidexterity.
\end{proof}

\begin{remark}[Nonformal constructibility statements]
\label{rem:nonformal-constructibility}
Preservation of constructible objects by \(f_!\), \(f_*\), or \(f^!\) for a
general separated morphism of finite presentation is not a consequence of
the identities established above.  Such preservation is an additional
constructibility theorem in ordinary motivic homotopy theory and can be
transported through stable comparison whenever it is available under the
chosen hypotheses on the classical bases.  It is therefore not included in
the unconditional six-operation theorem below.
\end{remark}

\begin{proposition}[Comparison of constructible objects]
\label{prop:constructible-classical-comparison}
Stable classical comparison restricts to an equivalence
\[
    \Theta_S^{\mathrm{st}}:
    \SH_c(S)
    \xrightarrow{\simeq}
    \SH_{\mathrm{cl}}(S_{\mathrm{cl}})^\omega.
\]
In particular, any ordinary theorem asserting preservation or duality of
compact objects transports to the spectral theory under the same classical
geometric hypotheses.
\end{proposition}

\begin{proof}
An equivalence of compactly generated presentable categories preserves and
reflects compact objects.  Apply
Theorem~\ref{thm:constructible-equals-compact} and stable classical
comparison.
\end{proof}


\section{Grothendieck--Verdier exchange}
\label{sec:six-verdier-duality}

Fix a quasi-compact quasi-separated spectral scheme \(S\).  Let
\[
    p_X:X\longrightarrow S,
    \qquad
    p_Y:Y\longrightarrow S
\]
belong to \(\mathcal E\), and let
\[
    f:X\longrightarrow Y
\]
belong to \(\mathcal E\) with a specified equivalence
\[
    p_Yf\simeq p_X.
\]
Put
\[
    \omega_{X/S}:=p_X^!\mathbf1_S,
    \qquad
    \omega_{Y/S}:=p_Y^!\mathbf1_S.
\]
Exceptional inverse-image composition gives
\[
    \omega_{X/S}
    \simeq
    f^!\omega_{Y/S}.
\]

\begin{definition}[Verdier duality functor]
\label{def:verdier-duality-functor}
Define
\[
    \mathbb D_{X/S}(A)
    :=
    \underline{\operatorname{Hom}}_X
    \bigl(
      A,
      \omega_{X/S}
    \bigr).
\]
When the base \(S\) is understood, write simply \(\mathbb D_X\).
\end{definition}

\begin{proposition}[Base change for relative dualizing objects]
\label{prop:dualizing-object-base-change}
Let \(g:T\to S\), put \(X_T=X\times_ST\), and write
\[
    g_X:X_T\longrightarrow X,
    \qquad
    p_{X_T}:X_T\longrightarrow T.
\]
Construction~\ref{constr:pullback-exceptional-exchange} gives a canonical
transformation
\[
    g_X^*\omega_{X/S}
    \longrightarrow
    \omega_{X_T/T}.
\]
It is an equivalence when \(g\) is smooth, separated, and of finite
presentation.  These transformations satisfy the base-change pasting laws.
\end{proposition}

\begin{proof}
Apply Construction~\ref{constr:pullback-exceptional-exchange} to the
Cartesian square obtained from \(p_X\) by base change along \(g\), and
evaluate it on \(\mathbf1_S\):
\[
    g_X^*p_X^!\mathbf1_S
    \longrightarrow
    p_{X_T}^!g^*\mathbf1_S
    \simeq
    p_{X_T}^!\mathbf1_T.
\]
Smooth invertibility and pasting are
Proposition~\ref{prop:smooth-pullback-exceptional-exchange}.
\end{proof}

\begin{theorem}[Grothendieck--Verdier exchange]
\label{thm:grothendieck-verdier-exchange}
For every \(A\in\SH(X)\), there is a canonical equivalence
\[
    \mathbb D_{Y/S}(f_!A)
    \xrightarrow{\simeq}
    f_*\mathbb D_{X/S}(A).
\]
It is compatible with composition in \(f\) and with localization.

For every base change \(g:T\to S\), write
\[
    X_T:=X\times_ST,
    \qquad
    Y_T:=Y\times_ST,
\]
and let
\[
    g_X:X_T\to X,
    \qquad
    g_Y:Y_T\to Y,
    \qquad
    f_T:X_T\to Y_T
\]
be the induced morphisms.  The displayed equivalence is compatible with the
canonical base-change transformations for direct image, internal Hom, and
the relative dualizing objects.  If \(g\) is smooth, separated,
and of finite presentation and \(A\) is strongly dualizable, all of these
base-change comparisons are equivalences; under them, the pullback of the
displayed exchange equivalence is the Grothendieck--Verdier exchange
equivalence over \(T\).
\end{theorem}

\begin{proof}
Apply Theorem~\ref{thm:exceptional-hom-projection} with
\[
    B=\omega_{Y/S}.
\]
It gives
\[
\begin{aligned}
    \mathbb D_{Y/S}(f_!A)
    &=
    \underline{\operatorname{Hom}}_Y
    \bigl(
      f_!A,
      \omega_{Y/S}
    \bigr) \\
    &\simeq
    f_*
    \underline{\operatorname{Hom}}_X
    \bigl(
      A,
      f^!\omega_{Y/S}
    \bigr) \\
    &\simeq
    f_*
    \underline{\operatorname{Hom}}_X
    \bigl(
      A,
      \omega_{X/S}
    \bigr) \\
    &=
    f_*\mathbb D_{X/S}(A).
\end{aligned}
\]
Composition and localization compatibility follow from those of the
exceptional Hom--projection identity, exceptional inverse-image
composition, and the localization triangles.

For arbitrary \(g\), pullback has canonical transformations through direct
image and internal Hom, and
Proposition~\ref{prop:dualizing-object-base-change} supplies the comparison
of dualizing objects.  Naturality of the Hom--projection identity gives the
asserted comparison diagram.  If \(g\) is smooth,
Theorem~\ref{thm:smooth-base-change-arbitrary-direct-image} gives
\[
    g_Y^*f_*
    \xrightarrow{\simeq}
    f_{T*}g_X^*.
\]
Proposition~\ref{prop:dualizing-object-base-change} identifies the pulled-back
dualizing object.  Finally, strong dualizability of \(A\) gives
\[
    g_X^*
    \underline{\operatorname{Hom}}_X
    \bigl(
      A,
      \omega_{X/S}
    \bigr)
    \xrightarrow{\simeq}
    \underline{\operatorname{Hom}}_{X_T}
    \bigl(
      g_X^*A,
      g_X^*\omega_{X/S}
    \bigr).
\]
Under these equivalences the comparison diagram becomes the defining
Hom--projection calculation over \(T\).
\end{proof}

\begin{corollary}[Proper Verdier exchange]
\label{cor:proper-verdier-exchange}
If \(f:X\to Y\) is proper and belongs to \(\mathcal E\), then
\[
    \mathbb D_{Y/S}(f_*A)
    \xrightarrow{\simeq}
    f_*\mathbb D_{X/S}(A).
\]
\end{corollary}

\begin{proof}
Use the proper normalization \(f_!\simeq f_*\) in
Theorem~\ref{thm:grothendieck-verdier-exchange}.
\end{proof}

\begin{proposition}[Smooth dualizing object]
\label{prop:smooth-dualizing-object}
Let \(p:X\to S\) be smooth, separated, and of finite presentation.  Then
\[
    \omega_{X/S}
    \simeq
    \operatorname{th}_X(T_p).
\]
Consequently,
\[
    \mathbb D_{X/S}(A)
    \simeq
    \underline{\operatorname{Hom}}_X
    \bigl(
      A,
      \operatorname{th}_X(T_p)
    \bigr).
\]
\end{proposition}

\begin{proof}
Smooth purity gives
\[
    p^!\mathbf1_S
    \simeq
    \Sigma^{T_p}p^*\mathbf1_S
    \simeq
    \Sigma^{T_p}\mathbf1_X.
\]
This is the Thom object \(\operatorname{th}_X(T_p)\).
\end{proof}

\subsection{Biduality and dualizing coefficients}

\begin{definition}[Constructibly dualizing coefficient]
\label{def:constructibly-dualizing-coefficient}
Let \(K\in\SH(S)\).  For \(p:X\to S\) in \(\mathcal E\), put
\[
    \omega_{X/S}(K):=p^!K
\]
and define
\[
    \mathbb D_{X/S,K}(A)
    :=
    \underline{\operatorname{Hom}}_X
    \bigl(
      A,
      \omega_{X/S}(K)
    \bigr).
\]
The object \(K\) is \emph{constructibly dualizing over \(S\)} if, for every
\(p:X\to S\) in \(\mathcal E\), the functor \(\mathbb D_{X/S,K}\) preserves
constructible objects and the canonical biduality transformation
\[
    A
    \longrightarrow
    \mathbb D_{X/S,K}\mathbb D_{X/S,K}(A)
\]
is an equivalence for every constructible object \(A\in\SH_c(X)\).
\end{definition}

\begin{proposition}[Conditional constructible biduality]
\label{prop:conditional-constructible-biduality}
Let \(K\in\SH(S)\) be constructibly dualizing over \(S\).  Then, for every
\(p:X\to S\) in \(\mathcal E\), the functor
\[
    \mathbb D_{X/S,K}:
    \SH_c(X)^{\mathrm{op}}
    \longrightarrow
    \SH_c(X)
\]
is an anti-equivalence with square canonically equivalent to the identity.
For every \(f:X\to Y\) in \(\mathcal E\) over \(S\), there is a
canonical exchange equivalence
\[
    \mathbb D_{Y/S,K}(f_!A)
    \simeq
    f_*\mathbb D_{X/S,K}(A).
\]
\end{proposition}

\begin{proof}
The exchange formula is proved exactly as
Theorem~\ref{thm:grothendieck-verdier-exchange}, replacing \(\mathbf1_S\)
by \(K\).  Biduality is part of the definition of a constructibly dualizing
coefficient.  The biduality transformation makes the duality functor fully
faithful, and applying it once more gives essential surjectivity.
\end{proof}

\begin{remark}[Scope of biduality]
\label{rem:biduality-scope}
No assertion is made that \(\mathbf1_S\) is constructibly dualizing for an
arbitrary qcqs spectral scheme \(S\).  Such a statement requires a genuine
duality theorem in the ordinary theory and is not a formal consequence of
the six operations.  Whenever the ordinary coefficient
\(\Theta_S^{\mathrm{st}}K\) is known to be constructibly dualizing, stable
comparison transports that result to \(K\).
\end{remark}

\section{The coherent six-operation theorem}
\label{sec:coherent-six-operation-theorem}

\begin{theorem}[The six operations in brave new motivic homotopy theory]
\label{thm:coherent-six-operations}
The assignment
\[
    S\longmapsto\SH(S)
\]
on quasi-compact quasi-separated spectral schemes has the following
structures and properties.

\begin{enumerate}
    \item \textbf{Stable symmetric monoidal coefficients.}
    For every \(S\), the category \(\SH(S)\) is compactly generated,
    stable, presentable, closed symmetric monoidal, and equipped with virtual
    Thom twists
    \[
        \Sigma^\xi,
        \qquad
        \xi\in K(S).
    \]

    \item \textbf{Ordinary adjunction.}
    For every morphism \(f:X\to Y\), there is an adjunction
    \[
        f^*:\SH(Y)
        \rightleftarrows
        \SH(X):f_*,
    \]
    the functor \(f^*\) is colimit-preserving and strong symmetric
    monoidal, and \(f_*\) preserves all small colimits.  Ordinary direct
    image satisfies base change along smooth separated morphisms of finite
    presentation:
    \[
        g^*f_*
        \simeq
        f'_*(g')^*.
    \]

    \item \textbf{Exceptional adjunction.}
    For every separated morphism of finite presentation
    \(f:X\to Y\), there is an adjunction
    \[
        f_!:\SH(X)
        \rightleftarrows
        \SH(Y):f^!.
    \]
    The functor \(f_!\) preserves all small colimits.

    \item \textbf{Coherent composition.}
    There are coherently unital and associative equivalences
    \[
        g_!f_!\simeq(gf)_!,
        \qquad
        (gf)^!\simeq f^!g^!.
    \]

    \item \textbf{Exceptional base change.}
    For every Cartesian square
    \[
    \begin{tikzcd}[column sep=large,row sep=large]
        X' \arrow[r,"g'"] \arrow[d,"f'"']
            & X \arrow[d,"f"] \\
        Y' \arrow[r,"g"']
            & Y
    \end{tikzcd}
    \]
    with \(f\) separated and of finite presentation, there are canonical
    equivalences
    \[
        g^*f_!
        \simeq
        f'_!(g')^*,
    \]
    \[
        (g')_*f'^!
        \simeq
        f^!g_*.
    \]
    There is also a canonical transformation
    \[
        (g')^*f^!
        \longrightarrow
        f'^!g^*,
    \]
    which is an equivalence when \(g\) is smooth, separated, and of finite
    presentation.

    \item \textbf{Projection and internal Hom.}
    There are canonical equivalences
    \[
        f_!(A\otimes_Xf^*B)
        \simeq
        f_!A\otimes_YB,
    \]
    \[
        \underline{\operatorname{Hom}}_Y(f_!A,B)
        \simeq
        f_*\underline{\operatorname{Hom}}_X(A,f^!B),
    \]
    \[
        f^!\underline{\operatorname{Hom}}_Y(B,C)
        \simeq
        \underline{\operatorname{Hom}}_X(f^*B,f^!C).
    \]

    \item \textbf{Open and proper normalization.}
    If \(j\) is a quasi-compact open immersion, then
    \[
        j_!\simeq j_\sharp,
        \qquad
        j^!\simeq j^*.
    \]
    If \(p\) is proper and separated of finite presentation, then
    \[
        p_!\simeq p_*.
    \]

    \item \textbf{Smooth purity.}
    If \(p\) is smooth, separated, and of finite presentation, then
    \[
        p^!
        \simeq
        \Sigma^{T_p}p^*,
        \qquad
        p_!
        \simeq
        p_\sharp\Sigma^{-T_p}.
    \]
    If \(p\) is \'{e}tale, separated, and of finite presentation, these reduce to
    \[
        p^!\simeq p^*,
        \qquad
        p_!\simeq p_\sharp.
    \]

    \item \textbf{Closed purity transformation.}
    For every quasi-smooth closed immersion \(i:Z\hookrightarrow X\) with
    finite locally free conormal module, there is a canonical transformation
    \[
        \Sigma^{-\nu_i}i^*
        \longrightarrow
        i^!,
    \]
    compatible with composition, Thom twists, and the fundamental-class
    base-change and excess formulas.  It is an
    equivalence on every coefficient object pure for \(i\).

    \item \textbf{Localization and recollement.}
    For every complementary closed--open pair
    \[
        Z\xrightarrow{i}X\xleftarrow{j}U
    \]
    with \(j\) quasi-compact, there are exact triangles
    \[
        j_!j^*A
        \longrightarrow
        A
        \longrightarrow
        i_*i^*A,
    \]
    \[
        i_*i^!A
        \longrightarrow
        A
        \longrightarrow
        j_*j^*A.
    \]

    \item \textbf{Descent and comparison.}
    The coefficient system satisfies spectral Nisnevich descent, derived
    nilpotent invariance, and stable comparison with ordinary motivic stable
    homotopy theory.  Stable comparison identifies all six operations and
    their canonical transformations with their ordinary counterparts.  It
    carries the virtual closed-purity transformation to the ordinary
    specialization class attached to the classical truncation of the
    spectral deformation space.

    \item \textbf{Correspondences.}
    The assignment
    \[
        \bigl(
          X\xleftarrow{u}K\xrightarrow{v}Y
        \bigr)
        \longmapsto
        v_!u^*
    \]
    extends to a lax symmetric monoidal \((\infty,2)\)-functor on
    exceptional correspondences, with proper reversed maps of apices as
    \(2\)-morphisms.
\end{enumerate}

All composition, base-change, projection, tensor, localization, and
comparison transformations belonging to the six-operation formalism satisfy
the corresponding higher pasting laws.  The closed-purity transformations
satisfy the binary composition, derived base-change, excess, Thom-twist, and
coefficient-naturalities proved above.
\end{theorem}

\begin{proof}
The stable symmetric monoidal coefficient system, ordinary adjunction,
Thom twists, localization, and Nisnevich descent are the output of
Chapter~\ref{chap:stable-brave-new-motivic}.  Exceptional direct image and
its cocontinuity are Definition~\ref{def:exceptional-direct-image} and
Corollary~\ref{cor:exceptional-direct-image-cocontinuous}.  Coherent
composition is Theorem~\ref{thm:exceptional-composition}; exceptional
inverse image and its composition are
Definition~\ref{def:exceptional-inverse-image} and
Proposition~\ref{prop:exceptional-inverse-composition}.  Exceptional base
change is Theorem~\ref{thm:exceptional-base-change}, its right-adjoint mate
is Theorem~\ref{thm:right-adjoint-exceptional-base-change}, and the
pullback--exceptional exchange transformation is
Construction~\ref{constr:pullback-exceptional-exchange} with smooth
invertibility from
Proposition~\ref{prop:smooth-pullback-exceptional-exchange}.  Smooth base
change for arbitrary ordinary direct image is
Theorem~\ref{thm:smooth-base-change-arbitrary-direct-image}.  Projection
and internal-Hom formulas are
Theorems~\ref{thm:exceptional-projection-formula},
\ref{thm:exceptional-hom-projection}, and
\ref{thm:exceptional-pullback-hom}.  Open and proper normalization are
Theorems~\ref{thm:exceptional-open-normalization} and
\ref{thm:exceptional-proper-normalization}.  Smooth purity is
Theorem~\ref{thm:six-smooth-purity} and
Corollary~\ref{cor:smooth-exceptional-direct-image}.  The closed-purity
transformation is
Construction~\ref{constr:absolute-closed-purity-transformation} and
Proposition~\ref{prop:closed-purity-formal-properties}.  Classical
comparison is
Theorem~\ref{thm:comparison-exceptional-direct-image} and
Corollary~\ref{cor:comparison-exceptional-inverse-image}.  The
correspondence formulation is
Theorem~\ref{thm:exceptional-correspondence-functor}.  The higher coherence
assertions for the six-operation formalism are supplied by finite
compactification grids, Beck--Chevalley pasting, projection-formula pasting,
and the calculus of adjoint mates.  The closed-purity assertions are those
of
Theorem~\ref{thm:virtual-class-composition-coherence},
Proposition~\ref{prop:virtual-class-base-change-excess}, and
Proposition~\ref{prop:closed-purity-formal-properties}.
\end{proof}

\section{The final brave new motivic package}
\label{sec:final-brave-new-motivic-package}

\begin{theorem}[The completed foundational construction]
\label{thm:completed-foundational-construction}
For every quasi-compact quasi-separated spectral scheme \(S\), brave new
motivic homotopy theory assigns a compactly generated stable presentably
symmetric monoidal \(\infty\)-category
\[
    \SH(S).
\]
For every morphism there is the adjunction
\[
    f^*\dashv f_*,
\]
and for every separated morphism of finite presentation there is the
exceptional adjunction
\[
    f_!\dashv f^!.
\]
Together with
\[
    -\otimes-
    \dashv
    \underline{\operatorname{Hom}},
\]
these functors form a coherent six-operation formalism with localization,
proper and exceptional base change, smooth base change for ordinary direct
image, projection formulas, smooth purity, the closed-purity
transformation and its equivalence on coefficients pure for the immersion,
virtual Thom twists, compactification, recollement, Borel--Moore objects,
and Grothendieck--Verdier exchange.

Under stable classical comparison, the six-operation formalism is canonically
equivalent to the ordinary six-operation formalism over the classical
truncation.  The virtual closed-purity transformation corresponds to the
ordinary specialization class attached to the classical truncation of the
spectral deformation space.
\end{theorem}

\begin{proof}
This is the combination of
Theorem~\ref{thm:coherent-six-operations},
Theorem~\ref{thm:grothendieck-verdier-exchange}, and
Theorem~\ref{thm:comparison-exceptional-direct-image}.
\end{proof}

The foundational construction is now complete.  The sequence of chapters
has produced
\[
    S\longmapsto\SH(S),
\]
\[
    (f^*,f_*),
    \qquad
    (f_!,f^!),
\]
\[
    (\otimes,\underline{\operatorname{Hom}}),
\]
together with virtual Thom twists, localization, smooth purity,
coefficientwise closed purity, proper and smooth base change, projection
formulas, compactification, recollement, and duality.  Within the standing
qcqs and separated finite-presentation scope, the six operations and their
basic composition, exchange, localization, smooth-purity, coefficientwise
closed-purity, and duality structures have now been constructed.
    